% !TeX encoding = UTF-8
% !TeX program = xelatex
\documentclass[UTF8,a4paper,11pt]{ctexart}

\usepackage{amsmath,amssymb,amsfonts}
\usepackage{booktabs}
\usepackage{longtable}
\usepackage{geometry}
\usepackage{xcolor}
\usepackage{hyperref}
\usepackage{enumitem}
\usepackage{array}
\usepackage{fancyhdr}
\usepackage{graphicx}
\usepackage[ruled,linesnumbered]{algorithm2e}
\usepackage{url}

\geometry{left=2.6cm,right=2.6cm,top=2.8cm,bottom=2.6cm}
\hypersetup{colorlinks=true,linkcolor=blue!60!black,urlcolor=blue!60!black,citecolor=blue!60!black}

\definecolor{boxbg}{RGB}{245,247,250}
\definecolor{boxframe}{RGB}{90,120,160}
\definecolor{warnbg}{RGB}{255,248,240}
\definecolor{warnframe}{RGB}{200,120,40}

\newenvironment{explainbox}[1]
{\begin{center}\begin{minipage}{0.94\textwidth}\colorbox{boxbg}{\parbox{0.92\textwidth}{\textbf{#1}\par\vspace{2pt}}}}
{\end{minipage}\end{center}}

\newcommand{\sym}[1]{\texttt{#1}}
\newcommand{\term}[1]{\textbf{#1}}
\newcommand{\hl}[1]{\textcolor{blue!70!black}{\textbf{#1}}}
\newcommand{\warn}[1]{\textcolor{red!70!black}{\textbf{#1}}}

\pagestyle{fancy}
\fancyhf{}
\rhead{\small Tender 零基础全流程教学手册 · V5.4}
\lhead{\small\thepage}
\renewcommand{\headrulewidth}{0.3pt}

\title{\textbf{Tender 心智化教育系统\\[6pt]零基础全流程教学手册·完整版}\\[10pt]
\large 第一至四期：做、验、防、接\\[4pt]
\large 增补：开花、结果、森林\\[4pt]
\large 版本 5.3}
\author{}
\date{2026 年 10 月 31 日}

\setCJKmainfont{SimSun}[BoldFont=SimHei]
\setCJKsansfont{Microsoft YaHei}
\setCJKmonofont{SimSun}

\begin{document}
\maketitle
\thispagestyle{empty}

\begin{center}
\fbox{\parbox{0.88\textwidth}{
\vspace{4pt}
\textbf{这本书写给谁}\\[4pt]
写给\textbf{没学过高等数学、没写过算法、没读过论文}的人。\\
读者只需要会：加减乘除、看百分比、用 Excel。\\[6pt]
每一个公式都会先说明\hl{它在说什么}，再说明\hl{每个符号是什么}，\\
最后说明\hl{它有什么用}，并配一个生活例子。\\[6pt]
\textbf{全文不假设任何先修知识。}
\vspace{4pt}
}}
\end{center}

\vspace{8pt}
\tableofcontents
\newpage


\part{项目总览与阅读指南}

\section{中心主旨：让 AI 成为森林的气候，让人成为生命树的作者}

\begin{quote}
\textbf{生命不是被某个框架评价出来的形状，而是在有限时间里，通过选择、关系、承担和创造长出来的年轮。}
\end{quote}

现有培养体系过于庞大，并且伴随人的一生。它提供秩序，也带来持续的压力和对未知的恐惧。外部标准一旦变成人的内在边界，一个人可能很会解决问题，却逐渐失去对自己生命的解释权。

Tender 要做的，是一层贯穿终身的心智生长环境：

\begin{itemize}[leftmargin=2em]
\item 在压力出现时，帮助人命名压力；
\item 在未知面前，帮助人把恐惧拆成可以理解和行动的风险；
\item 在选择面前，帮助人看见自己真正重视什么、愿意承担什么；
\item 在行动之后，帮助人保存结果、复盘和形成自己的价值证据；
\item 在自主能力增强之后，逐步退到背景，把解释权和决定权归还给人。
\end{itemize}

\begin{center}
\begin{longtable}{@{}p{2.2cm}p{5.5cm}p{6.3cm}@{}}
\toprule
\textbf{生命树} & \textbf{人的成长} & \textbf{Tender 如何承接} \\
\midrule
种子 & 尚未展开的可能性 & 不用一次测评定义一个人 \\
根系 & 安全、关系、身体状态与真实反馈 & 先理解支持条件，再讨论表现 \\
树干 & 长期选择与承担形成的连续性 & 记录跨时间的承诺、行动和结果 \\
枝叶 & 探索、试错、跨域连接与角色经验 & 提供多 Agent 视角，不替人规定方向 \\
开花 & 第一次有足够证据看见自己的整体成长 & 用长期证据形成可更正的状态假设 \\
结果 & 独立迁移、创造作品、影响他人并承担后果 & 关注真实结果和迁移，而非奖励等级 \\
森林 & 不同生命在同一生态中互相支撑 & 保持多样性，拒绝把人修剪成同一种形状 \\
\bottomrule
\end{longtable}
\end{center}

\begin{explainbox}{一句话记住}
\textbf{AI 可以模拟风雨，却不能替一棵树长出年轮；它可以提供地图，却不能替一个人决定要去哪里。}
\end{explainbox}

这里的“成熟”，指一个人能够理解自己的状态，选择值得追求的事情，承担选择后的结果，在需要时调整、求助并重新出发。

Tender 的中心主旨是：

\begin{quote}
\textbf{让 AI 成为森林的气候，让人成为生命树的作者。一个人不必成为最聪明的那个，也不必成为最有用的那个，他仍然可以通过自己的选择、关系、创造和承担，长成独一无二的生命形态。}
\end{quote}

生命树是叙事与建模语法。具体科学主张仍需通过真实数据、预注册实验和独立复核验证；Tender 把模型输出当作可更正的状态假设，不把它当作诊断，也不承诺替代考试或保证结果。


这一部分先回答三个问题：
\begin{enumerate}[leftmargin=2em]
\item 这个项目到底在做什么？
\item 整本书按什么架构组织？
\item 不同角色应该怎么读？
\end{enumerate}

\section{这份手册怎么用}

\subsection{一句话架构}

\begin{quote}
\textbf{Tender 是一个“采集行为 $\to$ 建模认知 $\to$ 验证结论 $\to$ 干预教学 $\to$ 连接群体与平台”的心智化教育系统。}
\end{quote}

\subsection{三层技术栈}

\begin{center}
\begin{longtable}{@{}p{2.6cm}p{4.2cm}p{7.0cm}@{}}
\toprule
\textbf{层} & \textbf{角色} & \textbf{回答什么} \\
\midrule
DAOQL 层（底层） & 事实与因果 & 发生了什么？因为什么？证据是什么？ \\
Tender 层（中层） & 心智化建模 & 学生现在是什么状态？不确定性多大？ \\
Agent 层（上层） & 交互与执行 & 现在该对谁说什么、做什么、怎么跟进？ \\
\bottomrule
\end{longtable}
\end{center}

\subsection{六层产品架构}

\begin{center}
\begin{longtable}{@{}p{2.2cm}p{4.0cm}p{4.4cm}p{4.2cm}@{}}
\toprule
\textbf{层} & \textbf{做什么} & \textbf{核心问题} & \textbf{对应章节} \\
\midrule
数据层 & 采集会话、事件、JOL、根号、犹豫、教师观察、互动，写入 DAOQL & 我们到底有什么数据？ & 第 3 部分 \\
建模层 & 用 $\theta$、节律、会话、知识追踪刻画学生，形成 Tender 状态 & 怎么把数据变成状态？ & 第 4 部分 \\
验证层 & 安慰剂、时间重标、换单位、AIC、功效、预注册 & 怎么防止自己骗自己？ & 第 5 部分 \\
应用层 & 元认知、犹豫、外部脚手架、教师 metameta，由 Agent 执行 & 怎么把状态变成教学干预？ & 第 6 部分 \\
群体/平台层 & 群体拓扑、平台逆向、推荐算法、供应链 & 怎么从个人扩展到群体和平台？ & 第 7 部分 \\
Agent 交互层 & 学生辅导、教师副驾驶、家长沟通、分组、运营 & 谁来执行、谁来跟进？ & 第 8 部分 \\
治理层 & 合规、隐私、伦理、翻车清单 & 哪些事不能做？ & 第 10 部分 \\
\bottomrule
\end{longtable}
\end{center}

\subsection{架构图（从上到下读）}

\begin{center}
\begin{longtable}{@{}c@{}}
\toprule
\textbf{项目架构} \\
\midrule
Agent 交互层：Student Tutor、Metacognition Coach、Teacher Copilot、Parent、Group Facilitator、Sales \\
$\downarrow$ \\
Tender 心智层：$\theta$、Bloch 球、情绪、元认知、知识追踪、节律、基线、群体拓扑 \\
$\downarrow$ \\
DAOQL 因果层：事件图、因果关系、证据、血缘、时间线、版本 \\
$\uparrow$ \\
数据来源：课堂、LMS、浏览器插件、教师卡、家长端、平台互动 \\
$\downarrow$ \\
治理层：合规、隐私、伦理、翻车清单 \\
\bottomrule
\end{longtable}
\end{center}

\subsection{按角色的阅读路径}

\begin{center}
\begin{longtable}{@{}p{2.6cm}p{4.2cm}p{8.0cm}@{}}
\toprule
\textbf{角色} & \textbf{先读什么} & \textbf{再读什么} \\
\midrule
老师 & 第 1 部分：零基础入门；第 5 部分：元认知与教学干预；第 7 部分：课堂操作手册 & 第 2 部分：数据怎么采；第 6 部分：群体与平台 \\
产品经理 & 第 0 部分：项目总览；第 5 部分：应用层；第 6 部分：群体/平台层 & 第 2、3 部分：数据与建模 \\
工程师 & 第 2 部分：数据层；第 3 部分：建模层；第 7 部分：实验复现 & 第 4 部分：验证层；第 8 部分：治理 \\
研究者 & 第 3 部分：建模层；第 4 部分：验证层；第 7 部分：实验复现 & 第 5 部分：元认知实验结果；第 6 部分：群体实验 \\
\bottomrule
\end{longtable}
\end{center}

\section{项目愿景与四个关键词}

\subsection{四个关键词}

\begin{center}
\begin{longtable}{@{}p{2.2cm}p{4.0cm}p{8.4cm}@{}}
\toprule
\textbf{关键词} & \textbf{白话} & \textbf{在本书哪里} \\
\midrule
做 & 先建一个能算的模型 & 第 2、3 部分：数据层与建模层 \\
验 & 再检查这个模型是不是骗自己 & 第 4 部分：验证层 \\
防 & 事先写死规则，防止事后找理由 & 第 4 部分：功效与预注册 \\
接 & 把信号接回学生、教师、群体、环境 & 第 5、6 部分：应用层与群体/平台层 \\
\bottomrule
\end{longtable}
\end{center}

\subsection{项目的一句话}

\begin{explainbox}{一句话}
\textbf{不是给学生贴标签，而是找到最省力的那一个杠杆。}\\
学生、教师、同伴、排课、任务，互相影响；
Tender 要做的，是把这个耦合系统里可测量、可干预的部分找出来。
\end{explainbox}

\section{术语统一说明}

原来的四本分册是按时间写的，所以会出现“第一期/第二期”“这一期/上一期”。
本完整版按项目架构重排后，请按下面的对应关系理解：

\begin{center}
\begin{longtable}{@{}p{4.2cm}p{4.4cm}p{6.2cm}@{}}
\toprule
\textbf{旧写法} & \textbf{新位置} & \textbf{怎么理解} \\
\midrule
第一期 & 第 2--3 部分：数据层与建模层 & 最初的数据采集和核心公式 \\
第二期 & 第 4 部分：验证层 & 安慰剂、时间重标、换单位 \\
第三期 & 第 4--5 部分：验证层与应用层 & 功效、预注册、元认知、metameta \\
第四期 & 第 5--6 部分：应用层与群体/平台层 & 犹豫、外部脚手架、耦合基线、群体拓扑 \\
“这一期” & “本章” & 按当前章节上下文理解 \\
“上一期” & “前面章节” & 按架构顺序回看 \\
\bottomrule
\end{longtable}
\end{center}

\section{全书地图}

\begin{center}
\begin{longtable}{@{}p{1.6cm}p{4.0cm}p{4.6cm}p{4.6cm}@{}}
\toprule
\textbf{部分} & \textbf{名称} & \textbf{回答什么} & \textbf{关键产出} \\
\midrule
第 1 部分 & 数学、符号与零基础入门 & 我不懂数学，怎么读？ & 生活类比、符号词典、小明例子 \\
第 2 部分 & 数据与基础设施层 & 我们有什么数据？怎么存？ & 会话、事件、数据表、采集工具 \\
第 3 部分 & Tender 建模核心 & 怎么把数据变成状态？ & $\theta$、Bloch 球、节律、知识追踪 \\
第 4 部分 & 验证、防错与预注册 & 怎么防止自己骗自己？ & 安慰剂、时间重标、AIC、功效、预注册 \\
第 5 部分 & 元认知、教学干预与教师观察者 & 怎么把状态变成干预？ & JOL、根号、犹豫四象限、外部脚手架、metameta \\
第 6 部分 & 群体、平台与供应链 & 怎么从个人扩展到群体和平台？ & 群体拓扑、平台逆向、推荐算法、供应链 \\
第 7 部分 & 实验复现与附录 & 怎么跑、怎么复现？ & E1--E6、脚本、结果、文件索引 \\
第 8 部分 & 治理与合规 & 哪些事不能做？ & 十条红线、隐私、伦理、翻车清单 \\
第 9 部分 & 统一符号与总结 & 所有符号和一句话 & 符号词典、最重要的 20 个数字 \\
\bottomrule
\end{longtable}
\end{center}

\section{四期地图：做、验、防、接}

\begin{center}
\begin{longtable}{@{}p{1.6cm}p{3.0cm}p{4.4cm}p{5.0cm}@{}}
\toprule
\textbf{期} & \textbf{动词} & \textbf{解决什么} & \textbf{核心工具} \\
\midrule
第一期 & 做 & 怎么建模、怎么算指标 & 会话、节律、知识追踪、$\theta$ \\
第二期 & 验 & 怎么检查模型是不是骗自己 & 安慰剂、时间重标、换单位 \\
第三期 & 防 & 怎么防止事后找理由 & 功效、预注册、元认知、metameta \\
第四期 & 接 & 怎么把信号接回系统 & 犹豫、外部脚手架、耦合基线、群体拓扑 \\
\bottomrule
\end{longtable}
\end{center}


\section{本书的三种证据}

\begin{center}
\begin{longtable}{@{}p{3.0cm}p{4.0cm}p{7.0cm}@{}}
\toprule
\textbf{证据} & \textbf{例子} & \textbf{能说什么} \\
\midrule
真实公开数据 & EdNet、OULAD、Stanford MOOCPosts & 能说“在这个公开数据上，效应方向和大小是什么” \\
真实课堂数据 & 9 篇 56--65 题、实验 B & 能说“在这 6--19 人的班级里，初步信号是什么” \\
合成自检 & E1/E5/E6 模拟 & 只能验证“估计器能不能恢复真值”，不能说学生真实结论 \\
\bottomrule
\end{longtable}
\end{center}

\begin{explainbox}{读本书的硬规则}
看到任何一个数字，先问：\\
\textbf{这是真实数据，还是模拟？是相关，还是因果？}
\end{explainbox}


\section{五个最容易犯的错}

\begin{enumerate}[leftmargin=2em]
\item \textbf{把相关当因果}：两件事一起变，不代表一个导致另一个。
\item \textbf{把不显著当没有效应}：可能只是样本太小、功效不够。
\item \textbf{把基线当性格}：基线是耦合稳定点，环境变了它也会变。
\item \textbf{把犹豫当不会}：犹豫只是信号，必须先验证 AUC。
\item \textbf{把黄金比当结论}：目前没有证据，只能当隐喻。
\end{enumerate}



\part{数学、符号与零基础入门}

\section{这本书怎么读}

\subsection{三遍法}

\begin{enumerate}[leftmargin=2em]
\item \textbf{第一遍：只读白话。}跳过所有公式和表格，先知道每一章在讲什么故事。
\item \textbf{第二遍：看公式和符号。}每个公式下面都有逐符号解释，配一个具体数字例子。
\item \textbf{第三遍：跑脚本、看数据。}照着附录里的命令跑一遍，看到真实的图和数字。
\end{enumerate}

\begin{explainbox}{一句话}
\textbf{先懂意思，再看符号，最后动手。}
不要一上来就背公式。
\end{explainbox}

\subsection{这本书的四个关键词}

\begin{center}
\begin{longtable}{@{}p{2.4cm}p{4.6cm}p{6.8cm}@{}}
\toprule
\textbf{关键词} & \textbf{白话} & \textbf{在本书哪里} \\
\midrule
做 & 先建一个能算的模型 & 第一期：会话、节律、知识追踪 \\
验 & 再检查这个模型是不是骗自己 & 第二期：安慰剂、时间重标、换单位 \\
防 & 事先写死规则，防止事后找理由 & 第三期：功效、预注册、元认知 \\
接 & 把信号接回学生、教师、群体、环境 & 第四期：犹豫、基线、群体拓扑 \\
\bottomrule
\end{longtable}
\end{center}


\section{什么是模型：地图不是地形}

\subsection{生活类比}

你手里有一张北京地铁图。
地铁图不是北京，它把真实的城市简化成一堆线和点。
但有了它，你就能规划路线。

\textbf{模型就是地铁图。}
它不是现实本身，它是现实的简化版。
好的模型不是“最像现实”，而是“在你关心的问题上够用”。

\subsection{形式定义}

模型是一个映射：

\begin{equation}
\text{输入}\ \rightarrow\ \text{模型}\ \rightarrow\ \text{输出}.
\end{equation}

符号说明：

\begin{itemize}[leftmargin=2em]
\item 输入：你观察到的数据，例如学生做对了哪几题；
\item 模型：你写的规则，例如“做对越多，能力越高”；
\item 输出：模型给出的判断，例如“这个学生掌握度 0.7”。
\end{itemize}

\subsection{在本书哪里用到}

第一期的 Tender $\theta$、知识追踪 BKT、会话模型，都是模型。
模型不是真理，只是暂时可用的地图。


\section{什么是指标：体温计和体重秤}

\subsection{生活类比}

体温计把“发烧”变成一个数字：38.5℃。
体重秤把“胖瘦”变成一个数字：60 kg。
数字不完美，但比“我感觉有点热”更有用。

\textbf{指标就是把模糊的感觉变成可比较的数字。}

\subsection{形式定义}

指标是一个函数：

\begin{equation}
\text{指标}=\text{函数}(\text{原始数据}).
\end{equation}

例如：

\begin{equation}
\text{正确率}=\frac{\text{答对题数}}{\text{总题数}}.
\end{equation}

\subsection{好指标的三条标准}

\begin{enumerate}[leftmargin=2em]
\item \textbf{可重复}：同样的数据，今天算和明天算结果一样；
\item \textbf{可比较}：不同学生、不同时间能放在一起比；
\item \textbf{不容易骗}：不能只看一个数就下结论。
\end{enumerate}

\subsection{在本书哪里用到}

AUC、偏差、Brier、效应量、功效，都是指标。
指标不是结论，指标是证据的一部分。


\section{什么是概率：下雨概率 30\%}

\subsection{生活类比}

天气预报说“明天降雨概率 30\%”。
这不是说“明天一定下 30\% 的雨”，也不是说“明天一定不下雨”。
它说的是：在类似情况下，100 次里有 30 次会下雨。

\subsection{形式定义}

概率 $P(A)$ 是事件 $A$ 发生的可能性，范围是：

\begin{equation}
0\le P(A)\le 1.
\end{equation}

\begin{itemize}[leftmargin=2em]
\item $P(A)=0$：不可能；
\item $P(A)=1$：一定；
\item $P(A)=0.3$：100 次里大约 30 次。
\end{itemize}

\subsection{最容易误解的地方}

概率不是命运。
一个学生“掌握度 0.7”，不是说这次一定对 7 题，
而是说他这类题的长期正确率大约 70\%。


\section{什么是基线：体温 36.5℃}

\subsection{生活类比}

医生不会只问“你体温多少”，还会问“你平时体温多少”。
36.5℃ 是正常基线，38.5℃ 才是异常。

\textbf{没有基线，就没有异常。}

\subsection{形式定义}

基线是“在正常情况下，这个指标通常是多少”。

\begin{equation}
\text{偏离}=\text{当前值}-\text{基线}.
\end{equation}

\subsection{本书的最重要升级}

第三期之前，基线常被当成“这个人固定的一个数”。
第四期把它升级为：

\begin{equation}
x^*=f(x^*,u^*,e^*;\theta).
\end{equation}

翻译成白话：

\begin{quote}
\textbf{基线 = 学生 × 环境 的耦合稳定点。}
\end{quote}

环境变了，基线也会移动。
所以“改安排”往往比“强制改习惯”更有效。


\section{什么是 AUC：医生诊断的区分能力}

\subsection{生活类比}

医生用一项检查区分“有病”和“没病”。
如果这项检查完全瞎猜，AUC = 0.5；
如果它完美区分，AUC = 1.0。

\subsection{形式定义}

\begin{equation}
\mathrm{AUC}
=
\Pr(\text{正例得分}>\text{负例得分})
+\frac12\Pr(\text{相等}).
\end{equation}

\begin{itemize}[leftmargin=2em]
\item 正例：你想找出来的东西，例如“错题”“有病”“不会”；
\item 负例：其余情况，例如“对题”“没病”“会”；
\item AUC = 0.5：没有区分能力；
\item AUC = 1.0：完美区分。
\end{itemize}

\subsection{在本书哪里用到}

元认知敏感性、犹豫、教师观察，都用 AUC。
但 AUC 不是准确率，它只看“能不能把两类分开”。


\section{什么是效应量和功效：放大镜够不够大}

\subsection{生活类比}

你想看一只蚂蚁。
用肉眼看，看不清；
用 10 倍放大镜，勉强；
用 100 倍放大镜，很清楚。

\textbf{效应量就是蚂蚁有多大；功效就是你的放大镜够不够。}
样本量就是放大镜的倍数。

\subsection{形式定义}

效应量：真实差别有多大，例如 $d=0.2$ 是小效应，$d=0.8$ 是大效应。

功效：如果真实差别存在，你的实验能把它检测出来的概率。

\begin{equation}
\text{功效}=P(\text{检测到效应}\mid\text{效应真的存在}).
\end{equation}

\subsection{本书最重要的功效数字}

\begin{center}
\begin{longtable}{@{}p{5.0cm}p{4.0cm}p{4.0cm}@{}}
\toprule
\textbf{场景} & \textbf{样本} & \textbf{能测出多大} \\
\midrule
相关分析 & 19 人 & $|r|\ge0.605$ \\
OULAD 最大效应 & 26{,}074 人 & 中位 $0.053$、最大 $0.289$ \\
OULAD 最大效应所需样本 & 约 92 人 & 检测 $0.289$ \\
\bottomrule
\end{longtable}
\end{center}

\begin{explainbox}{一句话}
\textbf{不显著，不等于没有效应；可能只是你的放大镜不够。}
\end{explainbox}


\section{什么是预注册：先写规则，再看数据}

\subsection{生活类比}

考试前先把答案写下来，考完后不能改。
如果你考完再写“我早就知道是这个答案”，那不算数。

\textbf{预注册就是在看数据之前，先把分析规则写死。}

\subsection{为什么要预注册}

因为人很会事后找理由：

\begin{itemize}[leftmargin=2em]
\item 本来想测 A，结果 A 不显著，就改测 B；
\item 本来想看整体，结果整体不显著，就挑一个子群体；
\item 本来想用 0.05 标准，结果刚好 0.06，就说“接近显著”。
\end{itemize}

预注册就是防止这些“事后找理由”。

\subsection{预注册锁死什么}

\begin{enumerate}[leftmargin=2em]
\item 主假设只有一个；
\item 样本量先算好；
\item 主要指标先写死；
\item 排除标准先写死；
\item 如果没显著，也要报告功效。
\end{enumerate}


\section{什么是元认知：知不知道自己会不会}

\subsection{生活类比}

考试时，有些学生知道“我这题不会”，
有些学生“以为自己会，结果错了”。
后者不是知识问题，是\textbf{元认知问题}。

\textbf{元认知 = 对自己认知过程的监控和判断。}

\subsection{三个一阶指标}

\begin{enumerate}[leftmargin=2em]
\item \textbf{JOL}（学习判断）：测试前估计自己能对几题；
\item \textbf{根号自标}：学生自己标“哪里不会/是重点”；
\item \textbf{犹豫}：响应时间、修改次数、停顿。
\end{enumerate}

\subsection{两个关键量}

\begin{itemize}[leftmargin=2em]
\item \textbf{校准偏差} $b=\hat p-p$：预测正确率减实际正确率；$>0$ 过度自信；
\item \textbf{敏感性 AUC}：能不能把“会”和“不会”分开。
\end{itemize}

\begin{explainbox}{最重要的一句话}
\textbf{相关不等于校准。}\\
一个学生能预测自己考得好不好（相关高），\\
不代表他预测的分数一定准（校准好）。
\end{explainbox}


\section{什么是 trait、state 和基线}

\subsection{三个词的白话解释}

\begin{center}
\begin{longtable}{@{}p{2.4cm}p{4.2cm}p{7.0cm}@{}}
\toprule
\textbf{词} & \textbf{白话} & \textbf{例子} \\
\midrule
trait & 比较稳定的倾向 & MBTI 连续分数、节律峰型 \\
state & 当前的状态 & 今天注意力差、刚考完试很累 \\
baseline & 当前耦合出来的稳定点 & 平时晚上 20:00 学习效率最高 \\
\bottomrule
\end{longtable}
\end{center}

\subsection{为什么不能把基线当 trait}

因为基线会随环境移动：

\begin{itemize}[leftmargin=2em]
\item 学校把晚自习从 20:00 调到 19:00；
\item 基线可能跟着移动；
\item 如果你把基线当固定性格，就会误判学生“变差了”。
\end{itemize}

\subsection{扰动实验}

想判断一个量是 trait 还是 baseline：

\begin{enumerate}[leftmargin=2em]
\item 改变环境；
\item 看它动不动；
\item 动得多，是 baseline/state；
\item 几乎不动，才更像 trait。
\end{enumerate}


\section{一个贯穿全书的例子：小明学语法填空}

我们用一个学生“小明”把四期串起来。

\subsection{第一期：做}

小明每次做语法填空，我们记录：

\begin{itemize}[leftmargin=2em]
\item 做了几题、对了几题；
\item 什么时候做（节律）；
\item 一次学习会话有多长；
\item 他的 $\theta$、掌握度。
\end{itemize}

\subsection{第二期：验}

我们发现小明的学习有“自激”现象：做得越多，接下来做得越快。
但要先做安慰剂检验和时间重标检验，确认这不是假象。
最后发现，真正的基本单位是\textbf{一次会话}。

\subsection{第三期：防}

我们事先算好：19 个学生只能测出中等以上效应。
然后让小明每次测试前写“估计能对几题”（JOL），
再让他自己标“哪里不会”（根号自标）。
我们发现：他可能预测得挺准，但不一定校准得好。

\subsection{第四期：接}

我们继续记录：

\begin{itemize}[leftmargin=2em]
\item 他做题时犹不犹豫（响应时间、修改次数）；
\item 教师对他的预测和信心（metameta）；
\item 他和谁讨论（群体拓扑）；
\item 他的节律和排课匹不匹配（耦合基线）。
\end{itemize}

最后我们发现：

\begin{itemize}[leftmargin=2em]
\item 顺应他的节律排课，比强制他改作息更有效；
\item 小圈子和桥接各有优势，取决于要触达还是要认可；
\item 教师观察需要折扣，不能当 ground truth；
\item 黄金比目前只是隐喻，真正可检验的是闭环和尺度不变。
\end{itemize}

\subsection{小明的故事告诉我们}

\begin{quote}
\textbf{教育不是一个数字，而是一个耦合系统。}\\
学生、教师、同伴、排课、任务，互相影响。\\
我们要做的不是给学生贴标签，而是找到最省力的那一个杠杆。
\end{quote}



\section{数学基础（零基础版）}
% ============================================================

这一章只讲后面会用到的数学。每个概念都用“生活例子”解释。

\subsection{角度与弧度}

\subsubsection{角度}

圆的一圈是 $360^\circ$。转半圈是 $180^\circ$，转四分之一圈是 $90^\circ$。
这个大家都熟。

\subsubsection{弧度}

数学里更常用的是\term{弧度}。定义很简单：

\begin{equation}
1 \text{ 圈} = 360^\circ = 2\pi \text{ 弧度}
\end{equation}

\begin{explainbox}{符号说明}
\begin{itemize}[leftmargin=1.4em,itemsep=1pt]
\item $\pi$（读作“派”）：圆周率，约等于 $3.14159$。
\item $2\pi \approx 6.283$：一整圈的弧度数。
\item 换算方法：$180^\circ = \pi$ 弧度，所以 $90^\circ = \pi/2 \approx 1.571$，$60^\circ = \pi/3 \approx 1.047$。
\end{itemize}
\end{explainbox}

\textbf{为什么用弧度}：因为弧度直接对应圆上的“弧长”。
半径为 1 的圆上，转过 1 弧度，弧长就是 1。这让后面的公式更简洁。

\subsubsection{三角函数：cos 和 sin}

在一个直角三角形里，取一个锐角 $\theta$（读作“西塔”），那么：

\begin{equation}
\cos\theta = \frac{\text{邻边}}{\text{斜边}}, \qquad
\sin\theta = \frac{\text{对边}}{\text{斜边}}
\end{equation}

\begin{explainbox}{符号说明}
\begin{itemize}[leftmargin=1.4em,itemsep=1pt]
\item $\theta$：希腊字母 theta，读音“西塔”。本文中它表示\hl{融合角}，取值范围 $[0,\pi]$。
\item $\cos\theta$：余弦，输出范围 $[-1, 1]$。
\item $\sin\theta$：正弦，输出范围 $[-1, 1]$。
\item $\cos^2\theta$ 是 $(\cos\theta)^2$ 的简写，即余弦值的平方。
\end{itemize}
\end{explainbox}

\textbf{一个必须记住的等式}（后面会反复用到）：

\begin{equation}
\cos^2\!\left(\frac{\theta}{2}\right) + \sin^2\!\left(\frac{\theta}{2}\right) = 1
\end{equation}

\textbf{通俗理解}：$\cos^2(\theta/2)$ 是一个“占比”，它和 $\sin^2(\theta/2)$ 加起来永远是 100\%。
后面的“概率”就是从这个占比来的。

\subsection{概率}

\term{概率}就是“可能性”，取值在 0 到 1 之间。0 表示绝不可能，1 表示必然发生。
也可以写成百分比：$0.7 = 70\%$。

\textbf{本文中概率的两个角色}：

\begin{enumerate}[leftmargin=1.6em]
\item \textbf{预测概率}：模型说“这个学生有 0.7 的概率答对”。
\item \textbf{实际频率}：在一批题目里，实际做对了 70\%。
\end{enumerate}

好的模型应该让这两个数接近。\hl{如果模型说 0.9，实际只有 0.6，这叫“校准不好”。}
这是第 5 章的一个核心发现。

\subsection{指数与对数}

\subsubsection{指数}

$e^x$ 表示 $e$ 的 $x$ 次方，其中 $e \approx 2.71828$，是一个特殊的常数。

\subsubsection{sigmoid 函数}

\begin{equation}
\sigma(x) = \frac{1}{1 + e^{-x}}
\end{equation}

\begin{explainbox}{符号说明}
\begin{itemize}[leftmargin=1.4em,itemsep=1pt]
\item $x$：任意实数（可正可负）。
\item $\sigma(x)$：输出范围是 $(0,1)$，正好可以当概率用。
\item 特性：$x$ 很大时输出接近 1；$x$ 很小时输出接近 0；$x=0$ 时输出 $0.5$。
\end{itemize}
\end{explainbox}

\textbf{通俗理解}：sigmoid 是一个“压缩器”。不管输入多大或多小，
它都会把结果压到 0 到 1 之间。所以它常被用来把“分数”变成“概率”。

第 4 章的平台评分公式最后一步就用了它。

\subsubsection{对数}

$\log(x)$ 表示“$e$ 的多少次方等于 $x$”。

\textbf{例}：$\log(1) = 0$，$\log(2.718) \approx 1$，$\log(7.389) \approx 2$。

\begin{explainbox}{为什么要用对数}
对数的作用是\hl{压制大数值}。比如：
\begin{itemize}[leftmargin=1.4em,itemsep=1pt]
\item 一个人参与 10 次，$\log(1+10) = 2.40$
\item 一个人参与 1000 次，$\log(1+1000) = 6.91$
\end{itemize}
参与次数差了 100 倍，但取对数后只差了 2.9 倍。
这样可以防止“刷次数的人”获得过大的权重。第 6 章的共识权重就用了这个技巧。
\end{explainbox}

\subsection{黄金比}

\begin{equation}
\varphi = \frac{1 + \sqrt{5}}{2} \approx 1.618
\end{equation}

\begin{explainbox}{符号说明}
\begin{itemize}[leftmargin=1.4em,itemsep=1pt]
\item $\varphi$：希腊字母 phi，读音“斐”。本文中表示黄金比。
\item $\sqrt{5}$：5 的平方根，约 $2.236$。
\item 它有个特殊的性质：$\varphi^2 = \varphi + 1$，而且 $1/\varphi = \varphi - 1 \approx 0.618$。
\end{itemize}
\end{explainbox}

\textbf{在本文中的用途}：$\varphi$ 是系统里“所有阈值和衰减率”的来源。
具体用法见第 3 章。

\subsection{哈希}

\term{哈希}是一种“指纹”算法。你把任意长度的内容丢进去，
它输出一串固定长度的乱码。

\begin{equation}
H(\text{内容}) = \text{一串固定长度的十六进制数}
\end{equation}

\begin{explainbox}{符号说明}
\begin{itemize}[leftmargin=1.4em,itemsep=1pt]
\item $H$：哈希函数（Hash function）。
\item 输入：任意内容（一段文字、一个数字、一个文件）。
\item 输出：例如 \sym{a3f5c9...}（通常 64 个十六进制字符）。
\item 三个关键性质：\textbf{(1) 同样的输入永远得到同样的输出}；
      \textbf{(2) 输入变一点点，输出完全变样}；\textbf{(3) 无法从输出反推出输入}。
\end{itemize}
\end{explainbox}

\textbf{通俗理解}：像人的指纹——能证明“是你”，但指纹本身不能还原出你这个人。

\subsubsection{一个真实存在的陷阱：低熵值哈希可以被暴力破解}

这一点非常重要，第 7 章会详细讲。简单说：

\begin{quote}
\textbf{“正确率 62\%” 这个值的可能性只有 101 种（0\% 到 100\%）。}\\
所以就算你把它哈希了，别人只要枚举 101 次就能反推出原值。
\end{quote}

\hl{解决办法是加盐}：把数值和一个随机数拼在一起再哈希。

\begin{equation}
C = H(\text{数值} \;\|\; \text{随机数})
\end{equation}

\begin{explainbox}{符号说明}
\begin{itemize}[leftmargin=1.4em,itemsep=1pt]
\item $C$：承诺值（Commitment）。上链的就是它。
\item $\|$：拼接符号，表示把两段内容接在一起。
\item \textbf{随机数}（nonce）：本地生成的、别人猜不到的随机串，只保存在自己手里。
\item 验证时，出示“数值 + 随机数”，对方重算一次哈希，对得上就证明“我没改过数据”。
\item 不验证时，别人就算知道 $C$，也推不出数值——因为随机数有无数种可能。
\end{itemize}
\end{explainbox}

\newpage
% ============================================================

\section{全部符号总索引}
% ============================================================

\begin{longtable}{@{}p{2.3cm}p{2.1cm}p{9.9cm}@{}}
\toprule
\textbf{符号} & \textbf{读法} & \textbf{含义与取值范围} \\
\midrule
\endhead
$\theta$ & 西塔 & 融合角。$[0,\pi]$。0 = 完全区分（清醒），$\pi$ = 完全融合（迷糊） \\
$\phi$ & 斐 & 认知相位。$[0,2\pi]$。标记历史路径，不影响概率大小 \\
$\gamma$ & 伽马 & 认知惯性。$[0,1]$。越大越固执 \\
$\alpha$ & 阿尔法 & (1) 颜色权重；(2) 冲击幅度。看上下文 \\
$\beta$ & 贝塔 & 文字权重，默认 0.30 \\
$\delta$ & 德尔塔 & 生理（HRV）权重，动态 \\
$\eta$ & 伊塔 & 学习率，每步 $\theta$ 的更新幅度。\textbf{文档默认 0.3，实验建议 0.02} \\
$\xi$ & 克西 & 场景切换惩罚。\textbf{文档默认 0.15，实验建议 0} \\
$\rho$ & 柔 & 情绪惯性 \\
$\varphi$ & 斐（大写） & 黄金比，$\approx 1.618$ \\
$\pi$ & 派 & 圆周率，$\approx 3.14159$ \\
$\sigma$ & 西格玛 & (1) 标准差；(2) sigmoid 函数。看上下文 \\
$\Delta$ & 德尔塔（大写） & 增量 / 变化量。$\Delta t$ = 时间步长 \\
$\mu$ & 缪 & 平均值（本文用横线表示） \\
$\bar{x}$ & x-bar & $x$ 的平均值 \\
$|D\rangle$ & D 态 & 区分态（清晰） \\
$|F\rangle$ & F 态 & 融合态（模糊） \\
$|\psi\rangle$ & 普西 & 认知状态向量 \\
$i$ & —— & 虚数单位，$i^2=-1$。仅作数学记号 \\
$H(\cdot)$ & H & 哈希函数 \\
$\|$ & —— & 拼接符号 \\
$\sum$ & 西格玛（求和） & 求和 \\
$\prod$ & 派（连乘） & 连乘 \\
$\log$ & —— & 自然对数（以 $e$ 为底） \\
$e$ & —— & 自然常数，$\approx 2.71828$ \\
$\cos, \sin$ & —— & 余弦、正弦 \\
$\arccos$ & —— & 反余弦 \\
$P(\cdot)$ & —— & 概率，取值 $[0,1]$ \\
$P(L_0)$ & —— & BKT 的先验掌握概率 \\
$P(T)$ & —— & BKT 的学习转移概率 \\
$P(G)$ & —— & BKT 的猜测概率 \\
$P(S)$ & —— & BKT 的失误概率 \\
$\hat{p}$ & p-hat & 模型预测的概率（帽子表示“估计值”） \\
$y$ & —— & 实际结果（答对 1，答错 0） \\
$N$ & —— & 样本数量 \\
$B_{i,t}$ & —— & 用户 $i$ 在时刻 $t$ 的 Being 记录 \\
$B^{\text{group}}$ & —— & 群体聚合记录 \\
$c$ & —— & 情绪色相（颜色） \\
$e(c)$ & —— & 颜色的五维编码向量 \\
$M_i$ & —— & 心智化综合得分 \\
$W$ & —— & 心智化窗口是否打开 \\
$\theta_{\text{th}}$ & $\theta$-th & 临界融合角，$\approx 0.9046$ 弧度 $\approx 51.8^\circ$ \\
$r_1 \ldots r_5$ & —— & 漏斗各环节转化率 \\
\bottomrule
\end{longtable}

\newpage

% ============================================================

\section{第二期新增的符号索引}
% ============================================================

\begin{longtable}{@{}p{2.6cm}p{2.2cm}p{9.8cm}@{}}
\toprule
\textbf{符号} & \textbf{读法} & \textbf{含义} \\
\midrule
\endhead
$\lambda(t)$ & 兰姆达 & 瞬时强度。$t$ 时刻事件发生的急迫程度 \\
$\mu$ & 缪 & 基线强度。没有历史事件时的发生率 \\
$\alpha$ & 阿尔法 & 激发增量。每个历史事件抬高多少强度 \\
$\beta$ & 贝塔 & 衰减速率。激发效应衰减多快 \\
$n=\alpha/\beta$ & --- & \textbf{分支比}。一个事件平均引发多少个后续事件 \\
$t_j$ & --- & 第 $j$ 个历史事件的时刻 \\
$\Delta t$ & 德尔塔 t & 时间间隔 \\
$\Lambda(t)$ & 大写兰姆达 & 累计强度（补偿器），$\int_0^t\lambda(s)\,ds$ \\
$\tau_i$ & 陶 & 重标间隔，$\Lambda(t_i)-\Lambda(t_{i-1})$ \\
$\tau_{\text{gap}}$ & 陶 gap & 会话切分阈值，默认 30 分钟 \\
$m(h)$ & --- & 昼夜调制函数 \\
$\mu_0$ & --- & 昼夜基线的整体水平 \\
$\phi(\Delta t)$ & 斐 & 核函数，描述“一个事件对后续影响多大” \\
$\chi^2$ & 卡方 & 衡量实际分布与理论分布差多远 \\
$k$ & --- & Weibull 形状参数。$k<1$ 表示间隔聚簇 \\
AIC & --- & 模型综合评价，越小越好 \\
KS 统计量 & --- & 累积分布的最大垂直距离，越小越好 \\
$p$ 值 & --- & 显著性。$p>0.05$ 通常视为通过检验 \\
Exp(1) & --- & 标准指数分布，平均单位时间 1 次且完全随机 \\
\bottomrule
\end{longtable}

\newpage

% ============================================================

\section{第三期新增的符号索引}
% ============================================================

\begin{longtable}{@{}p{2.8cm}p{2.2cm}p{9.6cm}@{}}
\toprule
\textbf{符号} & \textbf{读法} & \textbf{含义} \\
\midrule
\endhead
$\alpha$（幂律） & 阿尔法 & 幂律指数。$P(k) \propto k^{-\alpha}$ \\
$\chi^2$ & 卡方 & 实际分布与理论分布的距离度量 \\
$m(h)$ & --- & 相对活跃度调制函数，平均值归一化为 1 \\
$Z$ & --- & 归一化常数，$\frac{1}{P}\int_0^P \exp(s)\,dh$ \\
$s(h)$ & --- & 对数线性调制的指数部分（傅里叶级数） \\
$R$（会符号） & --- & 合向量长度，作息集中度，取值 $[0,1]$ \\
$R_{24}$ & --- & \textbf{日内强度}（Rayleigh）。时间点是否固定，取值 $[0,1]$ \\
$R_{168}$ & --- & \textbf{周内强度}（Rayleigh）。是否集中在周内特定位置 \\
$B$（节律平衡） & --- & $R_{168}/R_{24}$。$B$ 小 = 每日型；$B$ 大 = 间歇型 \\
$q_{\text{量}}$ & --- & 量偏离：今日活动量的\textbf{拟泊松偏差残差} \\
$q_{\text{时}}$ & --- & 时偏离：今日时相与个人基线的圆形距离（以 6 小时为标度）\\
$q_{\text{型}}$ & --- & 型偏离：今日小时轮廓与基线轮廓的\textbf{余弦距离} \\
$\phi$（过散布） & 斐 & 拟泊松的过散布因子 $=$ 方差$/ $均值，$\ge 1$ \\
trait & 特质 & \textbf{稳定}的个体属性（性格），应预测节律的\textbf{水平} \\
MetaConfidence & --- & \textbf{综合可信度}，$=1-\sum_k\beta_k\,\mathrm{alert}_k$，越大越可信 \\
$\beta_k$ & 贝塔 & 第 $k$ 个元元检测器的权重，$\sum_k\beta_k=1$ \\
$\mathrm{alert}_k$ & --- & 第 $k$ 个检测器是否报警，取 $0$ 或 $1$ \\
$s_k$ & --- & 研究评分里第 $k$ 个维度的\textbf{得分}，$s_k=1-\mathrm{alert}_k$ \\
$w_k$ & --- & 研究评分里第 $k$ 个维度的权重，$\sum_k w_k=1$ \\
$M_{\text{meta}}$ & --- & \textbf{弱测量算符}，$=I-\eta(1-\text{MC})P_{\text{模糊}}$ \\
$I$ & --- & 单位算符（什么都不做）\\
$\eta$（弱测量） & 伊塔 & 弱测量强度系数，L2 取 $0.01$ \\
$P_{\text{模糊}}$ & --- & 投影到模糊态的算符 \\
$V_{\text{veto}}$ & --- & \textbf{一票否决}的封顶值 \\
$\min$ & --- & 取较小值，用于\textbf{不可补偿}的护栏 \\
$\text{AUC}_{\text{meta}}$ & --- & \textbf{元认知敏感性}。标记能否区分错与对，$0.5$ 为无效 \\
$\text{bias}_{\text{meta}}$ & --- & \textbf{元认知偏差} $=\bar M-\bar E$。$>0$ 过度谨慎，$<0$ 过度自信 \\
$H,\ FA$ & --- & 命中（标记且错）与虚报（标记但对）的计数 \\
$\bar M,\ \bar E$ & --- & 标记率与错误率 \\
$d(\text{AUC})$ & --- & 自评可信度\textbf{折扣函数}，把 AUC 映射到 $[0.3,1.0]$ \\
$\text{MC}_{\text{adj}}$ & --- & 折减后的 MetaConfidence $=\text{MC}_{\text{meta}}\times d(\text{AUC})$ \\
JOL & --- & Judgment of Learning。测试前让学生预测自己会做对多少 \\
state & 状态 & \textbf{临时}的状态，影响节律的\textbf{波动}（偏离度） \\
$\rho$ & 柔 & Spearman 秩相关系数 \\
$d$ & --- & Cohen's $d$，两组均值的标准化差 \\
$N$ & --- & 样本量（学生人数） \\
power & 功效 & 真实效应存在时能检出它的概率 \\
$\alpha$（显著水平） & 阿尔法 & 假阳性率上限，习惯取 0.05 \\
Bonferroni & --- & 多重比较校正：$\alpha / (\text{检验次数})$ \\
AIC & --- & $-2\log L + 2k$，越小越好 \\
$k$（AIC 式） & --- & 参数个数 \\
H1--H8 & --- & 预注册登记的 8 条假设 \\
T1--T3 & --- & 数据采集的三个精度档位 \\
\bottomrule
\end{longtable}

\newpage

% ============================================================

\section{零基础符号与公式词典}
\begingroup\small
% ============================================================

\subsection{元认知}

\begin{longtable}{@{}p{3.0cm}p{4.8cm}p{6.4cm}@{}}
\toprule
\textbf{符号} & \textbf{读法} & \textbf{含义} \\
\midrule
$n$ & en & 测试总题数 \\
$y$ & wai & 实际答对题数 \\
$\hat y$ & y-hat & 预测答对题数 \\
$p=y/n$ & p & 实际正确率 \\
$\hat p=\hat y/n$ & p-hat & 预测正确率 \\
$b=\hat p-p$ & bias & 校准偏差；$>0$ 过度自信，$<0$ 过度保守 \\
$a=|\hat p-p|$ & absolute error & 校准误差，只看大小 \\
$\mathrm{Brier}$ & 布里尔分数 & 概率预测的评分，越小越好 \\
$\mathrm{AUC}$ & area under curve & 区分能力；$0.5$ 无效，$1.0$ 完美 \\
$d'$ & d prime & 信号检测论的区分力 \\
$\mathrm{MC}$ & metameta confidence & 对判断本身的信心 \\
\bottomrule
\end{longtable}

\subsection{犹豫与外部脚手架}

\begin{longtable}{@{}p{3.0cm}p{4.8cm}p{6.4cm}@{}}
\toprule
\textbf{符号} & \textbf{读法} & \textbf{含义} \\
\midrule
$h$ & hesitation & 犹豫程度（响应时间、修改次数等） \\
$m_{\text{internal}}$ & m-internal & 内部监控信号 \\
$m_{\text{report}}$ & m-report & 报告出来的信心/JOL \\
$E$ & external input & 外部提示、反馈、暗示 \\
$g$ & 映射函数 & 从内部信号到报告的映射 \\
\bottomrule
\end{longtable}

\subsection{耦合基线}

\begin{longtable}{@{}p{3.0cm}p{4.8cm}p{6.4cm}@{}}
\toprule
\textbf{符号} & \textbf{读法} & \textbf{含义} \\
\midrule
$x_t$ & x-t & 第 $t$ 期学生状态 \\
$u_t$ & u-t & 第 $t$ 期学生行为 \\
$e_t$ & e-t & 第 $t$ 期环境 \\
$\theta$ & theta & 慢参数（trait） \\
$f$ & 耦合函数 & 状态转移函数 \\
$x^*$ & x-star & 基线/稳定点 \\
$\varepsilon_t$ & epsilon-t & 噪声 \\
\bottomrule
\end{longtable}

\subsection{节律与习惯}

\begin{longtable}{@{}p{3.0cm}p{4.8cm}p{6.4cm}@{}}
\toprule
\textbf{符号} & \textbf{读法} & \textbf{含义} \\
\midrule
$\text{peak}_i$ & peak & 学生 $i$ 的日内节律峰型 \\
$h_s$ & h-s & 会话 $s$ 的开始小时 \\
$\mathrm{fit}_s$ & fit & 安排与节律的匹配度，$+1$ 最匹配 \\
$\text{habit\_change}$ & habit change & 周内节奏改变量，$0$ 稳定，$1$ 全变 \\
$R_{24}$ & R-24 & 日内节律强度 \\
$R_{168}$ & R-168 & 周内节律强度 \\
$B=R_{168}/R_{24}$ & B & 每日型（小）vs 间歇型（大） \\
\bottomrule
\end{longtable}

\subsection{群体拓扑}

\begin{longtable}{@{}p{3.0cm}p{4.8cm}p{6.4cm}@{}}
\toprule
\textbf{符号} & \textbf{读法} & \textbf{含义} \\
\midrule
$G=(V,E)$ & graph & 班级互动图 \\
$C_i$ & clustering & 局部聚类系数，小圈子程度 \\
$Q$ & modularity & 模块度，班级分裂程度 \\
$B_i$ & bridge ratio & 桥接比例，跨圈连接 \\
$D_{\text{div}}$ & diversity & 群体多样性 \\
\bottomrule
\end{longtable}

\subsection{尺度}

\begin{longtable}{@{}p{3.0cm}p{4.8cm}p{6.4cm}@{}}
\toprule
\textbf{符号} & \textbf{读法} & \textbf{含义} \\
\midrule
$T$ & period & 回路周期 \\
$\lambda$ & lambda & 尺度因子 \\
$\alpha$ & alpha & 标度指数 \\
$\varphi$ & phi & 黄金比 $(1+\sqrt5)/2\approx1.618$ \\
\bottomrule
\end{longtable}

\newpage
% ============================================================
\endgroup


\section{统一符号词典}

\subsection{基础}

\begin{longtable}{@{}p{2.8cm}p{4.4cm}p{7.0cm}@{}}
\toprule
\textbf{符号} & \textbf{读法} & \textbf{含义} \\
\midrule
$n$ & en & 总题数或总样本数 \\
$y$ & wai & 实际结果，例如答对题数 \\
$p$ & p & 概率或正确率 \\
$\theta$ & theta & 认知态/能力参数 \\
$\alpha$ & alpha & 旋转角、标度指数或显著性水平 \\
$\varepsilon$ & epsilon & 噪声 \\
\bottomrule
\end{longtable}

\subsection{元认知}

\begin{longtable}{@{}p{2.8cm}p{4.4cm}p{7.0cm}@{}}
\toprule
\textbf{符号} & \textbf{读法} & \textbf{含义} \\
\midrule
$\hat p$ & p-hat & 预测正确率 \\
$b=\hat p-p$ & bias & 校准偏差；$>0$ 过度自信 \\
$a=|\hat p-p|$ & absolute error & 校准误差 \\
$\mathrm{AUC}$ & area under curve & 区分能力；$0.5$ 无效，$1.0$ 完美 \\
$d'$ & d prime & 信号检测论的区分力 \\
$\mathrm{Brier}$ & 布里尔分数 & 概率预测评分，越小越好 \\
$\mathrm{MC}$ & metameta confidence & 对判断本身的信心 \\
$h$ & hesitation & 犹豫程度 \\
\bottomrule
\end{longtable}

\subsection{耦合基线}

\begin{longtable}{@{}p{2.8cm}p{4.4cm}p{7.0cm}@{}}
\toprule
\textbf{符号} & \textbf{读法} & \textbf{含义} \\
\midrule
$x_t$ & x-t & 第 $t$ 期学生状态 \\
$u_t$ & u-t & 学生行为 \\
$e_t$ & e-t & 环境（任务、教师、同伴、排课） \\
$\theta$ & theta & 慢参数（trait） \\
$f$ & 耦合函数 & 状态转移函数 \\
$x^*$ & x-star & 基线/稳定点 \\
\bottomrule
\end{longtable}

\subsection{节律、习惯与群体}

\begin{longtable}{@{}p{2.8cm}p{4.4cm}p{7.0cm}@{}}
\toprule
\textbf{符号} & \textbf{读法} & \textbf{含义} \\
\midrule
$\mathrm{fit}$ & fit & 安排与节律匹配度，$+1$ 最匹配 \\
$\mathrm{habit\_change}$ & habit change & 周内节奏改变量，$0$ 稳定，$1$ 全变 \\
$R_{24}$ & R-24 & 日内节律强度 \\
$R_{168}$ & R-168 & 周内节律强度 \\
$C_i$ & clustering & 局部聚类系数，小圈子程度 \\
$Q$ & modularity & 模块度，班级分裂程度 \\
$B_i$ & bridge ratio & 桥接比例，跨圈连接 \\
\bottomrule
\end{longtable}

\subsection{尺度}

\begin{longtable}{@{}p{2.8cm}p{4.4cm}p{7.0cm}@{}}
\toprule
\textbf{符号} & \textbf{读法} & \textbf{含义} \\
\midrule
$T$ & period & 回路周期 \\
$\lambda$ & lambda & 尺度因子 \\
$\alpha$ & alpha & 标度指数 \\
$\varphi$ & phi & 黄金比 $(1+\sqrt5)/2\approx1.618$ \\
\bottomrule
\end{longtable}



\part{数据与基础设施层}

\section{数据怎么存}
% ============================================================

\subsection{一条个人记录长什么样}

系统把每一次记录存成一条 \term{Being 记录}：

\begin{equation}
B_{i,t} = \big\langle
\underbrace{\text{id}}_{\text{谁}},\;
\underbrace{\text{timestamp}}_{\text{何时}},\;
\theta_{i,t},\;
\phi_{i,t},\;
\gamma_{i,t},\;
\underbrace{c_{i,t}}_{\text{颜色}},\;
e(c_{i,t}),\;
\underbrace{\text{text}_{i,t}}_{\text{文字}},\;
\text{scenario}_{i,t}
\big\rangle
\label{p1:eq:being}
\end{equation}

\subsubsection{九个字段逐个解释}

\begin{longtable}{@{}p{2.2cm}p{2.4cm}p{10.2cm}@{}}
\toprule
\textbf{字段} & \textbf{符号} & \textbf{含义} \\
\midrule
标识 & $\text{id}$ & 全局唯一的编号，回答“这是谁” \\
时间戳 & $\text{timestamp}$ & 记录时间，回答“这是什么时候” \\
融合角 & $\theta_{i,t}$ & 范围 $[0, \pi]$。0 = 完全清醒，$\pi$ = 完全迷糊 \\
认知相位 & $\phi_{i,t}$ & 范围 $[0, 2\pi]$。标记“从哪条路过来” \\
认知惯性 & $\gamma_{i,t}$ & 范围 $[0, 1]$。越大越“固执”，越难被新信息改变 \\
情绪色相 & $c_{i,t}$ & 用户自标的颜色（如“红”、“蓝”） \\
情绪向量 & $e(c_{i,t})$ & 把颜色编成五个数字，方便计算 \\
反思文本 & $\text{text}_{i,t}$ & 用户写的话 \\
场景标签 & $\text{scenario}_{i,t}$ & 当前场景，例如“语法填空第 9 篇” \\
\bottomrule
\end{longtable}

\begin{explainbox}{一个重要的区别}
$\theta$ 是\hl{算出来的}，$c$ 是用户\hl{自己标的}。\\
一个是客观计算结果，一个是主观自报。两者可以互相验证——
如果一个人说自己“很平静”（蓝色），但算出来的 $\theta$ 很高（迷糊），
系统就会知道这里有矛盾。
\end{explainbox}

\subsection{群体记录}

把很多人的记录合并（脱敏后），得到群体记录：

\begin{equation}
B^{\text{group}}_{s,t} = \big\langle
s,\; \bar{\theta}_{s,t},\; \sigma_{\theta,s,t},\; \bar{\phi}_{s,t},\;
\bar{c}_{s,t},\; H_{s,t},\; C_{s,t},\; N_{s,t}
\big\rangle
\label{p1:eq:group}
\end{equation}

\begin{explainbox}{逐符号解释}
\begin{longtable}{@{}p{2.4cm}p{11.4cm}@{}}
$s$ & 场景标签（比如“早自习”） \\
$\bar{\theta}_{s,t}$ & 平均融合角。符号上的横线 $\bar{\;}$ 表示“平均值”。 \\
$\sigma_{\theta,s,t}$ & 希腊字母 sigma。标准差，衡量“大家差得多不多”。$\sigma$ 大 = 有人清醒有人迷糊 \\
$\bar{\phi}_{s,t}$ & 平均认知相位 \\
$\bar{c}_{s,t}$ & 平均情绪色相（群体情绪基调） \\
$H_{s,t}$ & 场景场强（这个场景的总体压力） \\
$C_{s,t}$ & 群体一致性（大家有多齐） \\
$N_{s,t}$ & 样本数量（有多少人） \\
\end{longtable}
\end{explainbox}

\begin{explainbox}{脱敏是什么意思}
“脱敏”就是把能识别到具体个人的信息去掉：
\begin{itemize}[leftmargin=1.4em,itemsep=1pt]
\item 删掉 ID；
\item 人数太少时不出结果（比如少于 5 人）；
\item 极端值合并成“其他”。
\end{itemize}
目的：\hl{任何人拿到这份数据，都不能反推出“某个具体的人状态不好”}。
\end{explainbox}

\subsection{教育场景的扩展字段}

在教育场景里，Being 记录额外挂一个“学业载荷”：

\begin{equation}
\begin{split}
\text{academic}_{i,t} = \big\langle\;&
\text{item\_id},\; \text{question\_no},\; \text{correct},\; \text{error\_type},\\
&\text{latency},\; \text{marked\_words},\; \text{marked\_sentences},\; \text{attempt\_no}
\;\big\rangle
\end{split}
\end{equation}

\begin{longtable}{@{}p{3.4cm}p{10.4cm}@{}}
\toprule
\textbf{字段} & \textbf{含义} \\
\midrule
\text{item\_id} & 哪一篇讲义，例如“语法填空第 9 篇” \\
\text{question\_no} & 第几题，范围 56–65 \\
\text{correct} & 对还是错 \\
\text{error\_type} & 错因分类：时态 / 语态 / 词性 / 搭配 / 从句 / 非谓语 \\
\text{latency} & 用了多长时间（秒） \\
\text{marked\_words} & 这个学生标记的重点单词 \\
\text{marked\_sentences} & 标记的重点句子 \\
\text{attempt\_no} & 第几次做这道题 \\
\bottomrule
\end{longtable}

\subsection{出域白名单：五个字段}

在校内模式下，学生数据保存在学校。只有这五项允许离开：

\begin{center}
\begin{tabular}{p{4.0cm}p{9.8cm}}
\toprule
\textbf{允许出域} & \textbf{说明} \\
\midrule
\texttt{item\_id} & 题目标识 \\
\texttt{question\_no} & 题号 \\
$N$ & 作答人数 \\
\texttt{correct\_rate} & 正确率 \\
\texttt{error\_distribution} & 错因分布 \\
\bottomrule
\end{tabular}
\end{center}

\textbf{绝不出域}（这些必须留在学校）：

\begin{equation}
\theta,\; \phi,\; \gamma,\; c,\; e(c),\; \text{text},\; \text{所有学号}
\end{equation}

\begin{explainbox}{为什么这条线要划在这里}
\begin{itemize}[leftmargin=1.4em,itemsep=1pt]
\item 正确率、错因分布这类数据，说“这个班这道题 70\% 做对”是教学信息；
\item 而 $\theta$、情绪、文本属于\textbf{心理状态信息}，一旦流出就可能被滥用；
\item 最危险的是“反推”：如果一个班只有 8 个人，光看“有 1 人情绪异常”很容易猜到是谁。
\end{itemize}
所以规则是：\hl{聚合到班级以上、且人数足够多时，学业数据可以出域；心智数据永不出域。}
\end{explainbox}

\newpage

% ============================================================

\section{换单位：会话建模}
% ============================================================

\subsection{转折点}

把三次尝试并排：

\begin{center}
\begin{tabular}{p{6.0cm}rp{4.6cm}}
\toprule
\textbf{加了什么} & \textbf{AIC} & \textbf{时间重标通过率} \\
\midrule
什么也不加（单尺度） & 基准 & 36.7 \% \\
+ 第二个时间尺度 & 略好 & 36 \% \\
+ 昼夜周期 & 明显好 & 36.7 \% \\
\bottomrule
\end{tabular}
\end{center}

\textbf{两个假设都让模型拟合得更好，但没有一个能让时间重标通过。}

\begin{explainbox}{这个现象指向什么}
当模型越改越复杂、却仍然通不过检验时，问题通常在\hl{建模对象的定义}上，
而不在模型的容量上。\\[4pt]
换句话说：\textbf{不是模型不够强，是“什么算一个事件”选错了。}
\end{explainbox}

\subsection{新假设：基本单位应该是“一次会话”}

\begin{quote}
学生行为不是“自激点过程”，而是\textbf{“离散学习会话的到达过程”}。
\end{quote}

拆开就是：

\begin{enumerate}[leftmargin=2em]
\item 学生坐下来学习 $\to$ 30 分钟内密集作答 20~30 题（\textbf{会话内}）；
\item 然后离开 $\to$ 可能隔几小时、也可能隔几天（\textbf{会话间}）；
\item \textbf{会话的发生由作息、课表决定，不是由上一个作答激发的。}
\end{enumerate}

\subsection{会话怎么定义}

\begin{equation}
\text{相邻两次事件间隔} > \tau_{\text{gap}} \;\Longrightarrow\; \text{换了一个会话}
\end{equation}

\begin{explainbox}{符号说明}
\begin{itemize}[leftmargin=1.4em,itemsep=1pt]
\item $\tau_{\text{gap}}$（读作“tau gap”）：切分阈值，默认取 \textbf{30 分钟}。
\item 意思是：如果两道题之间隔了超过 30 分钟，就认为“这不是同一次学习”。
\item 这个定义简单到近乎粗糙——但后面会看到它对结论不敏感。
\end{itemize}
\end{explainbox}

\subsection{会话长什么样}

对 41 名学生切分后：

\begin{center}
\begin{tabular}{p{6.6cm}p{6.6cm}}
\toprule
\textbf{指标} & \textbf{中位数} \\
\midrule
每个人的会话数 & \textbf{14} \\
每次会话的作答数 & \textbf{13} \\
每次会话时长 & \textbf{15.4 分钟} \\
\textbf{会话内作答速率} & \textbf{51 次 / 小时} \\
\bottomrule
\end{tabular}
\end{center}

\textbf{通俗读数}：学生一次学习通常持续 15 分钟、做 13 道题，
速率约每分钟 0.85 题。会话之间可能隔几小时甚至几天。

\subsection{核心结果}

\subsubsection{配对比较（24 名学生）}

\begin{center}
\begin{tabular}{p{5.4cm}rr}
\toprule
\textbf{指标} & \textbf{事件层面} & \textbf{会话层面} \\
\midrule
分支比（中位） & 0.929 & \hl{0.543} \\
\textbf{分支比更低的用户数} & --- & \hl{\textbf{24 / 24}} \\
\bottomrule
\end{tabular}
\end{center}

\textbf{每一名学生的分支比都下降了，没有一个例外。}
这不是抽样波动，是系统性变化。

\subsubsection{加入时间重标检验（10 名会话数足够的学生）}

\begin{center}
\begin{tabular}{p{5.4cm}rr}
\toprule
\textbf{指标} & \textbf{事件层面} & \textbf{会话层面} \\
\midrule
分支比（中位） & 0.917 & \hl{0.543} \\
KS 统计量（中位） & 0.122 & 0.135 \\
\textbf{KS 通过率} & \warn{\textbf{10 \%}} & \hl{\textbf{100 \%}} \\
分支比更低的用户数 & --- & \hl{10 / 10} \\
\bottomrule
\end{tabular}
\end{center}

\begin{explainbox}{最关键的一行}
\textbf{KS 通过率从 10\% 升到 100\%。}\\[4pt]
也就是说，“会话是基本事件”这个模型，\hl{通过了时间重标检验}---
而事件层面的 Hawkes 做了三次改进都无法通过。
\end{explainbox}

\subsection{切分阈值的敏感性}

会话怎么切，会不会影响结论？扫描 $\tau_{\text{gap}}$ 从 6 分钟到 4 小时：

\begin{center}
\begin{tabular}{rrrr}
\toprule
$\tau_{\text{gap}}$ & \textbf{会话数中位} & \textbf{分支比中位} & \textbf{KS 通过率} \\
\midrule
0.10 h（6 分） & 20 & 0.492 & 93 \% \\
0.25 h（15 分） & 16 & 0.500 & 100 \% \\
\textbf{0.50 h（30 分）} & \textbf{14} & \textbf{0.497} & \textbf{100 \%} \\
1.00 h & 12 & 0.447 & 100 \% \\
2.00 h & 11 & 0.475 & 100 \% \\
4.00 h & 9 & 0.000 & 100 \% \\
\bottomrule
\end{tabular}
\end{center}

\begin{explainbox}{读数}
在 \textbf{6 分钟到 2 小时}这个很宽的范围内：
\begin{itemize}[leftmargin=1.4em,itemsep=1pt]
\item 分支比稳定在 0.45 ~ 0.50；
\item KS 通过率始终在 93\%~100\%。
\end{itemize}
\textbf{结论不依赖阈值的具体取值}---这是很强的稳健性证据。
\end{explainbox}

\subsection{三层结构终于闭合}

\begin{longtable}{@{}p{1.0cm}p{4.4cm}p{3.4cm}p{6.6cm}@{}}
\toprule
\textbf{层次} & \textbf{问题} & \textbf{答案} & \textbf{证据} \\
\midrule
{(1)}& 数据有结构吗？ & \hl{有} & 0.93 vs 泊松 0.09 \\
{(2)}& 是“单次作答自激”吗？ & \warn{不是} & 24/24 人换成会话后分支比下降 \\
{(3)}& 是“会话到达”吗？ & \hl{是} & 会话层面 KS 通过率 \textbf{100\%} \\
\bottomrule
\end{longtable}

\subsection{剩下的 0.54 是什么}

会话层面的分支比仍有 0.54，说明\textbf{会话之间还有一定的自激}，可能是：

\begin{itemize}[leftmargin=1.6em]
\item 一次学习之后更容易再学一次（习惯 / 惯性）；
\item 或者作息、周计划带来的规律性。
\end{itemize}

\textbf{0.54 是一个“中等自激”的合理读数}，与“学生会连着几天学习、
也会连着几天不学”的直觉一致。

\newpage

% ============================================================

\section{会话本身长什么样}
% ============================================================

上一期把单位换成会话就停了，没有细看会话本身。
这一节先把它看清楚。

\subsection{会话有多大：幂律}

对 653 次会话（41 名学生）的大小做分布拟合：

\begin{center}
\begin{tabular}{p{4.0cm}rrr}
\toprule
\textbf{分布} & \textbf{参数} & \textbf{对数似然} & \textbf{AIC} \\
\midrule
\textbf{幂律} & $\alpha = 1.342$ & $-1929.8$ & \textbf{3861.5} \\
对数正态 & $\mu=2.23,\ \sigma=1.19$ & $-2497.7$ & 4999.4 \\
几何分布 & $p = 0.057$ & $-2501.4$ & 5004.7 \\
\bottomrule
\end{tabular}
\end{center}

\textbf{幂律以约 1,140 的 AIC 优势胜出}——这不是微弱差别。

\begin{explainbox}{幂律长什么样}
幂律就是 $P(k) \propto k^{-\alpha}$：\textbf{大部分会话很小，但偶尔出现极长的会话}，
而且这种长尾不是异常值，是分布的固有部分。\\[4pt]
拟合出 $\alpha \approx 1.34$，属于长尾很重的一档。
通俗说：\textbf{少数几次学习，占了总题量的大部分。}
\end{explainbox}

\textbf{产品含义}：如果按“平均每次做 10 题”来设计提醒和任务，
会完全错过那几次决定性的长会话。

\subsection{会话内有疲劳吗：没有，反而有热身}

把会话内部的作答间隔按相对位置分箱（0 = 开始，1 = 结束）：

\begin{center}
\begin{tabular}{rrr}
\toprule
\textbf{相对位置} & \textbf{作答间隔中位数（秒）} & \textbf{样本数} \\
\midrule
0.1（开始） & \textbf{42.19} & 2,259 \\
0.3 & 41.40 & 2,087 \\
0.5 & 38.60 & 2,079 \\
0.7 & 38.63 & 2,093 \\
0.9（结束） & \textbf{38.53} & 1,908 \\
\bottomrule
\end{tabular}
\end{center}

\begin{explainbox}{这不是疲劳，是热身}
如果存在疲劳，应该看到间隔\textbf{变长}（越做越慢）。\\
实测是\textbf{变短}（42.2 秒降到 38.5 秒，降幅约 9\%）。\\
说明学生需要一个“进入状态”的过程：前几十秒在回忆规则、适应题型，之后速度稳定下来。
\end{explainbox}

\textbf{产品含义}：会话开头的题目\textbf{不应该太难}——那正是学生还在热身的时候。

\subsection{有没有周周期：有，而且很强}

把全会话开始时刻映射到“周内小时”（0\textasciitilde 167）：

\begin{center}
\begin{tabular}{lrr}
\toprule
\textbf{检验} & $\chi^2$ 统计量 & \textbf{自由度} \\
\midrule
星期分布 & \textbf{59} & 6 \\
周内 168 小时分布 & \textbf{781} & 167 \\
\bottomrule
\end{tabular}
\end{center}

若完全均匀，$\chi^2$ 应约等于自由度。实测远大于临界值，\hl{周结构显著存在}。

\newpage
% ============================================================


\part{Tender 建模核心}

\section{心智化系统的核心公式}
% ============================================================

这一章回答一个问题：\hl{怎么用数据表示“一个人的认知状态”？}

\subsection{认知态：一个球面上的点}

系统用一个叫 \term{Bloch 球}（布洛赫球）的几何对象来表示认知状态。
不要被名字吓到，它就是一个普通的球。

想象一个地球仪：

\begin{itemize}[leftmargin=1.6em]
\item 球面上每一个点，代表一种认知状态；
\item 用两个角度就能定位一个点：
  \begin{itemize}
  \item $\theta$（融合角）：决定\hl{纬度}——点在上半球还是下半球；
  \item $\phi$（认知相位）：决定\hl{经度}——点在东半球还是西半球。
  \end{itemize}
\end{itemize}

\subsubsection{认知态的数学写法}

\begin{equation}
|\psi\rangle = \cos\!\left(\frac{\theta}{2}\right)|D\rangle
\;+\; e^{i\phi}\sin\!\left(\frac{\theta}{2}\right)|F\rangle
\end{equation}

\begin{explainbox}{逐符号解释}
\begin{longtable}{@{}p{2.6cm}p{11.2cm}@{}}
$\psi$ & 希腊字母 psi，读音“普西”。$|\psi\rangle$ 表示“这个人的认知状态”。 \\
$|D\rangle$ & 区分态（Differentiated）。代表“思路清晰、能分清事实与感受”。 \\
$|F\rangle$ & 融合态（Fused）。代表“纠缠在一起、分不清”。 \\
$\theta$ & 融合角，范围 $[0, \pi]$。$\theta=0$ 表示完全区分，$\theta=\pi$ 表示完全融合。 \\
$\phi$ & 认知相位，范围 $[0, 2\pi]$。表示“走的是哪条路径”。 \\
$e^{i\phi}$ & 复数相位因子。本文只用到它的一个性质：模长恒为 1，不影响概率大小。 \\
$i$ & 虚数单位，$i^2 = -1$。这里只是数学记号，\warn{不代表真实物理量子效应}。 \\
\end{longtable}
\end{explainbox}

\begin{explainbox}{用一句话理解}
$\theta$ 越小，$|D\rangle$（清晰态）的成分越多，人越清醒；\\
$\theta$ 越大，$|F\rangle$（融合态）的成分越多，人越迷糊。
\end{explainbox}

\subsubsection{从认知态推出“答对概率”}

这是本文最关键的映射之一。在 Bloch 球表示下：

\begin{equation}
P(\text{区分态}) = \cos^2\!\left(\frac{\theta}{2}\right)
\label{p1:eq:bloch-map}
\end{equation}

\begin{explainbox}{符号说明与直觉}
\begin{itemize}[leftmargin=1.4em,itemsep=1pt]
\item 左边：处于“清晰态”的概率，取值范围 $(0,1)$。
\item 右边：$\theta/2$ 的余弦的平方。
\item \textbf{代入验证}：
  \begin{itemize}
  \item $\theta = 0$ 时，$\cos^2(0) = 1$，即 100\% 清晰；
  \item $\theta = \pi$ 时，$\cos^2(\pi/2) = 0$，即 0\% 清晰；
  \item $\theta = \pi/2$ 时，$\cos^2(\pi/4) = 0.5$，即一半一半。
  \end{itemize}
\item \textbf{在第 5 章的实验中}，这个量被直接当作“答对概率”使用。
\end{itemize}
\end{explainbox}

\subsection{四源合成：$\theta$ 是怎么算出来的}

单个来源不够可靠，所以系统用四个来源加权合成：

\begin{equation}
\theta_{i,t} = \alpha_\theta(t)\,\theta^{\text{color}}_{i,t}
\;+\; \beta_\theta\,\theta^{\text{text}}_{i,t}
\;+\; \gamma_\theta\,\theta^{\text{seq}}_{i,t}
\;+\; \delta_\theta(t)\,\theta^{\text{HRV}}_{i,t}
\label{p1:eq:four-source}
\end{equation}

\begin{explainbox}{逐符号解释}
\begin{longtable}{@{}p{3.4cm}p{10.4cm}@{}}
$\theta_{i,t}$ & 用户 $i$ 在时刻 $t$ 的融合角（最终结果）。 \\
$\theta^{\text{color}}$ & 由用户自标的颜色推算出的融合角。 \\
$\theta^{\text{text}}$ & 由用户写的文字推算出的融合角。 \\
$\theta^{\text{seq}}$ & 由历史行为序列推算出的融合角。 \\
$\theta^{\text{HRV}}$ & 由心率变异性（生理信号）推算出的融合角。 \\
$\alpha_\theta(t)$ & 颜色权重，动态调整，默认范围 $0.10$–$0.35$。 \\
$\beta_\theta$ & 文字权重，默认 $0.30$。 \\
$\gamma_\theta$ & 序列权重，默认 $0.15$。 \\
$\delta_\theta(t)$ & 生理权重，动态调整，范围 $0$–$0.45$。 \\
\multicolumn{2}{@{}l@{}}{\textbf{约束}：四个权重之和恒为 1，即 $\alpha+\beta+\gamma+\delta = 1$。} \\
\end{longtable}
\end{explainbox}

\begin{explainbox}{为什么权重是动态的}
如果采集到的生理信号质量很差（比如手环没戴稳），系统会自动降低 $\delta_\theta$，
把权重让给更可靠的来源。这叫“按信号质量分配信任”。
\end{explainbox}

\subsection{序列分量：$\theta$ 怎么随时间变化}

这是第 5 章实验真正用到的公式（因为实验数据只有作答序列）：

\begin{equation}
\theta^{\text{seq}}_{i,t} = \theta_{i,t-1} \;+\; \eta \cdot \Delta t \cdot \text{趋势} \;+\; \xi \cdot \text{场景切换惩罚}
\label{p1:eq:theta-seq}
\end{equation}

\begin{explainbox}{逐符号解释}
\begin{longtable}{@{}p{3.2cm}p{10.6cm}@{}}
$\theta^{\text{seq}}_{i,t}$ & 用户 $i$ 在时刻 $t$ 的序列分量。 \\
$\theta_{i,t-1}$ & 上一个时刻的融合角。 \\
$\eta$ & 希腊字母 eta，读音“伊塔”。\hl{学习率}，即每一步 $\theta$ 变化多大。\textbf{文档默认 0.3}。 \\
$\Delta t$ & 时间步长。相邻两次作答时取 1。 \\
趋势 & 最近的变化方向。本文取：\textbf{答对 = $-1$，答错 = $+1$}。 \\
$\xi$ & 希腊字母 xi，读音“克西”。\textbf{场景切换惩罚}，文档默认 0.15。 \\
场景切换惩罚 & 当学生从一个知识点换到另一个知识点时，加一点不确定性。 \\
\end{longtable}
\end{explainbox}

\textbf{通俗理解}：

\begin{quote}
“如果你刚才就很迷糊，现在大概率也不会突然清醒。”\\
答对一次，$\theta$ 往下走一点（更清醒）；答错一次，$\theta$ 往上走一点（更迷糊）。
\end{quote}

\subsection{冲击响应饱和}

如果一道题错了，$\theta$ 应该跳多少？系统用了一个“饱和函数”：

\begin{equation}
\alpha = \frac{0.5 \cdot s}{0.618 + |s|}
\label{p1:eq:saturation}
\end{equation}

\begin{explainbox}{逐符号解释}
\begin{itemize}[leftmargin=1.4em,itemsep=1pt]
\item $s$：冲击信号强度。取绝对值 $|s|$。
\item $0.5$：分子系数。
\item $0.618$：分母常数，等于 $1/\varphi$（黄金比的倒数）。
\item $\alpha$：实际冲击幅度。
\end{itemize}
\end{explainbox}

\textbf{为什么要“饱和”}：如果冲击不设上限，一次极端事件就能把状态推到极端。
饱和函数的作用是\hl{让冲击有一个天花板}。

\textbf{代入计算}：当 $|s| = 1$ 时，
\begin{equation}
\alpha = \frac{0.5 \times 1}{0.618 + 1} = \frac{0.5}{1.618} = 0.309
\end{equation}

\subsection{临界纬线：什么时候算“该干预”}

系统画了一条线，$\theta$ 超过它就算“不适合干预”：

\begin{equation}
\theta_{\text{th}} = \arccos\!\left(\frac{1}{\varphi}\right) = \arccos(0.618) \approx 0.9046 \text{ 弧度} \approx 51.8^\circ
\label{p1:eq:theta-th}
\end{equation}

\begin{explainbox}{逐符号解释}
\begin{itemize}[leftmargin=1.4em,itemsep=1pt]
\item $\theta_{\text{th}}$：临界融合角（threshold），本文中读作“$\theta$ 阈值”。
\item $\arccos$：反余弦函数，是 $\cos$ 的反函数。已知余弦值，求对应的角度。
\item $1/\varphi = 0.618$：黄金比的倒数。
\item $0.9046$ 弧度：换算成角度约 $51.8^\circ$。
\end{itemize}
\end{explainbox}

\subsection{心智化窗口判据}

\begin{equation}
\text{Mentalizing\_Open} =
\begin{cases}
1, & \theta < \theta_{\text{th}} \;\text{且}\; \rho < \rho_{\text{th}} \\[4pt]
0, & \text{其他情况}
\end{cases}
\label{p1:eq:window}
\end{equation}

\begin{explainbox}{逐符号解释}
\begin{itemize}[leftmargin=1.4em,itemsep=1pt]
\item $\text{Mentalizing\_Open}$：心智化窗口是否打开。1 = 打开，0 = 关闭。
\item $\theta$：当前融合角。小于阈值表示“够清醒”。
\item $\rho$：希腊字母 rho，读音“柔”。\hl{情绪惯性}，值越大表示情绪越难平复。
\item $\rho_{\text{th}}$：情绪惯性的阈值。
\end{itemize}
\end{explainbox}

\begin{explainbox}{这就是整本书的核心}
当窗口打开时（$\theta$ 小、情绪稳），系统才推新知识；\\
当窗口关闭时，系统改做低负荷的事，或者只陪伴、不干预。\\
\hl{这一条判断，就是本系统区别于普通教育软件的全部要害。}
\end{explainbox}

\subsection{三端的窗口化身}

同一个判断逻辑，在三类对象上分别表现为：

\begin{center}
\begin{tabular}{p{2.1cm}p{4.6cm}p{6.0cm}}
\toprule
\textbf{对象} & \textbf{窗口问题} & \textbf{窗口关闭时的行为} \\
\midrule
学生端 & 现在适合学新知识，还是先巩固？ & 不推新知识点，改做低负荷练习 \\
教师端 & 现在适合讲新内容，还是先讲错题？ & 不讲新内容，先做错题干预 \\
销售端 & 现在适合推招生，还是先建信任？ & 不发招生内容，先发干货 \\
\bottomrule
\end{tabular}
\end{center}

\newpage

% ============================================================

\section{什么是“自激过程”}
% ============================================================

这一章开始进入这一期真正的主角。

\subsection{先说什么叫“点过程”}

\term{点过程}就是“一串事件发生的时间点”。

例如：一条视频发布后，点赞发生的时刻是
\texttt{第 3 分钟、第 5 分钟、第 5 分 20 秒、第 12 分钟、……}

这一串时间点就是一个点过程。

\begin{explainbox}{点过程和普通数据的区别}
普通数据是“一次一个值”（比如身高 175cm）。\\
点过程是“什么时候发生了一件事”。研究的是\textbf{时间和时间的间隔}，不是数值大小。
\end{explainbox}

\subsection{自激是什么意思}

\term{自激}（self-exciting）的直觉是：

\begin{quote}
\textbf{一件事发生之后，下一件事更容易发生。}
\end{quote}

例子：
\begin{itemize}[leftmargin=1.6em]
\item 地震：一次地震后余震更频繁（这是 Hawkes 过程的原始应用）；
\item 交易：一笔大单后常有更多交易；
\item 平台：有人说“一个赞会带来更多赞”。
\end{itemize}

\begin{explainbox}{为什么这件事重要}
如果“互动会引发更多互动”是真的，那运营策略就是\textbf{制造初始互动}。\\
如果它是假的（只是平台在推流），那策略就应该是\textbf{优化内容以通过平台筛选}。\\
\textbf{两种解释导致完全不同的做法。}
\end{explainbox}

\subsection{Hawkes 过程的三个符号}

\term{Hawkes 过程}是描述自激的标准数学模型。它有三个参数：

\begin{longtable}{@{}p{2.0cm}p{2.2cm}p{9.6cm}@{}}
\toprule
\textbf{符号} & \textbf{名称} & \textbf{含义} \\
\midrule
$\mu$ & 基线强度 & 没有历史事件时，单位时间的发生率 \\
$\alpha$ & 激发增量 & 每来一个事件，后续强度抬高多少 \\
$\beta$ & 衰减速率 & 激发效应衰减多快 \\
\bottomrule
\end{longtable}

\subsection{强度函数：核心公式}

\begin{equation}
\lambda(t) = \mu + \sum_{t_j < t} \alpha\, e^{-\beta (t - t_j)}
\label{p2:eq:intensity}
\end{equation}

\begin{explainbox}{逐符号解释}
\begin{longtable}{@{}p{2.2cm}p{11.6cm}@{}}
$\lambda(t)$ & $t$ 时刻的\hl{瞬时强度}——可以理解为“此刻事件发生的急迫程度” \\
$\mu$ & 基线强度（上面表格里的那个） \\
$\sum_{t_j<t}$ & 对所有“发生在 $t$ 之前”的事件求和 \\
$t_j$ & 第 $j$ 个历史事件的时刻 \\
$t - t_j$ & 距离那个事件过去了多久 \\
$e^{-\beta(t-t_j)}$ & \hl{衰减因子}。间隔越久，影响越小 \\
$\alpha$ & 每个历史事件贡献的强度 \\
\end{longtable}
\end{explainbox}

\textbf{通俗理解}：

\begin{quote}
“此刻的活跃程度 = 本来就有多活跃（$\mu$）
+ 之前每个事件留下的余温（$\alpha e^{-\beta\Delta t}$）”
\end{quote}

\subsection{分支比：判断会不会“自持”}

\begin{equation}
n = \frac{\alpha}{\beta}
\label{p2:eq:branching}
\end{equation}

\begin{explainbox}{分支比的三种情况}
\begin{center}
\begin{tabular}{p{2.6cm}p{10.2cm}}
\toprule
$n$ 的值 & 含义 \\
\midrule
$n < 1$ & \hl{平稳}。一次事件引发的后续影响会逐渐消失 \\
$n \to 1$ & \warn{接近临界}。一次事件可能引发很长的连锁反应 \\
$n > 1$ & \warn{发散}。过程会爆炸（现实中不可能持续） \\
\bottomrule
\end{tabular}
\end{center}
\textbf{直觉}：$n$ 就是“一个事件平均能引发多少个后续事件”。
\end{explainbox}

\newpage
% ============================================================

\section{第一次建模尝试：一个看起来很漂亮的结果}
% ============================================================

\subsection{用什么数据}

我用了 EdNet 数据集——一个公开的在线学习平台日志，里面有\textbf{毫秒级的时间戳}。

取其中 236 名学生的“作答”事件（respond），转换成小时，然后拟合 Hawkes。

\subsection{结果}

\begin{center}
\begin{tabular}{p{5.6cm}p{7.6cm}}
\toprule
\textbf{指标} & \textbf{数值} \\
\midrule
分析学生数 & 236 \\
\textbf{分支比中位数} & \textbf{0.941} \\
分支比均值 & 0.927 \\
分类为“接近临界”的学生 & 219 人 \\
分类为“强自持” & 16 人 \\
分类为“弱自持” & 1 人 \\
\bottomrule
\end{tabular}
\end{center}

\begin{explainbox}{看到这个结果的第一反应}
“分支比 0.94！\textbf{学生的作答几乎全部处于临界状态！}”\\
“临界状态意味着：一次作答能引发接近 1 次后续作答，”
“这是典型的自组织临界性，是个重大发现！”
\end{explainbox}

\warn{但这个结论是错的。}下一章讲怎么发现的。

\begin{explainbox}{先记住一条纪律}
得到一个漂亮结果时，第一件事不是写论文，而是问：
\begin{quote}
\textbf{“如果我把它打乱，还会得到同样的结果吗？”}
\end{quote}
这就是下一章的\term{安慰剂检验}。
\end{explainbox}

\newpage

% ============================================================

\section{分层模型：共享形状}
% ============================================================

修好上面两个问题后，才做成了这件事。

\subsection{为什么要共享}

每个用户单独估 7 个参数（3 个动力学 + 4 个傅里叶系数），
但每人只有十几个会话——\textbf{形状参数估不准}。

分层方案：把参数分两层。

\begin{center}
\begin{tabular}{p{3.6cm}p{5.4cm}p{5.0cm}}
\toprule
\textbf{层} & \textbf{参数} & \textbf{个数} \\
\midrule
\textbf{共享层}（所有人相同） & 傅里叶形状 $a_1,b_1,a_2,b_2$ & 4 \\
\textbf{个体层}（每人一套） & 幅度 $\mu_0$、自激 $\alpha$、$\beta$ & 3 / 人 \\
\bottomrule
\end{tabular}
\end{center}

总参数 $= 4 + 3N$。以 $N=21$ 为例：\textbf{67 个}，对比“各自 7 个”的 147 个。

\subsection{算法：交替优化}

一次性优化 67 维非凸问题很难。改成轮流：

\begin{enumerate}[leftmargin=2em]
\item 固定共享形状 → 每个用户的 $(\mu_0,\alpha,\beta)$ \textbf{互相独立}，可各自快速求解；
\item 固定个体参数 → 只优化 4 个共享系数；
\item 重复 1、2 直到收敛。
\end{enumerate}

实际跑 4 轮，似然变化：$-2205 \to -2177 \to -2173 \to -2173$（第 4 轮已收敛）。

\subsection{结果}

\begin{center}
\begin{tabular}{p{5.0cm}rrr}
\toprule
\textbf{模型} & \textbf{总对数似然} & \textbf{总 AIC} & \textbf{参数总数} \\
\midrule
无周期（3 / 人） & $-2198.7$ & 4523.5 & 63 \\
各自独立傅里叶（7 / 人） & $-2092.3$ & \textbf{4478.6} & 147 \\
\textbf{共享形状（4 共享 + 3 / 人）} & $-2173.1$ & \textbf{4480.2} & \textbf{67} \\
\bottomrule
\end{tabular}
\end{center}

\begin{explainbox}{核心读数}
\textbf{AIC 差 1.6，参数少 54\%。}\\
统计上 AIC 差小于 2 通常视为“无差别”。所以：\\
\textbf{共享形状用一半的参数，达到了和独立拟合一样的效果。}
\end{explainbox}

\subsection{拟合出的共享形状：一日双峰}

把共享系数转成一天内的相对活跃度 $m(h)$：

\begin{center}
\begin{tabular}{lrl}
\toprule
\textbf{时段} & \textbf{相对活跃度} & \\
\midrule
00 时 & 2.52 & 最高 \\
\textbf{13 时} & \textbf{1.78} & 次峰 \\
\textbf{06--07 时} & \textbf{0.13} & 谷底 \\
23 时 & 2.40 & 高 \\
\bottomrule
\end{tabular}
\end{center}

\textbf{峰值 2.56（约 24 时），谷值 0.13（6.5 时），峰谷比 20.2 倍。}

\begin{explainbox}{一个有意思的细节}
比较两个谐波的振幅：
\begin{center}
\begin{tabular}{lrr}
\toprule
谐波 & 周期 & 振幅 \\
\midrule
1 阶 & 24 小时 & 0.736 \\
\textbf{2 阶} & \textbf{12 小时} & \textbf{1.024} \\
\bottomrule
\end{tabular}
\end{center}
\textbf{12 小时周期的振幅比 24 小时还大}——意味着一日\textbf{两个}活跃高峰，
而不是“白天活跃、晚上安静”的单峰模式。\\
与实际一致：\textbf{下午（13 时）一个峰，深夜（23--24 时）另一个峰。}
\end{explainbox}

\subsection{产品价值：冷启动}

\begin{explainbox}{新用户只需要估 3 个参数}
共享形状已经有了，新用户只需要估 $\mu_0,\alpha,\beta$ 三个参数，
\textbf{十几个会话就能稳定}。\\
而独立拟合需要 7 个参数，同样数据量下估不准。\\[4pt]
\textbf{这是分层模型最直接的工程价值。}
\end{explainbox}

\newpage
% ============================================================

\section{从群体走向个体：MBTI 这条路}
% ============================================================

前面几节都在处理“群体的作息规律”。这一节开始处理\textbf{个体差异}。

\subsection{为什么想用 MBTI}

出发点很朴素：如果不同的人有不同的“节奏偏好”，
那给所有人推同样的学习安排就是浪费。

MBTI 和荣格八维提供了一个现成的语言体系来描述这种差异。

\subsection{但必须先说清楚它的证据基础}

\begin{longtable}{@{}p{3.6cm}p{10.2cm}@{}}
\toprule
\textbf{问题} & \textbf{具体表现} \\
\midrule
重测一致率低 & 同一人隔几周重测，得到同一“类型”的比例通常只有 50\%--75\% \\
不是双峰分布 & 四个维度上多数人落在中间，并不天然分成两类 \\
类型学不成立 & “16 型”分类框架的构念效度缺乏支持 \\
与五大人格的关系 & 四维与 Big Five 有中等相关，但\textbf{信息量不如 Big Five} \\
荣格八维 & \textbf{没有强心理测量学支持}，是解释性框架，不是测量工具 \\
\bottomrule
\end{longtable}

\begin{explainbox}{所以本项目的三条使用原则}
\begin{enumerate}[leftmargin=1.4em,itemsep=2pt]
\item \textbf{只用连续分数，不用类型}\\
      模型里存“外倾倾向 = 73 分”，不是“你是 E 型”。
\item \textbf{不宣称是“测量”}\\
      可以说“这个倾向帮你理解自己的学习节奏”，
      不能说“测出你是 INTJ”。
\item \textbf{必须能被证伪}\\
      登记具体假设，用数据检验。不通过就撤掉。
\end{enumerate}
\end{explainbox}

\subsection{把抽象维度变成可测指标}

MBTI 的维度本身测不了，但它的\textbf{行为表现}可以测。
所以每条假设都对应一个能从会话数据算出的指标：

\begin{longtable}{@{}p{1.0cm}p{2.6cm}p{3.4cm}p{1.4cm}p{5.4cm}@{}}
\toprule
\textbf{编号} & \textbf{性格维度} & \textbf{行为指标} & \textbf{方向} & \textbf{理由} \\
\midrule
H1 & J（判断端） & 作息集中度 $R$ & $+$ & 偏好计划的人作息更规整 \\
H2 & J（判断端） & 作息熵 & $-$ & 同上 \\
H3 & Si（内倾感觉） & 间隔变异 CV & $-$ & 习惯固定 \\
H4 & Ne（外倾直觉） & 间隔变异 CV & $+$ & 发散，间隔不稳 \\
H5 & S（感觉端） & 大小变异 CV & $-$ & 偏好稳定，任务量均匀 \\
H6 & N（直觉端） & 大小基尼 & $+$ & 容易一口气做很多 \\
H7 & E（外倾端） & 会话时长中位 & $-$ & 倾向短而频繁 \\
H8 & I（内倾端） & 会话时长中位 & $+$ & 倾向长而稀少 \\
\bottomrule
\end{longtable}

\subsubsection{四个行为指标怎么算}

\begin{longtable}{@{}p{3.4cm}p{6.6cm}p{3.8cm}@{}}
\toprule
\textbf{指标} & \textbf{定义} & \textbf{取值} \\
\midrule
作息集中度 $R$ & 活动时刻在 168 小时周期上的\textbf{合向量长度} & 0=完全分散，1=完全集中 \\
作息熵 & 周内分布归一化后的香农熵 & 0=集中，1=均匀 \\
间隔变异 CV & 相邻活动间隔的标准差 $\div$ 均值 & 越小越规律 \\
大小基尼 & 活动量的基尼系数 & 0=均匀，1=极不均 \\
\bottomrule
\end{longtable}

\begin{explainbox}{什么是“合向量长度”}
把每个活动时刻换算成圆上的一个角度，然后把这些角度的单位向量\textbf{加起来}，
取合向量的长度。\\[4pt]
\begin{itemize}[leftmargin=1.4em,itemsep=1pt]
\item 如果所有活动都在同一时段 → 向量全同向 → 合向量很长（$R \to 1$）；
\item 如果活动均匀散布在整周 → 向量相互抵消 → 合向量很短（$R \to 0$）。
\end{itemize}
这是一个纯粹的几何量，不需要任何参数，所以无法通过调参来“制造”结果。
\end{explainbox}

\newpage

% ============================================================

\section{把 MBTI 和节律指标接起来}
% ============================================================

第 6 节列了 8 条 MBTI 假设，第 8 节标定了效应量。
但中间还缺一环：\textbf{因变量选得对不对。}

\subsection{之前的因变量不够好}

第 6 节用的是“作息集中度 $R$”——它只有\textbf{一个数}，
把两种完全不同的模式混在一起：

\begin{center}
\begin{tabular}{p{6.2cm}p{6.2cm}}
\toprule
\textbf{模式} & \textbf{集中度 $R$} \\
\midrule
天天在固定时间学 & 高 \\
只在每周三学 & \textbf{也高} \\
\bottomrule
\end{tabular}
\end{center}

但这两类学生的干预方式\textbf{完全不同}：前者改习惯（几点学），后者改安排（哪几天有课）。

\begin{explainbox}{问题出在“一个数描述两件事”}
“集中度 $R$”衡量的是\hl{活动在周内是否聚集}——
但“天天同一时刻”和“只在某一天”都会让它变高。\\
\textbf{一个指标同时响应两个机制，就没法区分它们。}
\end{explainbox}

\subsection{新因变量：把集中度拆成三个}

用 	extbf{Rayleigh 周期图}把“集中度”拆开。这个方法很简单：\
对一个候选周期 $P$，把每个事件时刻换成	extbf{相位}：

\begin{equation}
\phi_i = \frac{t_i \bmod P}{P} \in [0, 1)
\end{equation}

如果活动确实以 $P$ 为周期聚集，$\phi_i$ 就会集中在某个值附近。
用\textbf{合向量长度}衡量集中程度：

\begin{equation}
C(P) = \frac{1}{n}\sum_i \cos(2\pi\phi_i), \qquad
S(P) = \frac{1}{n}\sum_i \sin(2\pi\phi_i)
\end{equation}

\begin{equation}
R(P) = \sqrt{C(P)^2 + S(P)^2} \in [0, 1]
\end{equation}

\begin{explainbox}{为什么用这个，不用普通功率谱}
普通功率谱（FFT / Lomb-Scargle）是给\textbf{连续信号}用的。
作息是\textbf{点过程}——有活动才有记录，一天可能只有一次。

对“每天一次”的脉冲做功率谱，会在 12h、8h、6h…… 产生大量\textbf{谐波}，
功率被摊开，反而看不清主节律。\\[4pt]
Rayleigh 周期图是点过程的标准工具，\textbf{直接问“相位是否聚集”}，
不产生谐波问题。\\[4pt]
显著性：零假设（相位均匀）下 $Z = nR^2 \sim \text{Exp}(1)$，故 $p = e^{-Z}$。
\end{explainbox}

代入两个关键周期，就得到三个指标：

\begin{longtable}{@{}p{2.6cm}p{11.2cm}@{}}
\toprule
\textbf{指标} & \textbf{含义} \\
\midrule
$R_{24}$ & 日内强度。\textbf{时间点是否固定}（每次都差不多同一时刻） \\
$R_{168}$ & 周内强度。\textbf{是否集中在周内的特定位置}（特定星期几 + 时间） \\
$B = R_{168}/R_{24}$ & \textbf{每日型还是间歇型}。$B$ 小 = 天天学；$B$ 大 = 只在某些天学 \\
\bottomrule
\end{longtable}

再加上节律基线给出的两个：

\begin{longtable}{@{}p{3.4cm}p{10.4cm}@{}}
\toprule
\textbf{指标} & \textbf{含义} \\
\midrule
偏离度中位 & 这个人\textbf{平时}偏离自己基线多少（水平） \\
偏离度标准差 & 这个人的偏离\textbf{稳不稳定}（波动） \\
\bottomrule
\end{longtable}

\subsubsection{实测示例}

在 EdNet 数据上，两个学生的对比很清楚地说明这个拆分有用：

\begin{center}
\begin{tabular}{lrrrl}
\toprule
\textbf{学生} & $R_{24}$ & $R_{168}$ & $B$ & \textbf{类型} \\
\midrule
u1 & 0.119 & 0.301 & \textbf{2.536} & \textbf{间歇型}（只在某些天） \\
u1007 & \textbf{0.942} & 0.600 & 0.637 & \textbf{每日型}（天天固定时刻） \\
\bottomrule
\end{tabular}
\end{center}

\textbf{u1 的 $R_{24}$ 只有 0.119}（时间点很不固定），但 $B$ 高达 2.54；
\textbf{u1007 的 $R_{24}$ 是 0.942}（时间点高度固定）。
在旧的“集中度 $R$”里，这两个人可能得到相近的分数——\textbf{但他们是完全不同的人。}

\subsection{核心理论点：trait 与 state 必须分开}

这是本次结合中\textbf{最重要的一个区分}。

\begin{explainbox}{trait（特质）vs state（状态）}
\begin{center}
\begin{tabular}{p{2.6cm}p{3.4cm}p{7.4cm}}
\toprule
 & \textbf{稳定性} & \textbf{影响什么} \\
\midrule
\textbf{性格（trait）} & 稳定 & 节律的\textbf{水平}：$R_{24}$、$B$ 的平均值 \\
\textbf{事件（state）} & 临时 & 节律的\textbf{波动}：偏离度的波动大小 \\
\bottomrule
\end{tabular}
\end{center}
\end{explainbox}

\begin{explainbox}{为什么必须分开}
如果把两者混在一起，就会把
\textbf{“这周考试所以作息乱了”}
错误归因成
\textbf{“这个人性格如此”}。\\[6pt]
这样的结论不仅错，而且会\textbf{给学生贴上一个不该贴的标签}——
而标签造成的自我实现预言，正是第 6 节开头警告过的事。\\[6pt]
所以假设必须分成两组：
\begin{itemize}[leftmargin=1.4em,itemsep=1pt]
\item \textbf{A 组（水平假设）}：性格 $\to$ 节律指标的\textbf{平均水平}；
\item \textbf{B 组（波动假设）}：性格 $\to$ 偏离度的\textbf{波动大小}。
\end{itemize}
\end{explainbox}

\subsection{更新后的假设表}

\begin{longtable}{@{}p{1.0cm}p{2.4cm}p{2.8cm}p{3.6cm}p{1.4cm}p{2.6cm}@{}}
\toprule
\textbf{编号} & \textbf{组} & \textbf{性格维度} & \textbf{因变量} & \textbf{方向} & \textbf{理由} \\
\midrule
H1 & A 水平 & JP（判断端） & $R_{24}$ & $+$ & 偏好计划 $\to$ 时间点固定 \\
H2 & A 水平 & JP（判断端） & 节律平衡 $B$ & $-$ & J 偏每日型，P 偏间歇型 \\
H3 & A 水平 & Si & $R_{24}$ & $+$ & 依赖既有习惯 $\to$ 高度重复 \\
H4 & A 水平 & Ne & 节律平衡 $B$ & $+$ & 发散探索 $\to$ 集中在少数时段 \\
H5 & B 波动 & JP（判断端） & 偏离度标准差 & $-$ & 计划性强 $\to$ 状态更稳 \\
H6 & B 波动 & Ne & 偏离度标准差 & $+$ & 发散 $\to$ 起伏更大 \\
H7 & \textbf{对照} & EI & 偏离度中位 & \textbf{0} & \textbf{阴性对照，理论上不应相关} \\
\bottomrule
\end{longtable}

\begin{explainbox}{H7 是“阴性对照”，为什么要有它}
H7 假设 EI（内外倾）与作息稳定性\textbf{没有}直接理论联系。\\
如果这个假设\textbf{显著了}，说明分析流程里可能有系统偏差——
比如某个共同的混杂变量，或者某个指标定义有问题。\\[4pt]
\textbf{阴性对照的作用，和前面“安慰剂检验”是同一个思路}：
用一个“本该没有效应”的检验，来检查工具本身。
\end{explainbox}

\subsection{这个结合带来了三样东西}

\begin{enumerate}[leftmargin=1.8em]
\item \textbf{因变量从“一个模糊的数”变成“两个可分别干预的维度”}——
      $R_{24}$ 对应改习惯，$B$ 对应改安排；
\item \textbf{引入了 trait/state 的区分}——避免把临时事件错误归因成性格；
\item \textbf{加了一条阴性对照}（H7）——可以检测系统性偏差。
\end{enumerate}

\subsection{对预注册的影响}

预注册文档里的因变量需要相应更新：

\begin{center}
\begin{tabular}{p{5.4cm}p{7.0cm}}
\toprule
\textbf{预注册 v1（原）} & \textbf{预注册 v2（更新）} \\
\midrule
主假设因变量：作息集中度 $R$ & 主假设因变量：\textbf{$R_{24}$}（时间点规律） \\
次要因变量：作息熵、间隔 CV 等 & 次要因变量：\textbf{$B$}、偏离度中位、偏离度标准差 \\
— & 新增阴性对照 \textbf{H7} \\
\bottomrule
\end{tabular}
\end{center}

\begin{explainbox}{必须在看数据之前改}
现在\textbf{还没有收任何 MBTI 数据}，所以现在改是\textbf{合法的}。\\
一旦收了数据再改因变量，那就变成了事后挑选——\\
\textbf{这正是预注册要防的事。}\\[4pt]
换句话说：\textbf{改设计要趁早，趁你还没看到答案。}
\end{explainbox}

\newpage

% ============================================================
% ============================================================

\section{基线：用户与环境的耦合}
% ============================================================

\subsection{为什么基线不是“固定性格”}

假设一个学生原来每天晚上 20:00 学习，表现稳定。
后来学校把晚自习调到 19:00，他连续两周状态变差。
这时他的“基线”变了吗？

如果基线是固定性格，答案是不变；
如果基线是耦合稳定点，答案是：\textbf{环境变了，基线也跟着移动。}

第四期采用后一种定义。

\subsection{形式定义：耦合动力系统}

把学生状态记作 $x_t$，学生行为记作 $u_t$，环境记作 $e_t$，慢参数记作 $\theta$：

\begin{equation}
x_{t+1}=f(x_t,u_t,e_t;\theta)+\varepsilon_t.
\end{equation}

符号逐一说明：

\begin{itemize}[leftmargin=2em]
\item $x_t$：第 $t$ 期的学生状态（注意力、情绪、认知负荷、节律）；
\item $u_t$：学生自己的行为（学习、休息、求助）；
\item $e_t$：外部环境（任务、教师、同伴、排课、家庭）；
\item $\theta$：慢参数，类似 trait（MBTI 连续分数、节律峰型、注意偏好）；
\item $f$：耦合函数；
\item $\varepsilon_t$：噪声。
\end{itemize}

基线 $x^*$ 是这个系统的稳定点：

\begin{equation}
x^*=f(x^*,u^*,e^*;\theta).
\end{equation}

所以：

\begin{equation}
\text{基线}=\text{用户节律}\times\text{任务结构}\times\text{教师安排}\times\text{同伴环境}
\end{equation}

的耦合结果，而不是“这个人单独拥有的一个数”。

\subsection{和 MBTI 的区别}

\begin{center}
\begin{longtable}{@{}p{2.4cm}p{4.6cm}p{6.8cm}@{}}
\toprule
 & \textbf{MBTI 类 trait} & \textbf{基线} \\
\midrule
数学角色 & 慢参数 $\theta$ & 耦合系统的稳定点 $x^*$ \\
来源 & 自评量表 & 行为轨迹估计 \\
稳定性 & 中等，重测信度弱 & 环境稳定时稳定，环境变则移动 \\
干预 & 很难直接改 & 可以通过改安排来移动 \\
正确用法 & 带不确定度的先验 & 动态监测的参照点 \\
\bottomrule
\end{longtable}
\end{center}

\begin{explainbox}{一句话区分}
MBTI 是“这个人倾向于怎样”；\\
基线是“这个人和当前环境耦合成了怎样”。\\
两者短期都稳定，但扰动环境后，基线会移动，trait 不会移动那么多。
\end{explainbox}

\subsection{扰动实验：怎么判断一个量是 trait 还是 baseline}

\begin{enumerate}[leftmargin=2em]
\item 改变环境：换时段、换任务、换分组；
\item 观察指标是否移动；
\item 如果移动明显，说明它是耦合/状态；
\item 如果几乎不动，才更像慢参数。
\end{enumerate}

这也是“顺应天性优于强制改习惯”的形式化：

\begin{equation}
\Delta x^*
=
\underbrace{\frac{\partial f}{\partial e}\Delta e}_{\text{改安排}}
+
\underbrace{\frac{\partial f}{\partial x}\Delta x}_{\text{改习惯}}
+
\underbrace{\frac{\partial f}{\partial \theta}\Delta \theta}_{\text{改性格，最慢}}.
\end{equation}

改安排通常成本最低、见效最快；改习惯其次；改性格最慢且最不确定。

\subsection{真实数据：E2 与 E3}

\subsubsection{E2 节律--安排匹配（EdNet）}

\begin{center}
\begin{longtable}{@{}p{4.6cm}rr@{}}
\toprule
\textbf{模型} & \textbf{系数} & \textbf{$p$} \\
\midrule
参与度 $\log(1+\text{事件数})\sim$ 匹配度 fit & $+0.164$ & $8.3\times10^{-10}$ \\
退出率 $\sim$ 匹配度 fit & $-0.00344$ & $0.0078$ \\
\bottomrule
\end{longtable}
\end{center}

对齐会话 vs 错位会话：$3.913$ vs $3.534$，配对差 $+0.379$，
$p=2.8\times10^{-15}$，$d_z=0.572$，有双条件的 $221$ 名学生。

\subsubsection{E3 习惯改变成本（OULAD）}

\begin{center}
\begin{longtable}{@{}p{6.0cm}rr@{}}
\toprule
\textbf{模型} & \textbf{系数} & \textbf{$p$} \\
\midrule
下一周参与度 $\sim$ 习惯改变 & $-1.134$ & $5.9\times10^{-247}$ \\
下一周退课 $\sim$ 习惯改变 & $+0.233$ & $1.2\times10^{-197}$ \\
控制当前周点击后的参与度 & $-0.431$ & $1.9\times10^{-40}$ \\
\bottomrule
\end{longtable}
\end{center}

高改变周 vs 低改变周：参与度差 $-0.623$（$d_z=-0.482$），退课率差 $+0.144$（$d_z=0.423$）。

\begin{explainbox}{边界}
E2/E3 都是观察性数据，不是随机强制改习惯的实验。\\
它们支持“耦合基线”的方向，但不能单独证明因果。
\end{explainbox}

\newpage

% ============================================================

\section{节律--安排匹配与习惯改变成本}
% ============================================================

\subsection{核心原则}

\begin{explainbox}{安排 $>$ 强制改习惯}
先改安排（低成本、低风险），再改习惯（高成本、高不确定）。\\
只有当预期收益超过改变成本时，才要求改习惯。
\end{explainbox}

\subsection{E2：节律--安排匹配}

\subsubsection{定义}

每个学生有一个日内峰型 $\text{peak}_i$。一次会话的开始小时记作 $h_s$。
定义匹配度：

\begin{equation}
\mathrm{fit}_s=
\cos\left(\frac{2\pi(h_s-\text{peak}_i)}{24}\right).
\end{equation}

符号说明：

\begin{itemize}[leftmargin=2em]
\item $\text{peak}_i$：学生 $i$ 的日内节律峰型（0--24 小时）；
\item $h_s$：会话 $s$ 的开始小时；
\item $\mathrm{fit}_s=+1$：正好在峰型；
\item $\mathrm{fit}_s=-1$：相差 12 小时。
\end{itemize}

\subsubsection{真实结果（EdNet）}

\begin{center}
\begin{longtable}{@{}lrr@{}}
\toprule
\textbf{模型} & \textbf{系数} & \textbf{$p$} \\
\midrule
参与度 $\log(1+\text{事件数})\sim\mathrm{fit}$ & $+0.164$ & $8.3\times10^{-10}$ \\
退出率 $\sim\mathrm{fit}$ & $-0.00344$ & $0.0078$ \\
\bottomrule
\end{longtable}
\end{center}

对齐会话 vs 错位会话：$3.913$ vs $3.534$，配对差 $+0.379$，
$p=2.8\times10^{-15}$，$d_z=0.572$。

\subsubsection{功效}

\begin{center}
\begin{longtable}{@{}lrrr@{}}
\toprule
\textbf{真实 $\beta$} & \textbf{19 人} & \textbf{50 人} & \textbf{100 人} \\
\midrule
$0.05$ & $0.219$ & $0.483$ & $0.774$ \\
$0.10$ & $0.657$ & $0.970$ & $1.000$ \\
$0.20$ & $0.997$ & $1.000$ & $1.000$ \\
\bottomrule
\end{longtable}
\end{center}

\subsection{E3：习惯改变成本}

\subsubsection{定义}

把前 4 周按 $date\bmod 7$ 的点击分布记作 baseline\_profile。
第 $w$ 周的分布记作 current\_profile。定义：

\begin{equation}
\text{habit\_change}_w
=
\frac12\sum_{k=0}^{6}
\left|
\mathrm{current\_profile}_{w,k}
-
\mathrm{baseline\_profile}_{k}
\right|.
\end{equation}

符号说明：

\begin{itemize}[leftmargin=2em]
\item $k=0,\ldots,6$：一周内的 7 个相位；
\item $0$：完全保持原有节奏；
\item $1$：完全换了一周节奏；
\item 这个量只用点击分布，不用内容，所以可跨课程比较。
\end{itemize}

\subsubsection{真实结果（OULAD）}

\begin{center}
\begin{longtable}{@{}lrr@{}}
\toprule
\textbf{模型} & \textbf{系数} & \textbf{$p$} \\
\midrule
下一周参与度 $\sim$ habit\_change & $-1.134$ & $5.9\times10^{-247}$ \\
下一周退课 $\sim$ habit\_change & $+0.233$ & $1.2\times10^{-197}$ \\
控制当前周点击后的参与度 & $-0.431$ & $1.9\times10^{-40}$ \\
\bottomrule
\end{longtable}
\end{center}

高改变周 vs 低改变周：

\begin{itemize}[leftmargin=2em]
\item 参与度：$2.060$ vs $2.683$，差 $-0.623$，$d_z=-0.482$；
\item 退课率：$0.286$ vs $0.142$，差 $+0.144$，$d_z=0.423$。
\end{itemize}

\subsubsection{边界}

这是观察性关联，不是随机强制改习惯的实验。\\
控制当前周点击后仍显著，说明不只是“低参与度的镜像”，但仍可能有第三变量。

\subsection{课堂怎么用}

\begin{enumerate}[leftmargin=2em]
\item 先测每个学生的节律峰型（可用 1--2 周的自报 + 行为日志）；
\item 排课时尽量把高难度任务放在峰型附近；
\item 如果要改变学生的习惯，先问：是安排不能改，还是我们不愿意改安排？
\item 如果必须改习惯，给小步、给反馈、给缓冲期，不要一次改到底。
\end{enumerate}

\newpage
% ============================================================

\section{黄金比、基线闭环与自相似}
% ============================================================

\subsection{三个概念要分开}

\subsubsection{基线闭环}

系统受扰动后回到基线，形成一个回路：

\begin{equation}
x_{t+T}\approx x_t.
\end{equation}

符号说明：

\begin{itemize}[leftmargin=2em]
\item $x_t$：第 $t$ 期状态；
\item $T$：回路周期；
\item $x_{t+T}\approx x_t$：经过一个周期后回到接近原来的状态。
\end{itemize}

\subsubsection{自相似}

同一结构在不同尺度重复出现：

\begin{equation}
f(\lambda x)=\lambda^\alpha f(x).
\end{equation}

符号说明：

\begin{itemize}[leftmargin=2em]
\item $\lambda$：尺度因子；
\item $\alpha$：标度指数；
\item 这是\textbf{幂律/尺度不变}，不是黄金比。
\end{itemize}

\subsubsection{黄金比}

\begin{equation}
\varphi=\frac{1+\sqrt5}{2}\approx1.618.
\end{equation}

\subsection{关键结论：自相似 $+$ 闭环不自动推出黄金比}

自相似只推出“没有特征尺度”（power law），不推出某一个特定常数。\\
黄金比只在\textbf{特定动力系统}里出现，例如某些 KAM 环面的黄金旋转数、Penrose 铺砌、部分最优觅食模型；它不是自相似系统的普遍常数。

\subsection{你项目里现在的数据}

\begin{center}
\begin{longtable}{@{}p{6.0cm}p{4.5cm}p{3.0cm}@{}}
\toprule
\textbf{量} & \textbf{当前值} & \textbf{是不是 $\varphi$} \\
\midrule
会话大小幂律指数 & $1.342$ & 不是 \\
周节律 $168/24$ & $7$ & 不是 \\
$24$ 小时 $/$ $90$ 分钟 & $\approx16$ & 不是 \\
\bottomrule
\end{longtable}
\end{center}

所以：\textbf{黄金比目前只能当隐喻，不能当结论。}

\subsection{如果一定要检验它：预注册比例检验}

\begin{enumerate}[leftmargin=2em]
\item 先定义两个嵌套尺度，例如：
  \begin{itemize}
  \item 日循环 / 周循环；
  \item 周循环 / 月循环；
  \item 个人 / 小组 / 班级。
  \end{itemize}
\item 计算比例 $r=T_{k+1}/T_k$；
\item 用 bootstrap 给 $r$ 的置信区间；
\item 比较三个模型：
  \begin{itemize}
  \item $H_\varphi: r=\varphi$；
  \item $H_1: r=1$；
  \item $H_2: r$ 自由估计。
  \end{itemize}
\item 用 Bayes factor 或交叉验证比较。
\end{enumerate}

\begin{explainbox}{必须预注册}
尺度定义必须先写死。\\
否则在大量比例里总能找到一个接近 $\varphi$ 的，属于事后挑选。
\end{explainbox}

\subsection{更稳妥的“和谐”判据}

与其追求 $\varphi$，不如用两个可检验的性质：

\begin{enumerate}[leftmargin=2em]
\item \textbf{回归性}：扰动后能回到基线，不发散；
\item \textbf{尺度不变性}：不同尺度上的形状相同，指数稳定。
\end{enumerate}

也就是：

\begin{equation}
\text{闭环}+\text{自相似}
\Rightarrow
\text{尺度不变的稳定回归}.
\end{equation}

这比“黄金比”更可证伪，也更贴合花/果/森林结构。

\newpage

% ============================================================


\part{验证、防错与预注册}

\section{关键工具一：安慰剂检验}
% ============================================================

\subsection{什么是安慰剂检验}

\term{安慰剂检验}（placebo test）来自医学：给病人吃没有药效的“假药”，
如果病人也变好了，说明之前的效果不是药带来的。

搬到数据分析上：

\begin{quote}
\textbf{把数据里的“结构”破坏掉，重新做一次分析。\\
如果结果不变，说明原来的结果不是那个结构带来的。}
\end{quote}

\subsection{三路对照}

我对 34 名学生做了三件事：

\begin{longtable}{@{}p{0.8cm}p{3.6cm}p{9.4cm}@{}}
\toprule
\textbf{序号} & \textbf{对照} & \textbf{做法} \\
\midrule
1 & \textbf{真实时间排布} & 原始数据，不动 \\
2 & \textbf{间隔重排} & 保持所有“间隔”的数值不变，只把它们打乱顺序 \\
3 & \textbf{均匀泊松} & 同样的事件数、同样的时间跨度，但时间完全随机 \\
\bottomrule
\end{longtable}

\begin{explainbox}{什么叫“均匀泊松”}
均匀泊松就是“完全随机、没有记忆”的事件流：
任何时刻发生的概率都相同，而且不受之前事件影响。\\
\textbf{它是“完全没有结构”的基准线。}
\end{explainbox}

\subsection{结果}

\begin{center}
\begin{tabular}{p{5.0cm}p{3.4cm}p{4.8cm}}
\toprule
\textbf{条件} & \textbf{中位分支比} & \textbf{说明} \\
\midrule
EdNet 真实 & \textbf{0.9318} & 原始数据 \\
间隔重排 & \textbf{0.9318} & 打乱间隔顺序 \\
均匀泊松 & \textbf{0.0866} & 完全随机 \\
\bottomrule
\end{tabular}
\end{center}

\begin{explainbox}{这三行怎么读（关键）}
\textbf{第一行 vs 第三行}：差 \hl{0.845} $\to$ 学生的作答\textbf{确实不是随机的}，有强结构。\\
\textbf{第一行 vs 第二行}：差 \hl{0.0001}（几乎没有）$\to$
\textbf{打乱间隔顺序后，估计值一点没变}。\\[4pt]
\textbf{第二行才是要点}：如果自激是真的，打乱先后顺序就该破坏它、让估计值掉下来。
但它没掉。这说明\warn{这个模型根本没有在用“先后顺序”这个信息}。
\end{explainbox}

\subsection{发现的第一个问题}

\begin{quote}
\textbf{模型拟合的是“间隔的分布”（有时隔几秒、有时隔几天），
而不是“事件之间的激发关系”。}
\end{quote}

在统计上这叫\term{识别性问题}：指数核 Hawkes 与“间隔过度离散的更新过程”
在这类数据上\textbf{无法区分}。

\newpage

% ============================================================

\section{关键工具二：时间重标检验}
% ============================================================

上一章只发现了问题，没有给出正确结论。这一章用更严格的工具。

\subsection{原理}

\term{时间重标定理}（Meyer 1971，Brown 等 2002）：

\begin{quote}
如果点过程的条件强度是 $\lambda(t)$，那么把每个事件时刻做如下变换：
\end{quote}

\begin{equation}
\Lambda(t_i) = \int_0^{t_i} \lambda(s)\, ds
\label{p2:eq:compensator}
\end{equation}

\begin{explainbox}{逐符号解释}
\begin{longtable}{@{}p{2.4cm}p{11.4cm}@{}}
$\Lambda(t_i)$ & 第 $i$ 个事件的“累计强度”（读作“大写的 Lambda”） \\
$\int_0^{t_i}$ & 从 0 到 $t_i$ 做积分——可以理解为“把这段时间里的强度全加起来” \\
$\lambda(s)$ & $s$ 时刻的瞬时强度（就是式 \eqref{p2:eq:intensity}） \\
\end{longtable}
\end{explainbox}

\begin{explainbox}{这个定理说了什么}
\textbf{变换之后，事件流应该变成一个“标准随机流”。}\\
具体说：重标后的间隔
\begin{equation*}
\tau_i = \Lambda(t_i) - \Lambda(t_{i-1})
\end{equation*}
应当服从\textbf{标准指数分布 Exp(1)}，而且彼此独立。\\[4pt]
\textbf{所以这就是一个检验}：
\begin{itemize}[leftmargin=1.4em,itemsep=1pt]
\item 如果 $\tau$ 真的像 Exp(1) $\to$ 模型对
\item 如果不像 $\to$ 模型设定有误
\end{itemize}
\end{explainbox}

\subsection{什么叫“像 Exp(1)”}

Exp(1) 是“平均每秒发生 1 次、且完全随机”的间隔分布。它的特点：

\begin{itemize}[leftmargin=1.6em]
\item 大量短间隔、少量长间隔，长尾很长；
\item 变异系数（标准差除以均值）等于 1。
\end{itemize}

检验方法用 \term{KS 检验}：

\begin{explainbox}{KS 检验零基础版}
把实际数据和理论分布画成两条累积曲线，看两条线最大的垂直距离。\\[4pt]
\begin{itemize}[leftmargin=1.4em,itemsep=1pt]
\item 距离小 $\to$ 很像 $\to$ 通过
\item 距离大 $\to$ 不像 $\to$ 不通过
\end{itemize}
这个距离叫 \term{KS 统计量}，越小越好。\\
同时给出 $p$ 值：\textbf{$p > 0.05$ 算通过}，$p < 0.05$ 算不通过。
\end{explainbox}

\subsection{引入一个真正的竞争模型}

要判断“自激”是不是多余的机制，需要一个更简单的对照：

\begin{longtable}{@{}p{3.6cm}p{4.4cm}p{5.8cm}@{}}
\toprule
\textbf{模型} & \textbf{强度依赖什么} & \textbf{参数个数} \\
\midrule
Poisson & 什么都不依赖 & 1 \\
\textbf{Weibull 更新过程} & \textbf{只依赖上一个事件} & 2 \\
Hawkes & \textbf{依赖全部历史} & 3 \\
\bottomrule
\end{longtable}

\begin{explainbox}{为什么这个对照很关键}
如果 Weibull 更新过程（\textbf{只依赖上一个事件}）与 Hawkes（\textbf{依赖全部历史}）
拟合得一样好，那“自激 / 长记忆”这个机制就是多余的。\\
这就是奥卡姆剃刀：\textbf{能简单解释，就不要用复杂解释}。
\end{explainbox}

\subsection{结果}

\subsubsection{模型比较（45 名学生）}

\begin{center}
\begin{tabular}{p{4.6cm}rr}
\toprule
\textbf{模型} & \textbf{中位对数似然} & \textbf{中位 AIC} \\
\midrule
Poisson & $-385.4$ & $+772.9$ \\
Weibull 更新 & $+336.4$ & $-668.9$ \\
\textbf{Hawkes} & \textbf{+481.0} & \textbf{$-956.0$} \\
\bottomrule
\end{tabular}
\end{center}

\begin{explainbox}{AIC 是什么}
AIC 是“模型的综合评价”，同时考虑两件事：
\begin{itemize}[leftmargin=1.4em,itemsep=1pt]
\item 拟合得好不好（对数似然越大越好）；
\item 参数多不多（参数越多越容易过拟合，要惩罚）。
\end{itemize}
公式：$\text{AIC} = -2\times(\text{对数似然}) + 2\times(\text{参数个数})$。\\
\textbf{AIC 越小越好。}
\end{explainbox}

\textbf{读数}：Hawkes 优于 Weibull 更新：\hl{43 / 45 人}。
中位 AIC 差距约 \textbf{287}——不是微小差别。

而且 Weibull 的形状参数中位是 $k = 0.299$（$k<1$ 表示间隔高度聚簇）。

\subsubsection{时间重标检验}

\begin{center}
\begin{tabular}{p{4.6cm}rrr}
\toprule
\textbf{模型} & \textbf{中位 KS} & \textbf{通过率} & \textbf{Lag-1 自相关} \\
\midrule
\textbf{Hawkes} & \textbf{0.110} & \textbf{38 \%} & $-0.104$ \\
Weibull 更新 & 0.258 & 0 \% & $-0.006$ \\
Poisson & 0.884 & 0 \% & $-0.009$ \\
\bottomrule
\end{tabular}
\end{center}

\begin{explainbox}{三条读数}
\textbf{1. Poisson 彻底失败}（KS 0.884，通过率 0\%）：数据绝对不是常速率。\\
\textbf{2. Weibull 更新也失败}（KS 0.258，通过率 0\%）：只依赖上一个事件不够。\\
\textbf{3. Hawkes 明显更好，但也没通过}（KS 0.110，仅 38\% 通过）：
方向对了，\warn{设定还不对}。
\end{explainbox}

\subsection{修正后的结论}

到这里，故事变成了三层：

\begin{longtable}{@{}p{1.0cm}p{4.4cm}p{3.4cm}p{6.6cm}@{}}
\toprule
\textbf{层次} & \textbf{问题} & \textbf{答案} & \textbf{证据} \\
\midrule
{(1)}& 数据有结构吗？ & \hl{有，很强} & 0.93 vs 泊松 0.09 \\
{(2)}& 只依赖上一个事件吗？ & \hl{不是} & Hawkes 优于 Weibull：43/45 \\
{(3)}& Hawkes 够吗？ & \warn{不够} & 时间重标只有 38\% 通过 \\
\bottomrule
\end{longtable}

\textbf{所以上一章的表述要修正}：
“打乱间隔顺序结果不变”现在的解释是——
\textbf{单指数核在这个参数区间对顺序不敏感}，这是核函数的局限，
\textbf{不等于“没有长记忆”}。

\newpage
% ============================================================

\section{两次改进尝试（都失败）}
% ============================================================

既然 Hawkes 的设定不对，那就要改。我试了两个方向。

\subsection{尝试一：加一个时间尺度}

\subsubsection{动机}

单指数核只有一个时间尺度。而学生的行为至少有两个：

\begin{center}
\begin{tabular}{p{2.2cm}p{3.6cm}p{7.6cm}}
\toprule
\textbf{尺度} & \textbf{时间范围} & \textbf{机制} \\
\midrule
快 & 秒 ~ 分钟 & 一道题到下一道题（会话内） \\
慢 & 小时 ~ 天 & 一次学习到下一次学习（会话间） \\
\bottomrule
\end{tabular}
\end{center}

所以把核函数改成两个指数相加：

\begin{equation}
\phi(\Delta t) = \alpha_1 e^{-\beta_1 \Delta t} + \alpha_2 e^{-\beta_2 \Delta t}
\label{p2:eq:multiscale}
\end{equation}

\subsubsection{第一次尝试失败的原因}

第一次实现时，我直接在 $(\mu, \alpha_1, \beta_1, \alpha_2, \beta_2)$
上做粗网格 + Nelder-Mead。结果两尺度反而比单尺度更差——
这不可能，因为\textbf{单尺度是两尺度在 $\alpha_2=0$ 时的特例}，
似然不可能更低。

\begin{explainbox}{这说明什么}
\textbf{似然更低 = 优化没收敛}，不是模型不好。\\
问题有三个：
\begin{enumerate}[leftmargin=1.4em,itemsep=1pt]
\item $\beta$ 跨越 0.01~10 三个数量级，线性搜索空间尺度失衡；
\item 平稳性约束靠“拒绝”实现，优化器容易撞墙；
\item 没有 warm start，从随机点出发常常跑到边界。
\end{enumerate}
\end{explainbox}

\subsubsection{修法}

\begin{longtable}{@{}p{0.8cm}p{4.0cm}p{9.0cm}@{}}
\toprule
\textbf{序号} & \textbf{修法} & \textbf{说明} \\
\midrule
1 & 参数变换 & 在 log 空间优化 $(\log\mu, \log\beta_1, \log\beta_2)$；
分支比用两个 logit 保证 $n_1+n_2<0.98$，\textbf{约束内建}，不再撞墙 \\
2 & warm start & 先拟合单尺度，再用它的解分裂出快/慢两个尺度作为初值 \\
3 & 多起点 & 从多个初值出发，取似然最高者 \\
\bottomrule
\end{longtable}

\subsubsection{合成数据验证（先证明优化器能用）}

用已知参数生成数据，再拟合回来：

\begin{center}
\begin{tabular}{p{4.2cm}rr}
\toprule
\textbf{两个尺度的分离程度} & \textbf{2 尺度中位误差} & \textbf{误差 $<0.10$ 比例} \\
\midrule
分离明显（$\beta$ 差 60 倍） & \hl{0.019} & \hl{4/5} \\
分离中等（差 13 倍） & 0.287 & 1/5 \\
几乎单尺度（差 2 倍） & 0.341 & 0/5 \\
\bottomrule
\end{tabular}
\end{center}

\begin{explainbox}{读数}
\textbf{两个尺度分离明显时，新优化器能准确恢复参数（误差 0.019）。}\\
尺度接近时不可识别——这是数学上的必然（两条相近的指数和一条指数无法区分）。\\[4pt]
同时可见：当真值是两尺度时，\textbf{单尺度核会系统性低估}（误差 $-0.27$ 到 $-0.40$）。
\end{explainbox}

\subsubsection{EdNet 上的结果}

\begin{center}
\begin{tabular}{p{4.6cm}rr}
\toprule
\textbf{指标} & \textbf{1 尺度} & \textbf{2 尺度} \\
\midrule
中位 AIC & $-956.0$ & \hl{$-961.5$} \\
AIC 更好的人数 & --- & \hl{19 / 25} \\
中位 KS & 0.1051 & 0.1032 \\
\textbf{KS 通过率} & \textbf{36 \%} & \warn{\textbf{36 \%}} \\
中位对数似然提升 & --- & $+7.2$ \\
\bottomrule
\end{tabular}
\end{center}

拟合出的尺度结构：慢尺度 $\beta=0.82$、分支比 0.069；
快尺度 $\beta=24.77$、分支比 \textbf{0.886}；分离倍数 \textbf{28 倍}。

\begin{explainbox}{结论}
两尺度在 AIC 上确实更好（19/25），
但\textbf{时间重标通过率一点没改善}。\\[4pt]
\textbf{所以“缺尺度”不是根因。}
\end{explainbox}

\subsection{尝试二：加昼夜周期}

\subsubsection{先验证昼夜结构是否存在}

对 52 名学生的 49,576 条作答事件，按本地小时统计：

\begin{center}
\begin{tabular}{p{6.0cm}rr}
\toprule
\textbf{时段} & \textbf{事件数} & \textbf{相对水平} \\
\midrule
\textbf{21 时（峰值）} & 4,770 & 最高 \\
17 时 & 3,946 & 较高 \\
\textbf{5 时（谷值）} & 139 & 最低 \\
\bottomrule
\end{tabular}
\end{center}

\begin{equation}
\chi^2 = 23{,}155 \quad (\text{自由度} = 23)
\end{equation}

\begin{explainbox}{$\chi^2$ 是什么意思}
$\chi^2$（读作“卡方”）衡量“实际分布与均匀分布差多远”。
\begin{itemize}[leftmargin=1.4em,itemsep=1pt]
\item 如果事件在 24 小时里均匀分布，$\chi^2$ 应该约等于自由度 23；
\item 实测 \textbf{23,155} —— 大了 1000 倍。
\end{itemize}
\textbf{这是压倒性的周期结构。}峰谷比 \textbf{34.3 倍}。
\end{explainbox}

\subsubsection{昼夜调制 Hawkes}

把常数基线换成随时间变化的基线：

\begin{equation}
\lambda(t) = \underbrace{\mu_0\, m(\text{hour}(t))}_{\text{昼夜基线}} + \underbrace{R(t)}_{\text{自激}}
\label{p2:eq:circadian}
\end{equation}

\begin{explainbox}{逐符号解释}
\begin{longtable}{@{}p{2.4cm}p{11.4cm}@{}}
$\mu_0$ & 基线的整体水平（一个常数） \\
$m(\cdot)$ & 昼夜调制函数，表示“这个小时相对活跃多少” \\
$\text{hour}(t)$ & $t$ 时刻是一天中的第几小时 \\
$m$ 的约束 & 在一个周期上归一化：$\frac{1}{24}\int_0^{24} m(h)\,dh = 1$ \\
$R(t)$ & 自激分量，与式 \eqref{p2:eq:intensity} 的求和项相同 \\
\end{longtable}
$m$ 的做法：把一天切成 96 格（每 15 分钟一格），
统计每格里有多少事件，做完循环平滑后归一化。
\end{explainbox}

\subsubsection{结果}

\begin{center}
\begin{tabular}{p{4.6cm}rr}
\toprule
\textbf{指标} & \textbf{常数基线} & \textbf{昼夜基线} \\
\midrule
中位 AIC & $-978.8$ & \hl{$-1000.1$} \\
\textbf{AIC 更好的人数} & --- & \hl{\textbf{30 / 30}} \\
中位分支比 $n$ & 0.936 & 0.930 \\
中位 KS & 0.1099 & 0.1089 \\
\textbf{KS 通过率} & \textbf{36.7 \%} & \warn{\textbf{36.7 \%}} \\
\bottomrule
\end{tabular}
\end{center}

\begin{explainbox}{结论}
昼夜基线让 AIC \textbf{全面改善（30/30，平均改善 21.3）}
$\to$ 周期结构确实存在且重要。\\
但：\textbf{分支比几乎没动}（0.936 $\to$ 0.930），
\textbf{KS 通过率完全没变}（36.7\% $\to$ 36.7\%）。\\[4pt]
\textbf{所以“缺昼夜”也不是根因。}
\end{explainbox}

\newpage
% ============================================================

\section{方法论总结}
% ============================================================

\subsection{四次尝试的完整对照}

\begin{longtable}{@{}p{1.0cm}p{4.6cm}p{3.0cm}p{6.4cm}@{}}
\toprule
\textbf{序号} & \textbf{尝试} & \textbf{KS 通过率} & \textbf{结论} \\
\midrule
1 & 单尺度 Hawkes & 36 \% & 重标失败 \\
2 & + 第二个时间尺度 & 36 \% & AIC 略好，重标仍失败 \\
3 & + 昼夜周期 & 36.7 \% & AIC 明显好，重标仍失败 \\
4 & \textbf{换成会话为单位} & \hl{\textbf{100 \%}} & \hl{\textbf{通过}} \\
\bottomrule
\end{longtable}

\begin{explainbox}{一句话总结这四次尝试}
\textbf{前三步都是在“把模型改复杂”，第四步是“把单位换对”。}\\[4pt]
当模型越改越复杂却仍然通不过检验时，
问题通常在\hl{建模对象的定义}，而不在模型的容量。
\end{explainbox}

\subsection{这一期用到的两个检验工具}

\begin{longtable}{@{}p{4.2cm}p{9.6cm}@{}}
\toprule
\textbf{工具} & \textbf{作用} \\
\midrule
\textbf{对照实验}（三路安慰剂） & 用“打乱 / 随机”的版本做对照，看结果会不会变 \\
\textbf{时间重标检验} & 把拟合好的模型代回数据，看重标间隔是否像 Exp(1) \\
\bottomrule
\end{longtable}

\begin{explainbox}{为什么这两个工具重要}
在一个样本量小、结构不可观测的场景里，
\hl{“发现的规律”比“没发现规律”更可疑}。\\
这两个工具的价值就是：\textbf{在写结论之前，先证明它不是假的。}
\end{explainbox}

\subsection{对平台场景的落地建议}

\begin{longtable}{@{}p{0.8cm}p{12.8cm}@{}}
\toprule
\textbf{序号} & \textbf{建议} \\
\midrule
1 & \textbf{先剔除昼夜与周周期}再做动态建模（昼夜让 AIC 改善 30/30，
不剔除就会把周期当自激） \\
2 & \textbf{不要假设自激}（三次尝试都无法通过时间重标） \\
3 & \textbf{改用“会话到达”建模}---把基本单位从“单次互动”换成
“一次访问 / 一次曝光组” \\
4 & \textbf{保留时间重标检验作为验收标准}---它是唯一能识别“模型设定是否有误”的工具 \\
5 & \textbf{记住观测模型 $\neq$ 干预模型}---无论 Hawkes 拟合得多好，
“互动引发互动”与“平台推流”在时间序列上仍不可区分，
这需要\textbf{干预实验}，不是更复杂的模型 \\
\bottomrule
\end{longtable}

\newpage
% ============================================================

\section{错误一：AIC 漏计参数（推翻上一版结论）}
% ============================================================

这是这一期最重要的内容。

\subsection{上一期说了什么}

上一期我做了这么个实验：给会话层的 Hawkes 加一个“周课表”基线
（把一周切成 168 个自由格），然后看分支比怎么变。

结果是：分支比从 0.497 掉到 \textbf{0.000}，而且
\textbf{“17/17 人的 AIC 都变得更好了”}。

当时的结论是：

\begin{quote}
\textbf{“那 0.54 的自激完全来自课表，不是习惯。”}
\end{quote}

\subsection{错在哪里}

这一段代码是错的：

\begin{center}
\begin{tabular}{p{7.6cm}p{5.2cm}}
\toprule
\textbf{我写的} & \textbf{应该是} \\
\midrule
\texttt{aic = -2*loglik + 2*3} & \texttt{aic = -2*loglik + 2*(3+168)} \\
\bottomrule
\end{tabular}
\end{center}

\begin{explainbox}{AIC 的公式与含义}
\begin{equation}
\text{AIC} = -2 \times (\text{对数似然}) + 2 \times (\text{参数个数})
\end{equation}
\begin{itemize}[leftmargin=1.4em,itemsep=1pt]
\item 第一项：拟合得好不好（似然越大越好）；
\item 第二项：\textbf{参数越多，惩罚越大}（防止过拟合）；
\item \textbf{AIC 越小越好。}
\end{itemize}
周课表基线用了 168 个自由格，所以它\textbf{应该}被罚 168 个参数。
我只罚了 3 个——\textbf{等于让这个模型免费获得 168 个自由度}。
\end{explainbox}

\subsection{修正后结论完全反转}

\begin{center}
\begin{tabular}{p{4.6cm}rr}
\toprule
\textbf{指标} & \textbf{旧口径（错）} & \textbf{正确口径} \\
\midrule
无周期 AIC & 193.2 & 193.2 \\
自由格 AIC & 173.7 & \textbf{509.7} \\
“加周期更好”人数 & \textbf{17 / 17} & \textbf{0 / 17} \\
\bottomrule
\end{tabular}
\end{center}

\textbf{自由格模型不是“更好”，而是差了 316 个 AIC 单位}——属于严重过拟合。

它之所以能把分支比压到 0，是因为它有 \textbf{171 个参数去拟合十几个数据点}。
\textbf{不是发现了结构，是把噪声也拟合进去了。}

\subsection{换上参数化模型后的正确答案}

用傅里叶基代替 168 个自由格（第 4 节会讲为什么必须改形式）：

\begin{center}
\begin{tabular}{p{4.4cm}rrr}
\toprule
\textbf{模型} & \textbf{分支比} & \textbf{AIC} & \textbf{参数个数} \\
\midrule
无周期 & 0.445 & 185.3 & 3 \\
日内 K=1 & 0.377 & 182.5 & 5 \\
\textbf{日内 K=2} & \textbf{0.292} & \textbf{145.8} & \textbf{7} \\
日内 K2 + 周内 K1 & 0.238 & 147.4 & 9 \\
自由格 168 & 0.000 & \textbf{509.7} & 171 \\
\bottomrule
\end{tabular}
\end{center}

\textbf{三条修正后的结论}：

\begin{enumerate}[leftmargin=1.8em]
\item AIC 最优是“日内 K=2”，\textbf{只用 7 个参数}；
\item 作息解释了约 1/3 的自激（$0.445 \to 0.292$），\textbf{但没解释完}；
\item 自由格的“归零”是过拟合假象。
\end{enumerate}

\begin{explainbox}{这一节最值得记的一句话}
\textbf{“17/17 一致”这种完美结果，最该怀疑。}\\
真实的效应很少全票通过。\\
当所有样本都指向同一方向、而且方向特别干净时，
先去看代码，而不是先庆祝。
\end{explainbox}

\newpage

% ============================================================

\section{错误二：一个不可识别的自由度（振幅）}
% ============================================================

刚修完 AIC，又踩了第二个坑。

\subsection{症状}

做“跨用户共享节律”时（让所有人共享一套作息形状），
第一次拟合系数直接炸到 $10^{295}$，似然变成 $+\infty$。

\subsection{原因：正弦积分为零}

原来的调制写法是：

\begin{equation}
m(h) = 1 + \sum \big[a\cos(\omega h) + b\sin(\omega h)\big]
\end{equation}

问题出在 Hawkes 对数似然里的\textbf{补偿项}（积分项）：

\begin{equation}
\int_0^T \lambda(t)\,dt \supset \mu_0 \int_0^T m(h(t))\,dt
\end{equation}

\begin{explainbox}{关键事实}
\textbf{正弦和余弦在一个完整周期上的积分恰好为 0}：
\begin{equation}
\int_0^{P} \cos\!\Big(\frac{2\pi k h}{P}\Big) dh = 0, \qquad
\int_0^{P} \sin\!\Big(\frac{2\pi k h}{P}\Big) dh = 0
\end{equation}
所以调制项的积分只跟“周期数”有关：
\begin{equation}
\int_0^T m \approx 168 \times (\text{周期数})
\end{equation}
\textbf{这个积分完全不依赖振幅 $a, b$ 的大小。}
\end{explainbox}

后果：把振幅推到无穷大，
- 事件落在波峰时收益很大（似然上升）；
- 补偿项却完全不变（没有惩罚）。

\textbf{于是似然可以无限增长——振幅不可识别。}

\subsection{修法：对数线性形式}

\begin{equation}
s(h) = \sum \big[a\cos(\omega h) + b\sin(\omega h)\big]
\end{equation}

\begin{equation}
m(h) = \frac{\exp(s(h))}{Z},
\qquad
Z = \frac{1}{168}\int_0^{168}\exp(s)\,dh
\end{equation}

\begin{explainbox}{为什么这样就好了}
\begin{enumerate}[leftmargin=1.4em,itemsep=2pt]
\item \textbf{正性}：指数函数永远为正，$m(h)>0$ 自动满足，不用再裁剪；
\item \textbf{归一化}：除以 $Z$ 后平均值恒为 1，所以 $\mu_0$ 就是“平均强度”，含义清晰；
\item \textbf{振幅可识别}：$\exp$ 是\textbf{非线性}的——振幅越大，波形越\textbf{尖}。
      尖波形在高峰时段收益更大，但在低谷时段代价也更大，
      两者平衡后有唯一最优。
\end{enumerate}
\end{explainbox}

程序自检（代入任意系数）：

\begin{center}
\begin{tabular}{lrl}
\toprule
\textbf{检查项} & \textbf{实测} & \textbf{应该} \\
\midrule
周期积分 $M(168)$ & 168.0000 & $=168$ \\
$m$ 的均值 & 1.0004 & $\approx 1$ \\
$m$ 的最小值 & 0.586 & $>0$ \\
\bottomrule
\end{tabular}
\end{center}

\subsection{修正的副作用（必须诚实报告）}

\begin{center}
\begin{tabular}{p{5.4cm}rr}
\toprule
\textbf{版本} & \textbf{分支比中位} & \textbf{问题} \\
\midrule
修正前（自由振幅） & 0.29 & 振幅不可识别，\warn{估计偏乐观} \\
\textbf{修正后（对数线性）} & \textbf{0.64} & 可信 \\
\bottomrule
\end{tabular}
\end{center}

\textbf{之前的 0.29 有一部分是靠“不受约束的调制”吸收掉的}，不是真实的作息解释力。

\begin{explainbox}{这一节最值得记的一句话}
\textbf{模型里每一个不受约束的自由度，最后都会变成虚高的结论。}\\
加参数之前，先问：\textbf{这个参数的最优值是被数据决定的，还是可以随便跑？}\\
如果它可以随便跑，那它一定会跑到让你高兴的地方去。
\end{explainbox}

\newpage

% ============================================================

\section{功效分析：事前的防错工具}
% ============================================================

上面那张假设表列好了，下一个问题是：\textbf{要多少人才测得出来？}

\subsection{什么是功效}

\begin{explainbox}{功效（power）的定义}
假设某个效应\textbf{真实存在}，
你做的实验能把它检测成“显著”的概率，就叫\textbf{功效}。\\[4pt]
惯例标准：\textbf{功效 $\ge 0.80$} 才算够用。\\[4pt]
功效不足的后果是：\textbf{真效应被当成没效应}（假阴性）。
\end{explainbox}

\subsection{样本量 vs 可检测效应}

在显著水平 $\alpha=0.05$、功效 $0.80$ 下：

\begin{center}
\begin{tabular}{rr}
\toprule
\textbf{学生数 $N$} & \textbf{可检测到的最小 $\lvert r\rvert$} \\
\midrule
10 & 0.785 \\
15 & 0.669 \\
\textbf{19（两个班的规模）} & \textbf{0.605} \\
30 & 0.492 \\
47 & 0.400 \\
\textbf{92} & \textbf{0.289} \\
200 & 0.197 \\
\bottomrule
\end{tabular}
\end{center}

\subsection{反过来：想检测某个效应，需要多少人}

\begin{center}
\begin{tabular}{rr}
\toprule
\textbf{目标 $\lvert r\rvert$} & \textbf{所需学生数} \\
\midrule
0.3（弱） & \textbf{85} \\
0.4（中） & \textbf{47} \\
0.5（中强） & 30 \\
0.6（强） & 20 \\
0.7（很强） & 14 \\
\bottomrule
\end{tabular}
\end{center}

\subsection{用 19 人做实验，功效是多少}

\begin{center}
\begin{tabular}{rr}
\toprule
\textbf{真实 $\lvert r\rvert$} & \textbf{$N=19$ 时的功效} \\
\midrule
0.3 & \textbf{24.5\%} \\
0.4 & \textbf{39.1\%} \\
0.5 & 63.8\% \\
0.6 & 83.6\% \\
0.7 & 96.0\% \\
\bottomrule
\end{tabular}
\end{center}

\begin{explainbox}{怎么读这张表}
如果真实效应是 0.4（中等），用 19 人做实验，
\textbf{有 61\% 的概率测不出显著结果}——即使这个效应真实存在。\\[4pt]
所以：\textbf{测不出不等于不存在}。
\end{explainbox}

\subsection{模拟验证：19 人的精度有多差}

假设真实效应 $r=0.6$，模拟 19 名学生、重复 2000 次：

\begin{center}
\begin{tabular}{lr}
\toprule
\textbf{指标} & \textbf{数值} \\
\midrule
估计 $r$ 的均值 & 0.593 \\
估计 $r$ 的 \textbf{95\% 区间} & \textbf{[0.246, 0.835]} \\
检验出显著的比例（功效） & \textbf{81.2\%} \\
\bottomrule
\end{tabular}
\end{center}

\textbf{注意那个区间}：真实效应是 0.6，但估计值可能从 0.25 到 0.84。
\textbf{这就是 19 人的精度。}

\newpage

% ============================================================

\section{用公开数据标定“真实效应有多大”}
% ============================================================

功效分析告诉我们“多少样本能测多大的效应”。
但还有一个问题没回答：\textbf{真实世界里的效应，通常有多大？}

如果效应普遍很小，那需要的人数会比上面算的更多。

\subsection{为什么不自己估，而是找数据}

MBTI 数据项目里没有。但\textbf{同一个问题结构}的公开数据有：

\begin{center}
\begin{tabular}{p{3.4cm}p{10.4cm}}
\toprule
\textbf{需要的东西} & \textbf{OULAD 有什么} \\
\midrule
个体属性 & 性别、年龄段、教育背景、地区贫困指数、残障、已修学分、历史重考次数 \\
带日期的行为 & \texttt{date} + \texttt{sum\_click}，\textbf{逐日}记录 \\
样本量 & \textbf{26,074 名学生 / 10,655,280 行} \\
\bottomrule
\end{tabular}
\end{center}

\textbf{属性虽然不是 MBTI，但“属性 $\to$ 行为规律性”这个问题结构完全一样。}

\subsection{做法}

\begin{enumerate}[leftmargin=1.8em]
\item 对每个学生，把他的逐日活跃记录转成一组规律性指标
      （活跃密度、间隔中位、间隔变异、最长中断）；
\item 计算每个属性与每个指标的效应量：
  \begin{itemize}
  \item 二元属性（性别、残障）→ \textbf{Cohen's $d$}；
  \item 有序属性（年龄、教育、贫困）→ \textbf{Spearman 秩相关 $r$}；
  \end{itemize}
\item 看这些效应量的分布。
\end{enumerate}

\subsection{结果：真实效应非常小}

\begin{center}
\begin{tabular}{p{4.4cm}p{3.6cm}rr}
\toprule
\textbf{属性} & \textbf{指标} & \textbf{效应量} & \textbf{p 值} \\
\midrule
性别 F vs M & 活跃占比 & \textbf{$-0.289$} & 3.2e-118 \\
性别 F vs M & 间隔中位 & 0.221 & 2.1e-67 \\
性别 F vs M & 最长间隔 & 0.177 & 3.7e-45 \\
年龄段 & 活跃占比 & 0.101 & 1.2e-60 \\
残障 Y vs N & 间隔中位 & 0.071 & 2.2e-02 \\
教育背景 & 间隔中位 & $-0.067$ & 3.2e-27 \\
贫困指数 & 活跃占比 & 0.046 & 1.2e-13 \\
\bottomrule
\end{tabular}
\end{center}

\begin{center}
\begin{tabular}{lr}
\toprule
\textbf{汇总} & \textbf{数值} \\
\midrule
中位绝对效应 & \textbf{0.053} \\
最大绝对效应 & \textbf{0.289} \\
超过 0.3 的条数 & \textbf{0 / 20} \\
超过 0.2 的条数 & 2 / 20 \\
\bottomrule
\end{tabular}
\end{center}

\begin{explainbox}{这里有两个值得注意的地方}
\textbf{第一：效应量普遍很小。}\\
中位数只有 0.053——按行为科学惯例，$\lvert d\rvert < 0.2$ 算“可忽略”。
二十条里只有两条超过 0.2，最强的（性别）也只有 0.289。\\[6pt]
\textbf{第二：p 值全都极度显著。}\\
2.6 万人的样本下，$\lvert r\rvert = 0.05$ 也能得到 $p = 10^{-118}$。\\
\textbf{“显著”和“有用”是两回事}——大样本下什么都能显著。
\end{explainbox}

\subsection{翻译成“要多少样本”}

\begin{center}
\begin{tabular}{lr}
\toprule
\textbf{要检测的效应} & \textbf{所需学生数} \\
\midrule
0.289（本数据里的最大值） & \textbf{92 人} \\
0.2（“小效应”下限） & 194 人 \\
0.05（本数据里的中位值） & \textbf{2792 人} \\
\bottomrule
\end{tabular}
\end{center}

\textbf{用 19 名学生做这个研究：}

\begin{center}
\begin{tabular}{lr}
\toprule
\textbf{真实效应} & \textbf{$N=19$ 的功效} \\
\midrule
0.289（最乐观假设） & \textbf{约 23\%} \\
0.2 & 约 12\% \\
0.05 & 约 5\%（等于瞎猜） \\
\bottomrule
\end{tabular}
\end{center}

\begin{explainbox}{这意味着什么}
\textbf{即使 MBTI 的一个维度，对作息的预测力和“性别”一样强}
——这是相当乐观的假设——用 19 名学生也只有 23\% 的概率测得出来。\\[6pt]
换句话说：\textbf{如果测不出，你无法区分这两种可能}：
\begin{itemize}[leftmargin=1.4em,itemsep=1pt]
\item 效应不存在；
\item 效应存在，但样本不够。
\end{itemize}
\end{explainbox}

\subsection{一个反直觉的推论}

本数据里最强的预测变量是\textbf{性别}（0.289），而这个变量：
\begin{itemize}[leftmargin=1.6em]
\item 与“性格”的关系是间接的；
\item 与行为的关系更多来自\textbf{生活安排}（比如是否有照顾责任）。
\end{itemize}

\begin{explainbox}{一个由数据提示的猜想}
\textbf{影响学习作息规律性的，可能主要不是“性格”，而是“生活约束”。}\\[4pt]
如果这条成立，那么在做学生作息干预时，
\textbf{了解他的时间约束（几点下课、几点睡、要不要带弟妹）
可能比了解他的 MBTI 更有用。}\\[4pt]
（这只是一个猜想，需要专门设计研究才能确认。）
\end{explainbox}

\newpage
% ============================================================
% ============================================================

\section{预注册：把规则写在看到数据之前}
% ============================================================

功效分析解决了“要多少样本”。
但还有一个更隐蔽的问题：\textbf{收完数据之后，你总可以“找到”某个显著结果。}

\subsection{为什么需要预注册}

假设你有 8 条假设、4 个指标、2 个子组。
组合起来有几十种可能的检验，随便试总能碰到 $p < 0.05$ 的。

\begin{explainbox}{这在本项目里等价于前面踩过的坑}
\begin{center}
\begin{tabular}{p{6.4cm}p{6.0cm}}
\toprule
\textbf{前面几轮} & \textbf{本质} \\
\midrule
自由格把分支比压到 0 & 自由度太多 → 拟合噪声 \\
AIC 漏计 168 个参数 & 假装没有那些自由度 \\
\textbf{事后挑选显著结果} & \textbf{自由度太多 → 挑出噪声} \\
\bottomrule
\end{tabular}
\end{center}
\textbf{三个是同一件事。}
\end{explainbox}

\subsection{预注册锁死了什么}

本项目为此生成了一份完整的预注册文档，主要内容：

\begin{longtable}{@{}p{3.2cm}p{10.6cm}@{}}
\toprule
\textbf{项目} & \textbf{锁定的内容} \\
\midrule
\textbf{主假设} & \textbf{只有一个}：H1（JP 判断端 $\to$ 作息集中度，正相关） \\
次要假设 & 5 条，Bonferroni 校正到 $\alpha = 0.05/5 = 0.01$ \\
探索性分析 & 明确列出，\textbf{不得作为结论宣称} \\
指标公式 & 全部写死，不事后调整 \\
样本量 & 90 人（依据 OULAD 标定的最大效应 0.289） \\
判定规则 & 支持 / 方向相反 / 未检出，三种措辞都事先规定 \\
失效条件 & 有效样本 $<70$ 时不报相关性结论 \\
不做的事 & 不事后换指标、不事后增删样本、不亚组分析后单独宣称 \\
伦理 & 监护人同意、去标识、数据不出校、\textbf{不反馈个体性格分数} \\
\bottomrule
\end{longtable}

\subsection{为什么主假设只有一个}

\begin{explainbox}{多重比较问题}
如果你检验 20 个独立的假设，每个都用 $\alpha = 0.05$，
那么\textbf{至少有一个假阳性的概率}是：
\begin{equation}
1 - (1-0.05)^{20} \approx 64\%
\end{equation}
也就是说：\textbf{就算这些假设全是假的，你也有 64\% 的概率“发现”至少一个显著结果。}\\[4pt]
所以主假设只留一个，其余全部降级并做校正。
\end{explainbox}

\subsection{为什么“未检出”必须配上功效}

\begin{explainbox}{三种措辞，事先规定}
\begin{center}
\begin{tabular}{p{4.2cm}p{8.6cm}}
\toprule
\textbf{结果} & \textbf{必须怎么说} \\
\midrule
主假设 $p<0.05$ 且方向对 & “支持”，报告 $\rho$ 与置信区间 \\
主假设 $p<0.05$ 但方向反 & 报告反向结果，\textbf{不强行解释} \\
主假设 $p \ge 0.05$ & \textbf{“未检出显著相关”}，并注明功效不足 \\
任何情况 & \textbf{不得说“证明了 MBTI”或“推翻了 MBTI”} \\
\bottomrule
\end{tabular}
\end{center}
\end{explainbox}

\subsection{分析脚本的自检（含一次真实的翻车）}

预注册配了一份分析脚本。为了确认它真的能工作，我做了一个自检：
造一份“JP 越高活动越集中在固定星期几”的合成数据，看脚本能不能恢复。

\textbf{第一次自检失败}：真实效应是正的，脚本却给出 $-0.161$。

查出来是\textbf{我的数据生成器写错了}：
把相位信息转成日期时丢了，而且每周只生成 1 天活动
（结果全被“活动次数 $\ge 20$”的排除标准筛掉了）。

修完再跑：

\begin{center}
\begin{tabular}{p{4.0cm}rrl}
\toprule
\textbf{假设} & \textbf{$\rho$} & \textbf{p 值} & \textbf{判定} \\
\midrule
H1（真实存在） & \textbf{0.627} & 5.0e-10 & \textbf{支持} \\
H2（由 H1 推出） & $-0.699$ & 5.4e-13 & \textbf{支持} \\
H3--H6（合成数据里不存在） & $-0.13 \sim 0.13$ & $0.24 \sim 0.58$ & \textbf{未检出} \\
\bottomrule
\end{tabular}
\end{center}

\textbf{管线能找到真效应，且没有假阳性。}

\begin{explainbox}{这次翻车说明了什么}
如果我没做自检，直接拿脚本去跑真实数据，
\textbf{所有结果都会是错的}，而且看不出错在哪里。\\[4pt]
\textbf{自检的价值和预注册一样：它让你在“还不知道答案”的时候，
先确认工具是可用的。}
\end{explainbox}

\newpage

% ============================================================

\section{方法论主线：从做、验、防到接}
% ============================================================

\subsection{每期在解决什么问题}

\begin{longtable}{@{}p{1.6cm}p{4.6cm}p{4.4cm}p{3.2cm}@{}}
\toprule
\textbf{期} & \textbf{主题} & \textbf{核心工具} & \textbf{典型错误} \\
\midrule
第一期 & 怎么建模、怎么评价 & AUC / RMSE / LogLoss & 参数选错了优化目标 \\
第二期 & 怎么检验模型 & 安慰剂 / 时间重标 & 单位选错了 \\
\textbf{第三期} & \textbf{怎么防止自欺} & \textbf{约束 / 功效 / 预注册} & \textbf{自由度被滥用} \\
\bottomrule
\end{longtable}

\subsection{三期里“翻车”的完整清单}

这一期最值得看的其实是这张表——\textbf{这些都是真实发生过的}：

\begin{longtable}{@{}p{1.0cm}p{5.4cm}p{4.4cm}p{3.0cm}@{}}
\toprule
\textbf{序号} & \textbf{错误} & \textbf{怎么发现的} & \textbf{期} \\
\midrule
1 & $\eta$ 默认值偏大 15 倍 & 敏感性扫描 & 第一 \\
2 & 事件层的“自激”是会话结构 & 三路安慰剂 & 第二 \\
3 & 自由格“归零”是过拟合 & AIC 参数个数对照 & 第三 \\
4 & 傅里叶振幅不可识别 & 共享拟合发散 & 第三 \\
5 & 自检数据生成器写错 & 自检本身 & 第三 \\
6 & 差点用 19 人下结论 & 功效分析 & 第三 \\
\bottomrule
\end{longtable}

\begin{explainbox}{这张表的规律}
\textbf{每一次错误，都是被“另一个角度”发现的}：
\begin{itemize}[leftmargin=1.4em,itemsep=1pt]
\item 换成合成数据（能验证估计器）；
\item 换成竞争模型（能暴露过拟合）；
\item 换一个更受约束的模型（能暴露自由度滥用）；
\item 事前算功效（能避免样本不足）。
\end{itemize}
\textbf{没有一次是靠“再想一想”发现的。}
\end{explainbox}

\subsection{第三期的三条方法论}

\begin{enumerate}[leftmargin=2em]
\item \textbf{先问“这个自由度能不能被数据决定”，再决定要不要加它。}\\
      不可识别的自由度一定会被滥用（振幅那个例子）。

\item \textbf{“完美一致”是最该怀疑的信号。}\\
      “17/17 全票通过”后来被证明是算错参数个数。
      真实效应很少全票通过。

\item \textbf{防错的工具有“事前”和“事后”两类，两类都需要。}\\
      事后：安慰剂、竞争模型、更受约束的模型。\\
      事前：功效分析、预注册、管线自检。
\end{enumerate}

\newpage

% ============================================================

\section{验证层的翻车与边界}
% ============================================================

\subsection{四次翻车}

\begin{enumerate}[leftmargin=2em]
\item \textbf{E5 的 Fisher-$z$ 方差检验设定错误}。把相关估计的采样方差当成已知常数，导致功效低于名义水平；改用 ANOVA 后，5 位教师、$sd_\tau=0.3$ 时功效为 0.903。
\item \textbf{E4 的“抱团有害”假设被数据推翻了一半}。小圈子与触达正相关，桥接与点赞正相关；结果取决于因变量，不能一刀切。
\item \textbf{“单词测试零区分度”不能推广}。第 5/7 篇单词测试平均难度为 0.233 和 0.257，并非零区分度。
\item \textbf{MBTI 不能当天生硬约束}。重测信度弱，荣格八维也没有强心理测量学支持；只能用连续分数和节律峰型。
\end{enumerate}

\subsection{必须写清的边界}

\begin{center}
\begin{longtable}{@{}p{4.2cm}p{9.0cm}@{}}
\toprule
\textbf{边界} & \textbf{说明} \\
\midrule
观察性 vs 因果 & E2/E3/E4 都是观察性数据，不能单独证明因果 \\
公开数据推广 & EdNet 是韩国数据，课堂推广前需要本地节律基线 \\
19 人功效 & 19 人只能测中等以上效应，小效应测不出来 \\
犹豫的混淆 & 犹豫受疲劳、网络、题长影响，必须控制题目难度和题长 \\
外部脚手架伦理 & 不能伪造信心、不能用假反馈操纵学生 \\
网络指标 & 共帖网络不是真实社会关系；reads/up 不是学习结果 \\
黄金比 & 当前没有证据；只能当隐喻，不能当结论 \\
\bottomrule
\end{longtable}
\end{center}

\subsection{这一期最值得记住的三句话}

\begin{enumerate}[leftmargin=2em]
\item \textbf{犹豫是元认知的行为证据，但必须先验证 AUC，不能因为学生慢就默认他不会。}
\item \textbf{外部只能“引出和训练”元认知，不能硬植入；区分锚定和内部化的关键是迁移。}
\item \textbf{基线是用户与环境的耦合稳定点，不是 MBTI 式 trait；改安排通常比强制改习惯更优。}
\end{enumerate}

\newpage

% ============================================================


\part{元认知、教学干预与教师观察者}

\section{元认知：把“知道自己会不会”变成可测的指标}
% ============================================================

第 6 节讲 MBTI（性格），第 9 节讲节律。这一节讲另一个个体差异维度，
但它比前两个都更靠近“学习”本身：\textbf{元认知}。

\subsection{什么是元认知，和“元元”什么关系}

\begin{explainbox}{三个“元”的关系}
\begin{center}
\begin{tabular}{p{3.2cm}p{10.6cm}}
\toprule
\textbf{概念} & \textbf{含义} \\
\midrule
认知 & 学习、记忆、解题 \\
\textbf{元认知} & \textbf{对自己认知的认知}——“我知道我不会这个词” \\
元元 & 对测量的测量（第 11 节）——“这个推断有多可信” \\
\bottomrule
\end{tabular}
\end{center}
\textbf{元认知是“能力”层面的，元元是“测量”层面的。}\\
但两者结构同源：\textbf{都是对某一层再加一层}。
\end{explainbox}

\subsection{核心文献：元认知必须拆成两个独立的量}

本节最重要的文献是：

\begin{quote}
\textbf{Fleming \& Lau (2014)}, \textit{How to measure metacognition},
Frontiers in Human Neuroscience.\\
DOI \texttt{10.3389/fnhum.2014.00443}，\textbf{被引 1388}
\end{quote}

它的核心贡献是：\textbf{元认知不是一个数，至少是两个。}

\begin{center}
\begin{tabular}{p{3.4cm}p{4.2cm}p{6.2cm}}
\toprule
\textbf{量} & \textbf{含义} & \textbf{高值的表现} \\
\midrule
\textbf{元认知敏感性} & 能不能区分自己“会”与“不会” & 标记的重点正好是他做错的 \\
\textbf{元认知偏差} & 总体偏自信还是偏保守 & 决定他标得多还是少 \\
\bottomrule
\end{tabular}
\end{center}

\begin{explainbox}{为什么必须分开}
一个学生可能：\\
· \textbf{偏自信}（很少标记），但\textbf{很敏感}（标的那几个确实不会）；\\
· 或者\textbf{偏保守}（标很多），但\textbf{没准头}（标的和错的对不上）。\\[4pt]
把两者合成一个数，这两种完全不同的学生会被混为一类——
\textbf{和“作息集中度把每日型和间歇型混在一起”是同一类错误}。
\end{explainbox}

\subsection{项目里已经有元认知数据}

这是本节最实用的发现。讲义上的两类数据天然构成元认知测量：

\begin{center}
\begin{tabular}{p{3.6cm}p{3.4cm}p{6.8cm}}
\toprule
\textbf{讲义上的东西} & \textbf{认知含义} & \textbf{元认知含义} \\
\midrule
单词/句子前的\textbf{根号} & “这是重点” & \textbf{“我不会这个”} \\
56--65 题的\textbf{实际对错} & 客观表现 & 一阶结果 \\
\bottomrule
\end{tabular}
\end{center}

两者一列联，就是标准的元认知 $2\times2$ 表：

\begin{equation}
\begin{array}{c|cc}
 & \text{实际答错} & \text{实际答对} \\ \hline
\text{标记为重点} & \text{命中 } H & \text{虚报 } FA \\
\text{未标记} & \text{漏报} & \text{正确拒绝}
\end{array}
\end{equation}

\begin{explainbox}{这个表里的四个格子}
\begin{itemize}[leftmargin=1.4em,itemsep=1pt]
\item \textbf{命中}：标记了，而且确实错了 —— 元认知对了
\item \textbf{虚报}：标记了，但实际会 —— \textbf{过度谨慎}
\item \textbf{漏报}：没标记，但实际错了 —— \textbf{过度自信}
\item \textbf{正确拒绝}：没标记，也确实会 —— 元认知对了
\end{itemize}
\end{explainbox}

\subsection{指标定义（逐符号）}

\subsubsection{元认知敏感性 AUC}

\begin{equation}
\text{AUC} = P\Big(\text{错题被标记}\Big) \ \text{vs} \ P\Big(\text{对题被标记}\Big)
\end{equation}

形式化写为：

\begin{equation}
\text{AUC} = \frac{\#(\text{pos} > \text{neg}) + 0.5 \cdot \#(\text{pos} = \text{neg})}{n_{\text{pos}} \cdot n_{\text{neg}}}
\end{equation}

\begin{explainbox}{每一个符号}
\begin{longtable}{@{}p{3.0cm}p{10.8cm}@{}}
$\text{AUC}$ & 曲线下面积。$0.5$ = 标记与对错无关；$1.0$ = 标记完美对应错题 \\
$\#(\cdot)$ & 计数符号：满足条件的配对个数 \\
$\text{pos}$ & 正类样本，即\textbf{实际答错}的那些题的标记值 \\
$\text{neg}$ & 负类样本，即\textbf{实际答对}的那些题的标记值 \\
$>$ 与 $=$ & 逐对比较：每个错题的标记值与每个对题比 \\
$0.5 \cdot \#(\text{pos} = \text{neg})$ & 平局算一半 \\
$n_{\text{pos}}, n_{\text{neg}}$ & 正类与负类的样本数 \\
\end{longtable}
\textbf{读法}：随机取一道错题和一道对题，
错题被标记的概率更高的概率，就是 AUC。
\end{explainbox}

\subsubsection{元认知偏差}

\begin{equation}
\text{bias} = \bar{M} - \bar{E}
\end{equation}

\begin{explainbox}{每一个符号}
\begin{longtable}{@{}p{2.6cm}p{11.2cm}@{}}
$\text{bias}$ & 元认知偏差 \\
$\bar{M}$ & 标记率 = 被标记的题数 $/$ 总题数 \\
$\bar{E}$ & 错误率 = 答错的题数 $/$ 总题数 \\
\end{longtable}
$> 0$：标记得比实际错的多 $\to$ \textbf{过度谨慎}（宁可多标）\\
$< 0$：标记得比实际错的少 $\to$ \textbf{过度自信}（漏标）
\end{explainbox}

\newpage
\subsection{与 L2 元元层的连接}

\subsubsection{第 14 课的测谎模块做的是同一件事}

\begin{center}
\begin{tabular}{p{3.0cm}p{5.2cm}p{5.6cm}}
\toprule
 & \textbf{第 14 课（测谎模块）} & \textbf{本节} \\
\midrule
主观信号 & 自标颜色 & 标记的重点 \\
客观信号 & HRV 生理 & 实际对错 \\
一致性 & 自报--生理一致性 & 元认知校准度 \\
产出 & 测谎指数 $D$ & 敏感性 AUC \\
后果 & 自报权重打折 & 自评权重打折 \\
\bottomrule
\end{tabular}
\end{center}

\textbf{同一个元元逻辑，用在不同信号对上。}

\subsubsection{折减公式}

\begin{equation}
\text{MC}_{\text{adj}} = \text{MC}_{\text{meta}} \times d(\text{AUC})
\end{equation}

\begin{equation}
d(\text{AUC}) = \text{clip}\left(0.3 + 0.7 \cdot \frac{\text{AUC} - 0.5}{0.9 - 0.5},\; 0.3,\; 1.0\right)
\end{equation}

\begin{explainbox}{每一个符号}
\begin{longtable}{@{}p{3.4cm}p{10.4cm}@{}}
$\text{MC}_{\text{meta}}$ & L2 元元层算出的原始可信度 \\
$\text{AUC}$ & 该学生的元认知敏感性 \\
$d(\cdot)$ & 折扣函数，把 AUC 映射到 $[0.3, 1.0]$ \\
$\text{clip}(\cdot, a, b)$ & 截断函数：把值限制在 $[a,b]$ 内 \\
$\text{MC}_{\text{adj}}$ & 折减后的可信度 \\
\end{longtable}
\textbf{下限 0.3 是刻意的}：不把权重降到 0——
“标记不准”不等于“自报完全没信息”，只是可信度低。
\end{explainbox}

\subsection{公开数据集实验：在 1{,}936 名学生上验证}

\subsubsection{先踩了一个陷阱：求助行为不是元认知}

ASSISTments 里有一列 \texttt{hint\_count}（要了几次提示），看起来像元认知。
但列联表显示：

\begin{center}
\begin{tabular}{lrr}
\toprule
 & \textbf{答错} & \textbf{答对} \\
\midrule
没求助 & 52{,}301 & 180{,}892 \\
\textbf{求助过} & \textbf{40{,}033} & \textbf{9} \\
\bottomrule
\end{tabular}
\end{center}

\textbf{求助过的 40{,}042 条记录里只有 9 条答对（0.022\%）。}

\begin{explainbox}{为什么这个信号不能用}
在 ASSISTments 里，提示是做题过程中\textbf{随时可点}的。
学生点提示，是因为\textbf{已经卡住了}——那时结果基本已经注定。\\[4pt]
元认知信号必须发生在\textbf{结果未知之前}。
求助发生在结果“已经知道”之后，\textbf{时点不对}。\\[4pt]
如果直接拿它做指标，算出的 AUC 会接近 \textbf{1.0}（看似完美），
但那是\textbf{同义反复}——\textbf{好看的数字不等于有意义的数字。}
\end{explainbox}

\subsubsection{可用的信号：响应时间}

\texttt{ms\_first\_response} 是\textbf{首次作答之前}的耗时，发生在结果未知之前。
在元认知文献里，\textbf{响应时间是标准的不确定性指标}：犹豫 = 自己觉得不确定。

\begin{equation}
\text{AUC} = P\big(\text{答错题的响应时间} > \text{答对题的响应时间}\big)
\end{equation}

数据：273{,}235 条作答，1{,}936 名学生（每人 $\ge 20$ 条）。

\subsubsection{结果}

\begin{center}
\begin{tabular}{lr}
\toprule
\textbf{统计量} & \textbf{值} \\
\midrule
\textbf{元认知敏感性中位数} & \textbf{0.5845} \\
均值 & 0.5675 \\
标准差 & 0.1381 \\
\midrule
$\text{AUC} > 0.5$ 的学生比例 & \textbf{73.7\%} \\
$\text{AUC} > 0.6$ 的学生比例 & 44.2\% \\
\bottomrule
\end{tabular}
\end{center}

\begin{explainbox}{三条读数}
\textbf{(1) 中位 AUC $=0.585$，显著高于随机}——
响应时间确实携带元认知信息：\textbf{犹豫的题更容易错}。\\[4pt]
\textbf{(2) 但 26.3\% 的学生 AUC $\le 0.5$}——
这部分学生的“犹豫”不指向错误，响应时间不携带元认知信息。\\[4pt]
\textbf{(3) AUC $>0.6$ 的只有 44.2\%}——
\textbf{不到一半的学生能通过“犹豫”有效预测自己的错误。}
\end{explainbox}

\subsubsection{最重要的发现：与学业表现强相关}

\begin{equation}
r(\text{元认知敏感性},\ \text{总体准确性}) = +0.513
\end{equation}

\begin{center}
\begin{tabular}{p{6.0cm}rp{4.0cm}}
\toprule
\textbf{变量} & \textbf{效应量} & \textbf{来源} \\
\midrule
性别 & 0.289 & OULAD（26{,}074 人） \\
年龄段 & 0.101 & 同上 \\
教育背景 & 0.067 & 同上 \\
\midrule
\textbf{元认知敏感性} & \textbf{0.513} & \textbf{ASSISTments（1{,}936 人）} \\
\bottomrule
\end{tabular}
\end{center}

\begin{explainbox}{这张表修正了前面第 8 节的推断}
第 8 节用 OULAD 标定后发现“人口学属性的效应只有 0.05--0.29”，
由此推断“可能不是性格，而是生活约束”。\\[4pt]
现在看到元认知敏感性的相关是 \textbf{0.513}——比所有人口学属性都强得多。\\[4pt]
\textbf{这修正了那个推断}：不是“性格没用”，而是\textbf{之前测错了东西}。
人口学属性（性别、年龄）本来就是粗粒度标签，预测力自然弱；
\textbf{元认知是个体内部的心理过程，预测力强得多。}
\end{explainbox}

\subsection{三个实验的执行结果（2026-10-07 更新）}

这一节把上一版列出的“三个可做的实验”全部推进到\textbf{可复核}的状态。
每个数字都注明来源：\textbf{公开数据}、\textbf{真实课堂数据}，还是\textbf{合成自检}。
三者不能混着说。

\begin{center}
\begin{longtable}{@{}p{1.8cm}p{3.0cm}p{4.1cm}p{5.4cm}@{}}
\toprule
\textbf{实验} & \textbf{当前状态} & \textbf{数据} & \textbf{最重要的结果} \\
\midrule
A\ JOL & 公开数据已完成；课堂数据等待录入 & OSF yaguz：84+80+53+41 人 & JOL 与实际相关约 $0.46$--$0.51$；难易效应 $p<10^{-21}$；全局 JOL 不优于“只说班平均” \\
B\ 根号 & 讲义提取与教师判断已完成；学生自标等待录入 & 21 份讲义 PDF；9 篇 4 标记 vs 6 未标 & AUC $=0.5625$，置换 $p=0.795$；5/7 篇单词全部标记，无对照可识别 \\
C\ 难度基线 & 已完成（修正 v1 两处错误） & 4 套测试、30 个项目、6--7 人/题 & Wilson 宽度中位 $0.591$；10 题能力 SE $\approx 0.699$ logit \\
\bottomrule
\end{longtable}
\end{center}

\subsubsection{实验 A：JOL——学生能预测自己考多少吗？}

\textbf{形式定义（不是隐喻）}：对第 $s$ 个学生、第 $t$ 次测试，令

\begin{equation}
n_{st} = \text{本次测试总题数},\qquad
y_{st} = \text{实际答对题数},\qquad
\hat y_{st} = \text{测试前预测答对题数}.
\end{equation}

逐符号说明：

\begin{itemize}[leftmargin=2em]
\item $n_{st}$：第 $s$ 个学生在第 $t$ 次测试里一共有多少题；本项目中通常是 $10$。
\item $y_{st}$：批改后实际答对几题；范围是 $0\le y_{st}\le n_{st}$。
\item $\hat y_{st}$：测试\textbf{开始前}学生自己估计能对几题；范围同样 $0\le \hat y_{st}\le n_{st}$。
\item 实际正确率 $p_{st}=y_{st}/n_{st}$，预测正确率 $\hat p_{st}=\hat y_{st}/n_{st}$。
\end{itemize}

由此得到三个核心指标：

\begin{equation}
\text{校准偏差}\quad b_{st}=\hat p_{st}-p_{st},
\end{equation}

\begin{equation}
\text{校准误差}\quad a_{st}=|\hat p_{st}-p_{st}|,
\end{equation}

\begin{equation}
\text{测试级 Brier 分数}\quad B_{st}=\hat p_{st}^{2}(1-p_{st})+(1-\hat p_{st})^{2}p_{st}.
\end{equation}

符号含义：

\begin{itemize}[leftmargin=2em]
\item $b_{st}>0$：预测得比实际好，\textbf{过度自信}；$b_{st}<0$：预测得比实际差，\textbf{过度保守}。
\item $a_{st}$：只看偏差大小，不看方向；$0$ 最准，$1$ 最差。
\item $B_{st}$：Brier 是“概率预测”的标准评分；越小越好。它同时惩罚“方向错”和“过度自信”。
\item 基线：把 $\hat p_{st}$ 换成全班平均正确率 $\bar p$，得到“只说班平均”的基线。若 JOL 的 $B$ 或 $a$ 不优于基线，说明“学生会预测”没有超过“只说一个班平均分”。
\end{itemize}

\textbf{公开数据结果}：使用 OSF yaguz（Tekin \& Roediger 的延迟 JOL 实验）中的 84 名（Exp1）和 80 名（Pilot）做了 JOL 的被试。数据、脚本和复现路径在附录。

\begin{center}
\begin{longtable}{@{}lrrrrl@{}}
\toprule
\textbf{数据} & \textbf{人数} & \textbf{JOL--实际相关} & \textbf{总偏差} & \textbf{校准误差} & \textbf{是否优于班平均} \\
\midrule
Exp1 & 84 & $0.460$--$0.499$ & $+0.052$ & $0.120$ & 否（MAE 与 Brier 都不优） \\
Pilot & 80 & $0.324$--$0.475$ & $-0.178$ & $0.203$ & 否（Pilot 明显更差） \\
\bottomrule
\end{longtable}
\end{center}

\textbf{难易效应}（公开数据）：容易的项目上，学生偏保守；难的项目上，学生偏过度自信。

\begin{center}
\begin{longtable}{@{}lrrr@{}}
\toprule
\textbf{数据} & \textbf{易项偏差} & \textbf{难项偏差} & \textbf{配对检验} \\
\midrule
Exp1 & $-0.084$ & $+0.189$ & $p=5.13\times10^{-22}$ \\
Pilot & $-0.184$ & $-0.173$ & $p=0.696$（不显著） \\
\bottomrule
\end{longtable}
\end{center}

\textbf{条件实验}（公开数据）：Exp2a 中“做 JOL”和“提取练习”都显著高于“重学”，Exp2b 的效应小一些。

\begin{center}
\begin{longtable}{@{}lrrrr@{}}
\toprule
\textbf{条件} & \textbf{Exp2a 正确率} & \textbf{vs 重学 $p$} & \textbf{Exp2b 正确率} & \textbf{vs 重学 $p$} \\
\midrule
重学（Restudy） & $0.398$ & --- & $0.426$ & --- \\
只做 JOL & $0.496$ & $2.8\times10^{-5}$ & $0.452$ & $0.496$ \\
隐蔽 JOL & $0.524$ & $3.9\times10^{-5}$ & $0.506$ & $0.0177$ \\
公开 JOL & $0.536$ & $2.9\times10^{-6}$ & $0.535$ & $0.0034$ \\
公开提取 & $0.572$ & $1.6\times10^{-8}$ & $0.522$ & $0.0046$ \\
\bottomrule
\end{longtable}
\end{center}

\textbf{课堂设计功效}：用公开数据的效应量（偏差 SD $\approx0.15$）模拟“19 人 $\times$ 10 次 $\times$ 每次 10 题”。

\begin{center}
\begin{longtable}{@{}p{6.2cm}rr@{}}
\toprule
\textbf{要检测的效应} & \textbf{功效} & \textbf{判断} \\
\midrule
平均偏差 $0.05$ & $0.268$ & 不够 \\
平均偏差 $0.10$ & $0.790$ & 勉强（MDES $\approx0.10$） \\
个体差异相关 $r=0.5$ & $0.609$ & 不够 \\
难易效应斜率 $-0.55$ & $1.000$ & 足够 \\
校准改善 $\delta=0.01$/次 & $0.872$ & 足够 \\
个体差异 SD $=0.15$ & $1.000$ & 足够 \\
\bottomrule
\end{longtable}
\end{center}

\textbf{下一步（实验 A 的课堂部分）}：打印 \texttt{jol\_collection\_cards.html}，每名学生一张；
测试前先写“估计能对几题”，批改后由老师填“实际对了几题”；
录入 \texttt{jol\_observations.csv}。真实课堂数据一旦录入，重跑
\texttt{run\_jol\_analysis.py} 就会自动生成 \texttt{classroom\_analysis.json}。
\textbf{不要让学生自己填实际对题数}，否则会把“预测”和“事后合理化”混在一起。

\subsubsection{实验 B：根号——教师的重点判断准不准？学生自己知道吗？}

\textbf{先说结论}：21 份讲义 PDF 里的 $\surd$（U+221A）是\textbf{印在 PDF 里的}，
同一份讲义所有学生看到的都一样，所以它是\textbf{教师/教材的重点判断}，
不是学生元认知。上一版“必须先确认根号性质”的问题，现在可以回答：
\textbf{就这批 PDF 而言，是老师印的。}

\textbf{讲义提取结果}（真实文件，不是模拟）：

\begin{center}
\begin{longtable}{@{}lrrr@{}}
\toprule
\textbf{讲义} & \textbf{$\surd$ 总行数} & \textbf{$\surd$ 单词} & \textbf{$\surd$ 句子} \\
\midrule
五讲 & 17 & 14 & 3 \\
七讲 & 28 & 23 & 4 \\
八讲 & 38 & 34 & 4 \\
九讲 & 30 & 25 & 5 \\
十讲 & 52 & 48 & 4 \\
其余各讲 & 0 & 0 & 0 \\
\bottomrule
\end{longtable}
\end{center}

\textbf{教师判断效度（9 篇 56--65 题，真实课堂数据）}：
把 10 个空映射到它们所在的句子，4 个句子有 $\surd$（59、62、63、64），
6 个没有（56、57、58、60、61、65）。

\begin{equation}
\text{AUC}_{\text{teacher}}=\Pr(\text{标记句错误率}>\text{未标句错误率})+\tfrac12\Pr(\text{相等}).
\end{equation}

结果：标记句平均错误率 $0.625$，未标句 $0.583$，差 $+0.042$；
AUC $=0.5625$，置换检验 $p=0.795$，学生层面 bootstrap 95\% CI $=[0.333,0.812]$。
\textbf{没有证据表明“教师标为重点的句子”更难。} 但这只有 4 个标记句和 6 个未标句，
项目数太少，只能算方向性线索。

\textbf{第 5、7 篇单词测试：设计性退化}。检查后发现，
这两套单词测试的全部考题（10/10 和 5/5）都在教师 $\surd$ 标记内，
没有“未标记但同样考了”的对照项。\textbf{因此这两套数据无法识别教师判断效度}——
不是老师判断得准或不准，而是设计上没有对照。
\textbf{这是一个重要的产品/教研教训}：要检验“重点判断是否有效”，
测试必须同时包含\textbf{标记项}和\textbf{未标记的对照项}（foil）。

\textbf{学生自标元认知管线（等待真实数据）}：要测“学生知不知道自己哪里不会”，
必须让学生自己打 $\surd$。脚本已经建好：
\texttt{root\_mark\_observations\_template.csv} 是长表，
\texttt{根号自标采集卡.html} 是打印卡。
分析指标与上一节元认知一致：敏感性 AUC、偏差 $bias$、$d'$、L2 折扣。

\textbf{合成自检}（不是学生数据）：在已知真实 AUC 的情况下，
每人每次 10 题的估计 AUC 非常不稳定：

\begin{center}
\begin{longtable}{@{}lrrr@{}}
\toprule
\textbf{真实 AUC} & \textbf{个体 AUC 均值} & \textbf{个体 AUC SD} & \textbf{聚合 AUC} \\
\midrule
$0.50$ & $0.522$ & $0.171$ & $0.516$ \\
$0.65$ & $0.620$ & $0.182$ & $0.614$ \\
$0.75$ & $0.693$ & $0.141$ & $0.695$ \\
$0.85$ & $0.767$ & $0.148$ & $0.759$ \\
\bottomrule
\end{longtable}
\end{center}

零模型（真实 AUC $=0.5$）下，10 题个体 AUC 的 95\% 范围是 $0.143$--$0.798$。
\textbf{所以：单次 10 题的 AUC 只能当线索；要看个体差异，至少聚合 5 次测试。}

\subsubsection{实验 C：难度基线 v2——先修掉两个错误}

v1 有两个问题：
\begin{enumerate}[leftmargin=2em]
\item \textbf{题号类型错误}：9 篇和 7 篇句型的题号是整数，
脚本却用字符串去匹配，导致这两套题的难度 $p_i$ 被全部误算成 $0$。
\item \textbf{循环依赖}：学生残差用了“包含该学生自己”的题目难度，
自己会影响自己的基线。v2 改为 leave-one-out（留一法）难度。
\end{enumerate}

v2 的题目难度定义不变：

\begin{equation}
p_i=\frac{k_i}{N},\qquad
k_i=\text{第 }i\text{ 题做错的人数},\qquad
N=\text{参加该测试的总人数}.
\end{equation}

Wilson 95\% 区间用于小样本；它比常见的正态近似更稳。
v2 结果：

\begin{center}
\begin{longtable}{@{}lrrrr@{}}
\toprule
\textbf{测试} & \textbf{人数} & \textbf{项目数} & \textbf{平均难度} & \textbf{Wilson 宽度中位} \\
\midrule
第 9 篇 56--65 题 & 6 & 10 & $0.600$ & $0.603$ \\
第 7 篇 句型 1--5 & 6 & 5 & $0.500$ & $0.603$ \\
第 7 篇 单词 & 7 & 10 & $0.257$ & $0.523$ \\
第 5 篇 单词 & 6 & 5 & $0.233$ & $0.533$ \\
\bottomrule
\end{longtable}
\end{center}

\textbf{两种精度必须分开说}：

\begin{itemize}[leftmargin=2em]
\item \textbf{题目难度精度}取决于“每道题有多少学生”$N$。
现有 $N=6$--$7$，Wilson 宽度中位 $0.591$。
要把最坏宽度压到 $0.30$，需要每道题约 $39$ 人；$0.25$ 约 $58$ 人；$0.20$ 约 $93$ 人。
\item \textbf{学生能力精度}取决于“每个学生有多少题”$k$。
在 Rasch 模型下，$\mathrm{SE}(\hat\theta)\approx 1/\sqrt{\sum_i p_i(\theta)(1-p_i(\theta))}$。
现有 $k=10$ 时，解析 SE $\approx0.699$ logit；要到 $\approx0.308$，需要约 $50$ 题。
\textbf{学生人数再多也不能替代题数}。
\end{itemize}

\textbf{学生残差（v2，真实课堂数据）}：
用留一法难度算 $Z_{s}=(X_s-E_s)/\sqrt{V_s}$，
其中 $X_s$ 是实际错数，$E_s$ 是 LOO 预期错数，$V_s$ 是 LOO 预期方差；
再用 Monte Carlo 零模型给 $p$ 值。正残差 $=$ 比题目难度预期的更差；
负残差 $=$ 更好。合成的零模型自检显示：Wilson 覆盖率 $0.970$（目标 $0.95$），
但残差检验假阳性率 $0.088$（目标 $0.05$）——\textbf{偏高}。
所以下面的 $p<0.05$ 只能当\textbf{筛选线索}，不能当确证：

\begin{center}
\begin{longtable}{@{}llrrl@{}}
\toprule
\textbf{学生} & \textbf{测试} & \textbf{残差} & \textbf{Monte Carlo $p$} & \textbf{方向} \\
\midrule
陆怡婷 & 第 9 篇 56--65 & $+2.45$ & $0.0165$ & 比预期差 \\
张钊源 & 第 9 篇 56--65 & $-5.55$ & $<0.001$ & 比预期好 \\
刘泽宇 & 第 7 篇 句型 & $+3.20$ & $0.0040$ & 比预期差 \\
孙艺轩 & 第 7 篇 句型 & $+3.20$ & $0.0039$ & 比预期差 \\
孙艺轩 & 第 7 篇 单词 & $+3.62$ & $0.0008$ & 比预期差 \\
李陈浩然 & 第 5 篇 单词 & $+4.91$ & $<0.001$ & 比预期差 \\
\bottomrule
\end{longtable}
\end{center}

\textbf{一个必须修正的说法}：上一版“单词测试零区分度”来自某一批数据；
本次新增的第 5/7 篇单词测试\textbf{不是零区分度}（平均难度 $0.233$ 和 $0.257$）。
因此“单词测试零区分度”不能推广到所有单词测试，只能限定在原来那批数据上。

\subsubsection{三个实验怎么接回 Tender 和 DaoQL}

三个实验不是三份孤立报告，它们分别提供一种\textbf{状态测量}，
再写进 Tender 的会话和 DaoQL 的事件流：

\begin{center}
\begin{longtable}{@{}p{2.6cm}p{4.1cm}p{3.6cm}p{4.0cm}@{}}
\toprule
\textbf{实验} & \textbf{测量} & \textbf{Tender 节点} & \textbf{DaoQL 事件} \\
\midrule
A\ JOL & 预测 vs 实际 → 偏差、误差、Brier & monitoring（自我监控） & \texttt{metacognition\_jol} \\
B\ 根号 & 标记 vs 对错 → AUC、bias、$d'$、L2 折扣 & monitoring + L2 元元 & \texttt{metacognition\_root} \\
C\ 难度 & 逐题错误 → $p_i$、Wilson、学生残差 & content/item + student state & \texttt{item\_difficulty} + \texttt{student\_residual} \\
\bottomrule
\end{longtable}
\end{center}

\textbf{它们都不是情绪测量}。JOL 和根号回答的是“知不知道自己会不会”，
难度基线回答的是“这道题对这群人有多难”；情绪测量是另一条线，
不要把“元认知差”误读成“情绪差”，也不要用情绪模型替代元认知模型。

\subsection{一个意外发现：单词测试零区分度}

\begin{explainbox}{重要限定}
本节“单词测试零区分度”来自上一版的某一批单词数据。本次新增的第 5/7 篇单词测试
\textbf{并不是零区分度}（平均难度分别为 $0.233$ 和 $0.257$），所以这个发现只适用于
原来那批数据，\textbf{不能推广到所有单词测试}。
\end{explainbox}

在第 9 篇三项测试都参加的 6 名理科生上：

\begin{center}
\begin{tabular}{lrrr}
\toprule
\textbf{学生} & \textbf{填空错误率} & \textbf{单词错误率} & \textbf{句型错误率} \\
\midrule
陆怡婷 & 0.90 & 0.17 & 1.00 \\
孙欣 & 0.80 & \textbf{0.00} & 0.80 \\
王毅玮 & 0.50 & \textbf{0.00} & 1.00 \\
柴静轩 & 0.70 & \textbf{0.00} & 0.60 \\
\textbf{张钊源} & \textbf{0.00} & \textbf{0.00} & \textbf{1.00} \\
王昱琳 & 0.70 & \textbf{0.00} & 0.20 \\
\bottomrule
\end{tabular}
\end{center}

\begin{explainbox}{两个发现}
\textbf{(1) 单词测试几乎全体满分}——它目前在\textbf{提供零信息}，
既不能区分学生，也不能反映进步。这是\textbf{在浪费课堂时间}。\\[4pt]
\textbf{(2) 张钊源：填空全对、单词全对、句型全错}——
这不是“英语差”，是\textbf{明确的单点短板}。
而这类学生恰恰是元认知测量最有价值的对象：\textbf{他自己知道句型弱吗？}
\end{explainbox}

\subsection{边界与不做的事}

\begin{longtable}{@{}p{6.2cm}p{7.4cm}@{}}
\toprule
\textbf{不做} & \textbf{原因} \\
\midrule
不用“求助行为”当元认知信号 & 它是结果之后的反应，不携带额外信息 \\
不用自评问卷测元认知 & Double (2025)：问卷测元认知常常是假的 \\
不把相关当因果 & $r=0.513$ 不等于“训练元认知就能提分” \\
不给学生贴“元认知差”标签 & 元认知可训练（Huff \& Nietfeld 2009），标签会造成自证预言 \\
不在低敏感性时强推自评 & 敏感性低的学生，其自评应被降权而非直接采信 \\
\bottomrule
\end{longtable}

\newpage

\section{元元检测：给结论本身打可信度分}
% ============================================================

前面所有章节都在回答“值是多少”。这一节换一个问题：\textbf{“这个值有多可信？”}

这个问题在本项目里不是新问题——L2 层早就有一个专门的模块做这件事，叫\textbf{元元检测}。
本节做的，是把同一套方法用到“研究结论”上。

\subsection{什么是“元元”}

先把三个词分清楚：

\begin{center}
\begin{tabular}{p{2.4cm}p{4.6cm}p{6.6cm}}
\toprule
\textbf{词} & \textbf{做什么} & \textbf{例子} \\
\midrule
测量 & 观察一个对象，得到一个值 & “你今天 $\theta = 1.2$” \\
元测量 & 测量“这次测量有多可信” & “这个 $\theta$ 的可信度是 0.72” \\
\textbf{元元} & 对“测量的测量”再做一层 & 见下面的层次图 \\
\bottomrule
\end{tabular}
\end{center}

\begin{explainbox}{项目里的三层结构}
\begin{center}
\begin{tabular}{p{1.4cm}p{3.4cm}p{8.8cm}}
\toprule
\textbf{层} & \textbf{名字} & \textbf{做什么} \\
\midrule
L1 & 状态推断 & 从颜色、文本、序列合成 $\theta$ \\
\textbf{L2} & \textbf{元元检测} & \textbf{判断“这个 $\theta$ 推得准不准”} \\
L3 & 策略与干预 & 根据 $\theta$ \textbf{和}可信度，决定要不要干预 \\
\bottomrule
\end{tabular}
\end{center}
关键在 L2：它不产生新的状态值，它产生\textbf{一个关于可信度的评分}。
只有在可信度够高时，L3 才允许做出干预。
\end{explainbox}

\subsection{L2 的元元公式（原文，逐符号解释）}

\subsubsection{综合可信度}

\begin{equation}
\text{MetaConfidence} = 1 - \sum_{k=1}^{4} \beta_k \cdot \text{alert}_k
\label{p3:eq:mc-l2}
\end{equation}

\begin{explainbox}{每一个符号}
\begin{longtable}{@{}p{3.2cm}p{10.6cm}@{}}
$\text{MetaConfidence}$ & 综合可信度。取值 $[0,1]$，\textbf{越大越可信} \\
$k$ & 检测器的编号，从 1 到 4 \\
$\sum_{k=1}^{4}$ & 求和符号：把 $k=1,2,3,4$ 四项加起来 \\
$\beta_k$ & 第 $k$ 个检测器的\textbf{权重}。所有 $\beta_k$ 之和为 1 \\
$\text{alert}_k$ & 第 $k$ 个检测器是否报警。取值只有 0（没报）或 1（报了） \\
$1 - (\cdots)$ & 可信度 = 1 减去“所有报警的加权和” \\
\end{longtable}
\textbf{读法}：一个警报都没响时，MetaConfidence $= 1$（完全可信）；
警报越多、权重越大，可信度越低。
\end{explainbox}

\subsubsection{四个检测器与黄金比阈值谱系}

\begin{equation}
\varphi^{-3} \approx 0.236, \quad
\varphi^{-2} \approx 0.382, \quad
\varphi^{-1} \approx 0.618, \quad
\varphi^{1} \approx 1.618
\end{equation}

\begin{center}
\begin{tabular}{p{3.6cm}p{2.8cm}p{8.2cm}}
\toprule
\textbf{检测器} & \textbf{阈值} & \textbf{检查什么} \\
\midrule
数据充分性 & $\varphi^{-1} \approx 0.618$ & 数据量够不够 \\
群体一致性 & $\varphi^{-3} \approx 0.236$ & 群体里大家一致吗 \\
个人历史表现 & $\varphi^{-2} \approx 0.382$ & 这个人的历史是否自洽 \\
状态偏离度 & $\varphi^{1} \approx 1.618$ & 当前状态是否偏离太远 \\
\bottomrule
\end{tabular}
\end{center}

\begin{explainbox}{为什么阈值要用黄金比的幂}
黄金比 $\varphi$ 满足 $\varphi^2 = \varphi + 1$，且 $\varphi^{-1} = \varphi - 1$。
用它做阈值谱系有两个好处：
\begin{enumerate}[leftmargin=1.4em,itemsep=1pt]
\item \textbf{尺度自相似}：每一档都是上一档的 $\varphi$ 倍，
      所以“松/紧”的划分在任何尺度上看起来都一样；
\item \textbf{全系统统一}：L1 的临界角、L2 的阈值都出自同一个 $\varphi$，
      系统里没有第二套常数。
\end{enumerate}
\end{explainbox}

\subsubsection{弱测量算符}

\begin{equation}
M_{\text{meta}} = I - \eta \cdot \big(1 - \text{MetaConfidence}\big) \cdot P_{\text{模糊}}
\label{p3:eq:weak}
\end{equation}

\begin{explainbox}{每一个符号}
\begin{longtable}{@{}p{3.4cm}p{10.4cm}@{}}
$M_{\text{meta}}$ & 弱测量算符。作用在状态上产生一个\textbf{很轻的修正} \\
$I$ & 单位算符。“什么都不做”——作用在状态上原样返回 \\
$\eta$ & 弱测量强度系数。L2 里取 $\eta = 0.01$，\textbf{非常小} \\
$1 - \text{MetaConfidence}$ & \textbf{不可信程度}。可信度越低，这个数越大 \\
$P_{\text{模糊}}$ & 投影到“模糊态”的算符。把状态往“不确定”方向拉一点 \\
$\cdot$ & 乘法（标量乘算符） \\
\end{longtable}
\textbf{读法}：
\begin{itemize}[leftmargin=1.4em,itemsep=1pt]
\item 可信度 $=1$ 时：$1-\text{MetaConfidence}=0$，于是 $M_{\text{meta}}=I$，\textbf{什么也不做}；
\item 可信度低时：往“模糊态”收缩一点点。
\end{itemize}
\textbf{为什么叫“弱”测量}：$\eta$ 只有 0.01，修正极小。
它的目的是\textbf{标记不确定性}，而不是\textbf{改变状态}。
\end{explainbox}

\subsection{观测者效应：它本身就是元元问题}

第 6 节讲 MBTI 时提到三条设计（不告知假设、不反馈分数、不反馈偏离度）。
当时是以“伦理要求”给出的。这里补上它的\textbf{方法论根据}。

\begin{explainbox}{观测者效应的本质}
\begin{center}
\begin{tabular}{p{6.0cm}p{6.4cm}}
\toprule
\textbf{普通测量} & \textbf{观测者效应} \\
\midrule
观察者 $\to$ 被观察者 $\to$ 得到一个值 &
观察者 $\to$ 被观察者 $\to$ \textbf{被观察者知道自己被观察}\\
& $\to$ 被观察者改变 $\to$ \textbf{值已经不是原来的值} \\
\bottomrule
\end{tabular}
\end{center}
一句话：\textbf{测量过程回流进了被测量的系统。}
\end{explainbox}

\begin{explainbox}{所以它和元元是同一件事}
元元层要检测的，正是“测量是否被自身污染”。
所以观测者效应不是元元的“邻居”，\textbf{它就是元元的对象本身}：

\begin{center}
\begin{tabular}{p{4.4cm}p{8.8cm}}
\toprule
\textbf{元元层里的概念} & \textbf{在 MBTI 研究里的化身} \\
\midrule
弱测量（尽量不扰动系统） & \textbf{不告知假设、不反馈结果} \\
MetaConfidence & 对“研究结论有多可信”的评分 \\
阈值谱系 & 分级线：可作结论 / 只能作线索 / 不可作结论 \\
元元检测器 & \textbf{阴性对照 H7} \\
护栏 & 事后改动、系统性偏差 \\
\bottomrule
\end{tabular}
\end{center}
\end{explainbox}

\newpage
\subsection{把元元用到研究结论上}

\begin{explainbox}{为什么研究结论也需要 MetaConfidence}
一份研究报告会说“$\rho = 0.42$，$p = 0.01$”。
但它不会说“这个数是在 19 个人身上得到的，而且学生知道我们在验证什么”。\\[4pt]
\textbf{前者是测量值，后者是可信度。两个都要报。}
\end{explainbox}

\subsubsection{公式}

\begin{equation}
\text{MC}_{\text{study}} = \sum_{k=1}^{5} w_k \cdot s_k
\label{p3:eq:mc-study}
\end{equation}

\begin{explainbox}{每一个符号}
\begin{longtable}{@{}p{2.6cm}p{11.2cm}@{}}
$\text{MC}_{\text{study}}$ & 研究结论的元元可信度，取值 $[0,1]$ \\
$k$ & 维度编号，1 到 5 \\
$w_k$ & 第 $k$ 个维度的\textbf{权重}，且 $\sum_{k=1}^{5} w_k = 1$ \\
$s_k$ & 第 $k$ 个维度的\textbf{得分}，取值 $[0,1]$，越大越好 \\
\end{longtable}
注意 $s_k$ 和 L2 的 $\text{alert}_k$ \textbf{含义正好相反}：
$s_k$ 是“做得多好”，$\text{alert}_k$ 是“有没有报警”。
两者的关系是 $s_k = 1 - \text{alert}_k$。
\end{explainbox}

\begin{center}
\begin{tabular}{p{4.2cm}p{1.8cm}p{8.2cm}}
\toprule
\textbf{维度} & \textbf{权重 $w_k$} & \textbf{评什么} \\
\midrule
观测者效应风险 & 0.25 & 学生是否知情 / 是否反馈结果 / 记录是否可见 \\
功效 & 0.25 & 样本量 $N$ 够不够 \\
阴性对照 & 0.20 & 本该没效应的检验是否干净 \\
数据质量 & 0.15 & 覆盖天数、活动次数 \\
预注册依从 & 0.15 & 有没有事后改动 \\
\bottomrule
\end{tabular}
\end{center}

\subsubsection{与 L2 公式是同一个式子}

\begin{explainbox}{一个值得注意的等价关系}
把 $s_k = 1 - \text{alert}_k$ 代入式 \eqref{p3:eq:mc-study}：

\begin{equation}
\sum_{k} w_k s_k
= \sum_{k} w_k (1 - \text{alert}_k)
= \underbrace{\sum_{k} w_k}_{=1} - \sum_{k} w_k \text{alert}_k
= 1 - \sum_{k} w_k \text{alert}_k
\end{equation}

最后这个式子\textbf{就是 L2 的式 \eqref{p3:eq:mc-l2}}（把 $\beta_k$ 换成 $w_k$）。\\[4pt]
\textbf{所以本节的评分卡不是新公式，而是 L2 元元公式在“研究结论”这个对象上的应用。}
\end{explainbox}

\begin{explainbox}{唯一的区别：连续性}
\begin{center}
\begin{tabular}{p{3.6cm}p{4.2cm}p{5.4cm}}
\toprule
 & \textbf{L2} & \textbf{本研究} \\
\midrule
报警量 & $\text{alert}_k \in \{0,1\}$（二值） & $s_k \in [0,1]$（连续） \\
理由 & 状态推断只有“够/不够” & 设计质量可以“部分达标” \\
\bottomrule
\end{tabular}
\end{center}
当所有 $s_k$ 只取 0 或 1 时，本研究的形式\textbf{退化成 L2 的原始形式}。
\end{explainbox}

\begin{center}
\begin{tabular}{p{4.2cm}p{3.0cm}p{7.2cm}}
\toprule
$\text{MC}_{\text{study}}$ & \textbf{分级} & \textbf{可以说的话} \\
\midrule
$\ge 0.80$ & 可作结论 & 可以作为研究结论对外表述 \\
$0.50 \sim 0.80$ & 只能作线索 & 只能作为“值得进一步看”的线索 \\
$< 0.50$ & 不可作结论 & 结论无效，需重新设计或重新采集 \\
\bottomrule
\end{tabular}
\end{center}

\subsection{护栏：为什么加权平均会失效}

\begin{explainbox}{一个真实的失败}
第一版评分卡跑出来，“看完数据后改了指标”这个场景
\textbf{居然得到 0.800 —— 被判为“可作结论”}。\\[4pt]
原因：其他四项都满分，加权平均把\textbf{单项致命缺陷稀释掉了}：
\begin{equation}
0.25(0.80) + 0.25(1.0) + 0.20(1.0) + 0.15(1.0) + 0.15(0) = 0.800
\end{equation}
\textbf{这是评分卡自身的缺陷}——也就是“元元自己也需要元元”。
\end{explainbox}

\subsubsection{修正：一票否决}

\begin{equation}
\text{MC}_{\text{final}} = \min\Big(\text{MC}_{\text{study}},\; V_{\text{veto}}\Big)
\label{p3:eq:veto}
\end{equation}

\begin{explainbox}{每一个符号}
\begin{longtable}{@{}p{3.0cm}p{10.8cm}@{}}
$\text{MC}_{\text{final}}$ & 最终可信度 \\
$\min(\cdot,\cdot)$ & 取较小值。这是关键的\textbf{封顶算子} \\
$\text{MC}_{\text{study}}$ & 式 \eqref{p3:eq:mc-study} 算出的加权分 \\
$V_{\text{veto}}$ & 封顶值。触发致命缺陷时取一个小值 \\
\end{longtable}
\end{explainbox}

\begin{center}
\begin{tabular}{p{4.6cm}p{3.2cm}p{6.8cm}}
\toprule
\textbf{触发条件} & \textbf{封顶 $V_{\text{veto}}$} & \textbf{理由} \\
\midrule
看过数据后改过假设或指标 & \textbf{0.30} & 事后挑选，结论作废 \\
阴性对照显著（$p<0.01$） & \textbf{0.45} & 存在系统性偏差 \\
\bottomrule
\end{tabular}
\end{center}

\begin{explainbox}{为什么用 $\min$ 而不用减法}
如果用减法（$\text{MC} - c$），只要其他项足够高，仍然可能回到“可作结论”区间——
\textbf{缺陷仍然可补偿}。\\[4pt]
而 $\min$ 表达的是\textbf{“不可补偿”}：
无论其他维度多完美，只要触发红线，结论就封顶。\\[4pt]
这和项目《指标验证与商业路线图》里的\textbf{“护栏（一票否决）”}是同一个设计原则。
\end{explainbox}

\subsection{六个场景的实测对比}

\begin{longtable}{@{}p{4.0cm}rrrrrrl@{}}
\toprule
\textbf{场景} & \textbf{观测} & \textbf{功效} & \textbf{对照} & \textbf{质量} & \textbf{预注册} & \textbf{MC} & \textbf{分级} \\
\midrule
理想执行（$N=90$） & 0.80 & 1.00 & 1.00 & 1.00 & 1.00 & \textbf{0.950} & 可作结论 \\
现实规模（$N=19$） & 0.80 & 0.15 & 1.00 & 1.00 & 1.00 & \textbf{0.738} & 只能作线索 \\
告知了假设 & 0.35 & 0.15 & 1.00 & 1.00 & 1.00 & 0.625 & 只能作线索 \\
反馈了分数 & 0.00 & 0.15 & 1.00 & 1.00 & 1.00 & 0.537 & 只能作线索 \\
\textbf{阴性对照失败} & 0.80 & 1.00 & \textbf{0.00} & 1.00 & 1.00 & \textbf{0.450} & \textbf{不可作结论} \\
\textbf{看完数据改指标} & 0.80 & 1.00 & 1.00 & 1.00 & \textbf{0.00} & \textbf{0.300} & \textbf{不可作结论} \\
\bottomrule
\end{longtable}

\begin{explainbox}{三条读数}
\begin{enumerate}[leftmargin=1.4em,itemsep=2pt]
\item \textbf{理想执行与现实规模差 0.212，全部来自“功效”。}\\
      其余四项完美，样本不够仍然只能得到“线索”级结论。
\item \textbf{观测者效应是最容易失守的一环。}\\
      “告知假设”一项就掉 0.113；再加“反馈分数”掉到 \textbf{0.537}——
      已逼近 0.50 门槛。\textbf{只要既告知假设又反馈结果，这项研究基本就废了。}
\item \textbf{两个致命缺陷都被护栏拦住}——包括那个用加权平均差点漏掉的“事后改指标”。
\end{enumerate}
\end{explainbox}

\subsection{边界：可信度不等于效应量}

\begin{explainbox}{两者必须分开报}
\textbf{元元评分只评估“研究做得对不对”，不评估“效应大不大”。}
\begin{itemize}[leftmargin=1.4em,itemsep=1pt]
\item $\text{MC} = 0.95$ 的结论，可能效应只有 0.05，\textbf{没有实用价值}；
\item $\text{MC} = 0.55$ 的结论，可能效应很大，\textbf{只是样本不够}。
\end{itemize}
所以报告必须同时给出：

\begin{center}
\begin{tabular}{ll}
\toprule
\textbf{分量} & \textbf{示例} \\
\midrule
效应量 & $\rho = 0.42$，95\% CI $[0.22, 0.58]$ \\
元元可信度 & $\text{MC} = 0.74$（只能作线索：样本仅 19 人） \\
\bottomrule
\end{tabular}
\end{center}
\textbf{只报前者会误导}（好像很确定）；\textbf{只报后者也误导}（好像很小）。
\end{explainbox}

\subsection{小结：这一节连接了两件事}

\begin{center}
\begin{tabular}{p{5.0cm}p{8.2cm}}
\toprule
\textbf{连接} & \textbf{说明} \\
\midrule
观测者效应 $\leftrightarrow$ 元元 &
观测者效应的本质是“测量回流进系统”，这正是元元要检测的对象 \\
L2 元元公式 $\leftrightarrow$ 本研究评分卡 &
$s_k = 1-\text{alert}_k$ 时两者\textbf{完全等价}，本研究是其连续推广 \\
弱测量 $\leftrightarrow$ 研究设计 &
不告知假设、不反馈结果 = 对研究对象的“弱测量” \\
护栏 $\leftrightarrow$ 一票否决 &
两者都用\textbf{不可补偿的红线}，而不是加权扣分 \\
\bottomrule
\end{tabular}
\end{center}


\subsection{DAOQL 客观事件与生命树主观体验}

这一版新增一个关键分层：\textbf{DAOQL 记录客观世界，生命树记录主观体验。}

\begin{center}
\begin{longtable}{@{}p{3.2cm}p{5.0cm}p{5.6cm}@{}}
\toprule
\textbf{层} & \textbf{回答什么} & \textbf{保存什么} \\
\midrule
DAOQL 客观层 & 谁在什么时间做了什么，来源和证据是什么？ & 原始事件、时间、来源、证据、版本与审计 \\
metameta 主观体验层 & 这个事件在生命树中可能意味着什么？ & 体验节点、脑图区域、记忆类型、置信度与不确定性 \\
\bottomrule
\end{longtable}
\end{center}

新增 Agent \texttt{metameta\_cartographer} 读取客观事件，生成两类派生记录：

\begin{itemize}[leftmargin=2em]
\item \texttt{subjective\_experience\_node}：把一次事件拆解到工作记忆、海马体索引、长期记忆、内省反馈、情绪、代谢治理等区域；
\item \texttt{subjective\_experience\_summary}：保存主区域、最多 6 个标签、生命树阶段和体验节点 ID。
\end{itemize}

关系边明确写成：

\begin{verbatim}
客观事件 --experienced_as--> 主观体验节点
客观事件 --yields_subjective_experience--> 体验摘要
体验摘要 --contains_experience_node--> 体验节点
体验节点 --索引/编码/内省/表达等关系--> 体验节点
\end{verbatim}

每个体验节点都带 \texttt{epistemic\_status = subjective\_experience}、\texttt{derived = true}、客观来源事件 ID、匹配证据词、置信度、不确定性和生命树阶段。客观事件保持原样；体验节点可以更正、重算和撤回。

接口为：

\begin{verbatim}
POST /metameta/decompose/{student_id}
GET  /metameta/map/{student_id}
\end{verbatim}

\begin{explainbox}{边界}
这些区域是认知隐喻与检索标签，不是脑区扫描、临床诊断或客观事实。生命树回答“如何被体验”，DAOQL 回答“发生了什么”。两者通过来源边相连，不能互相冒充。
\end{explainbox}

\newpage

\section{犹豫：元认知的行为证据}
% ============================================================

\subsection{犹豫本身不是元认知}

犹豫是一个行为量。它可以是：

\begin{itemize}[leftmargin=2em]
\item 响应时间变长；
\item 停顿、回看、修改；
\item 鼠标或键盘反复操作；
\item 语言标记：“嗯”“可能”“应该是”“我不确定”。
\end{itemize}

但行为量不等于元认知。犹豫只有在满足两个条件时，才算\textbf{元认知信号}：

\begin{enumerate}[leftmargin=2em]
\item \textbf{发生在结果之前}：如果学生已经看到答案才犹豫，那是事后反应，不是事前监控；
\item \textbf{能预测错误，且超出题目难度本身}：如果只是题长、题难导致慢，那不是元认知。
\end{enumerate}

\subsection{形式定义：犹豫的元认知敏感性}

把“犹豫程度”记作 $h$，实际对错记作 $e\in\{0,1\}$，其中 $e=1$ 表示答错。
定义：

\begin{equation}
\mathrm{AUC}_{\text{hesitation}}
=
\Pr(h_{\text{错题}}>h_{\text{对题}})
+\frac12\Pr(h_{\text{错题}}=h_{\text{对题}}).
\end{equation}

逐符号说明：

\begin{itemize}[leftmargin=2em]
\item $h_{\text{错题}}$：答错题上的犹豫程度；
\item $h_{\text{对题}}$：答对题上的犹豫程度；
\item $\Pr(\cdot)$：概率；
\item $\mathrm{AUC}=0.5$：犹豫与对错无关，不是元认知信号；
\item $\mathrm{AUC}>0.5$：犹豫确实携带“我不确定”的信息；
\item $\mathrm{AUC}=1.0$：犹豫完美区分错题与对题。
\end{itemize}

\subsection{真实数据：ASSISTments 的响应时间}

第三期已经在 ASSISTments 上做过这个实验：

\begin{center}
\begin{longtable}{@{}p{5.0cm}p{8.5cm}@{}}
\toprule
\textbf{信号} & \textbf{真实结果} \\
\midrule
响应时间（首次作答前） & 元认知敏感性中位 $\mathrm{AUC}\approx0.585$；犹豫确实携带信息 \\
求助行为 hint\_count & 求助过的 40{,}042 条里只有 9 条答对；它是“卡住后的反应”，不是事前判断 \\
\bottomrule
\end{longtable}
\end{center}

\begin{explainbox}{关键区分}
\textbf{事前犹豫}（响应时间、首次作答前停顿）可以是元认知信号。\\
\textbf{事后求助}（hint、看答案、问老师）大多是结果之后的控制行为，不能当元认知判断。
\end{explainbox}

\subsection{怎么引导应用：犹豫 $\times$ 对错四象限}

这是第四期最实用的一张表。把每个学生每道题放进四象限：

\begin{center}
\begin{longtable}{@{}p{2.4cm}p{5.2cm}p{5.4cm}@{}}
\toprule
 & \textbf{实际答对} & \textbf{实际答错} \\
\midrule
\textbf{高犹豫} & \textbf{知识脆弱}：会，但不稳。安排变式练习、间隔复习、用自己的话复述。 & \textbf{知识缺口}：不会。先问“卡在哪一步”，再给最小提示、示范、拆步。 \\
\textbf{低犹豫} & \textbf{流畅掌握}：可以提速、进阶、做迁移题。 & \textbf{错误流畅}：很确定地错。用反例、对比题、同伴互评制造认知冲突。 \\
\bottomrule
\end{longtable}
\end{center}

\subsection{四种干预的详细脚本}

\subsubsection{高犹豫 + 答错：知识缺口}

不要直接给答案。按四步：

\begin{enumerate}[leftmargin=2em]
\item 让学生指出卡点：“你是哪一步开始不确定的？”
\item 给最小提示，不展开完整讲解；
\item 让学生自己补完这一步；
\item 换一道同结构题，验证是否真的补上。
\end{enumerate}

对应 T-ACR 的“注入 $\to$ 扰动 $\theta$”。

\subsubsection{高犹豫 + 答对：知识脆弱}

学生会，但靠的是临时回忆，不是稳定结构。安排：

\begin{itemize}[leftmargin=2em]
\item 间隔 1 天、3 天、7 天各做一次变式；
\item 要求学生解释“为什么选这个形式”；
\item 不要马上加难度，先把稳定性做出来。
\end{itemize}

对应“驻波耦合 $\to$ 裸传压缩”。

\subsubsection{低犹豫 + 答错：错误流畅}

这是最危险的一类。学生“很确定地错”，重复讲一遍往往没用。应该：

\begin{itemize}[leftmargin=2em]
\item 给反例，让学生自己发现矛盾；
\item 让同伴互评，先让学生听到不同答案；
\item 用“预测—实际”冲突：先写预测，再看实际结果。
\end{itemize}

对应“强制再观测”。

\subsubsection{低犹豫 + 答对：流畅掌握}

可以：

\begin{itemize}[leftmargin=2em]
\item 换情境、换题型；
\item 让学生当小老师；
\item 进入下一层布鲁姆层级（分析、评价、创造）。
\end{itemize}

对应“结果”——让 $\theta$ 基线台阶式上移。

\subsection{电子化采集要记什么}

\begin{center}
\begin{longtable}{@{}p{3.6cm}p{4.2cm}p{6.2cm}@{}}
\toprule
\textbf{字段} & \textbf{含义} & \textbf{注意} \\
\midrule
首次作答时间 & 从题目呈现到第一次输入 & 要排除网络延迟和题长 \\
修改次数 & 答案被改了几次 & 反复改是犹豫的强信号 \\
回看次数 & 回看题干或材料的次数 & 需要平台支持 \\
信心条 & 学生自报 0--100 & 与响应时间互补 \\
语言标记 & 可能/应该/不确定/嗯 & 中文可用 hedging 词表 \\
题目难度 & 该题错误率 $p_i$ & 必须控制，否则会把题难误判成犹豫 \\
\bottomrule
\end{longtable}
\end{center}

\begin{explainbox}{边界}
犹豫也受疲劳、网络、题目长度影响。\\
至少要控制：题目难度、题长、题序、学生基线反应时。\\
不能因为学生慢，就默认他不会。
\end{explainbox}

\newpage

% ============================================================

\section{外部脚手架：元认知能不能“植入”？}
% ============================================================

\subsection{先分清三件事}

“外部植入”这个词太粗。要拆成三层：

\begin{center}
\begin{longtable}{@{}p{2.8cm}p{4.4cm}p{6.8cm}@{}}
\toprule
\textbf{层次} & \textbf{做法} & \textbf{改变的是什么} \\
\midrule
引出 elicit & 问“你估计能对几题”“你有多确定” & 把内部信号变成可观测数据 \\
增强 augment & 给校准反馈、复盘、对照 & 可能真正改善监控敏感性 \\
锚定 anchor & 暗示“你应该有 80\% 信心” & 大多只改变报告，不一定改变内部监控 \\
\bottomrule
\end{longtable}
\end{center}

\subsection{形式定义：报告信号与内部信号}

把内部监控记作 $m_{\text{internal}}$，报告出来的信号记作 $m_{\text{report}}$，外部输入记作 $E$：

\begin{equation}
m_{\text{report}}=g(m_{\text{internal}},E).
\end{equation}

符号说明：

\begin{itemize}[leftmargin=2em]
\item $m_{\text{internal}}$：学生真实的“我有多确定”；
\item $m_{\text{report}}$：学生写出来的信心、JOL、标记；
\item $E$：外部提示、反馈、同伴信息、教师暗示；
\item $g$：从内部到报告的映射。
\end{itemize}

判断标准：

\begin{itemize}[leftmargin=2em]
\item 如果 $E$ 只改变 $m_{\text{report}}$，拿掉 $E$ 就恢复原样 $\to$ 这是\textbf{锚定}；
\item 如果 $E$ 改变了 $m_{\text{internal}}$，在\textbf{新任务、无提示}下仍然改善 $\to$ 这是\textbf{内部化}。
\end{itemize}

\begin{explainbox}{核心结论}
不能“植入”第一人称的元认知体验；\\
但可以外部\textbf{脚手架化}：引出、反馈、训练、迁移。
\end{explainbox}

\subsection{三臂实验：无提示 vs JOL vs JOL + 反馈}

这是可以直接在语法填空课堂做的实验。

\begin{center}
\begin{longtable}{@{}p{2.4cm}p{4.4cm}p{7.0cm}@{}}
\toprule
\textbf{组别} & \textbf{每次测试做什么} & \textbf{要检验什么} \\
\midrule
A 对照 & 正常测试，不提示 & 基线校准水平 \\
B JOL & 测前写“估计能对几题” & 光有提示能不能改善 \\
C JOL + 反馈 & 测前 JOL，批改后给“预测 vs 实际”反馈 & 反馈能不能改善敏感性，而不只是偏差 \\
\bottomrule
\end{longtable}
\end{center}

结果看三件事：

\begin{enumerate}[leftmargin=2em]
\item \textbf{偏差} $b=\hat p-p$：C 组是否降低过度自信？
\item \textbf{敏感性 AUC}：C 组是否提高对“会/不会”的区分？
\item \textbf{迁移}：撤掉反馈后，下一轮是否仍然更准？
\end{enumerate}

\begin{explainbox}{判读规则}
只有偏差改善、AUC 不变、撤掉反馈就恢复 $\to$ 那是锚定，不是元认知训练。\\
AUC 上升且迁移到新任务 $\to$ 才能说外部脚手架促进了内部监控。
\end{explainbox}

\subsection{伦理边界}

\begin{itemize}[leftmargin=2em]
\item 不要伪造信心；
\item 不要用假反馈操纵学生；
\item 不要把“外部提示”包装成“学生自己的判断”；
\item 外部植入只能做“引出 + 反馈 + 训练”，不能做欺骗性锚定。
\end{itemize}

\subsection{功效提醒}

第三期已经算过：19 名学生 $\times$ 5 次测试，相关 $r=0.5$ 时功效约 $0.61$，$r=0.6$ 时约 $0.82$。
所以三臂实验更适合看\textbf{组间方向}和\textbf{个体内变化}，不要指望一次就测出小效应。

\newpage
% ============================================================

\section{教师观察者与 metameta}
% ============================================================

\subsection{元认知的三层结构}

\begin{enumerate}[leftmargin=2em]
\item \textbf{一阶：学生自我监控}——学生自己预测、自己标重点（JOL、根号自标）；
\item \textbf{二阶：教师作为观察者}——教师判断“这个学生哪里不会、现在是什么状态”；
\item \textbf{三阶：metameta}——教师对自己判断的校准：教师不仅要有判断，还要知道自己判断得准不准。
\end{enumerate}

\subsection{教师观察的形式定义}

对教师 $O$ 的判断质量：

\begin{equation}
\mathrm{AUC}_{O}
=
\Pr(\text{教师标记的错题}>\text{教师标记的对题})
+\frac12\Pr(\text{相等}).
\end{equation}

\begin{equation}
\mathrm{MC}_{O}^{\text{adj}}
=
\mathrm{MC}_{O}\times \mathrm{discount}(\mathrm{AUC}_{O}).
\end{equation}

符号说明：

\begin{itemize}[leftmargin=2em]
\item $\mathrm{AUC}_O=0.5$：教师标记与对错无关，观察没有信息；
\item $\mathrm{AUC}_O=1.0$：教师标记完美对应错题；
\item $\mathrm{MC}_O$：教师自报的观察置信度；
\item $\mathrm{discount}(\cdot)$：把教师观察折进学生状态时使用的可信度折扣；
\item $\mathrm{MC}_O^{\text{adj}}$：折后置信度，才是 Tender 真正应该采信的教师观察权重。
\end{itemize}

\subsection{实验 B 的真实先验}

第三期实验 B 的教师题目级判断：

\begin{equation}
\mathrm{AUC}_{\text{teacher}}=0.5625,\quad
p=0.795,\quad
95\%\ \mathrm{CI}=[0.333,0.812].
\end{equation}

按当前 confidence\_discount 映射：

\begin{equation}
\mathrm{discount}(0.5625)
=
0.30+0.70\times\frac{0.5625-0.5}{0.9-0.5}
\approx 0.409.
\end{equation}

\begin{explainbox}{解释}
教师重点判断目前只能被当作\textbf{弱证据}，不能当作 ground truth。\\
这不是说教师不重要，而是说教师观察者本身必须被 metameta 校准。
\end{explainbox}

\subsection{E1：教师--学生--观察者三方对照}

\subsubsection{设计}

每次测试采集：

\begin{enumerate}[leftmargin=2em]
\item 学生预测：我估计能对几题（0--10）；
\item 教师预测：我估计该生能对几题（0--10）；
\item 教师信心：0--100；
\item 实际对题数。
\end{enumerate}

\subsubsection{指标}

\begin{itemize}[leftmargin=2em]
\item 学生偏差 $b_s=\hat p_s-p$；
\item 教师偏差 $b_t=\hat p_t-p$；
\item 学生校准误差 $a_s=|b_s|$；
\item 教师校准误差 $a_t=|b_t|$；
\item 教师 metameta $=\mathrm{corr}(\text{教师信心},-a_t)$。
\end{itemize}

\subsubsection{合成功效（19 名学生 $\times$ 5 次测试）}

\begin{center}
\begin{longtable}{@{}rr@{}}
\toprule
\textbf{教师 MAE 优势} & \textbf{功效} \\
\midrule
$0.00$ & $0.047$ \\
$0.02$ & $0.121$ \\
$0.04$ & $0.377$ \\
$0.06$ & $0.682$ \\
$0.08$ & $0.891$ \\
$0.10$ & $0.978$ \\
\bottomrule
\end{longtable}
\end{center}

相关功效：

\begin{center}
\begin{longtable}{@{}rr@{}}
\toprule
\textbf{真实相关 $r$} & \textbf{功效} \\
\midrule
$0.2$ & $0.131$ \\
$0.3$ & $0.249$ \\
$0.4$ & $0.406$ \\
$0.5$ & $0.612$ \\
$0.6$ & $0.816$ \\
$0.7$ & $0.948$ \\
\bottomrule
\end{longtable}
\end{center}

\subsection{E5：多教师 metameta}

现有项目只有一位教师，无法估计教师间差异。第四期做了合成功效。

\subsubsection{一次真实的翻车}

最初的 Fisher-$z$ 方差检验把相关估计的采样方差当成已知常数，导致功效低于名义水平。
后来改用 \textbf{ANOVA 检验教师间校准误差差异}，这才是正确做法。

\subsubsection{修正后的功效}

\begin{center}
\begin{longtable}{@{}lrrrr@{}}
\toprule
\textbf{教师间真实差异 $sd_\tau$} & \textbf{5 位} & \textbf{10 位} & \textbf{20 位} & \textbf{50 位} \\
\midrule
$0.3$ & $0.903$ & $0.987$ & $1.000$ & $1.000$ \\
$0.5$ & $0.970$ & $0.990$ & $1.000$ & $1.000$ \\
$0.8$ & $0.983$ & $1.000$ & $1.000$ & $1.000$ \\
\bottomrule
\end{longtable}
\end{center}

\begin{explainbox}{设计结论}
只要每位教师有 19 名学生 $\times$ 5 次测试（95 个观测），\\
\textbf{5 位教师就足以检出中等以上的教师间差异}。\\
门槛不是“几百位教师”，而是“至少 5 位、每位足够多观测”。
\end{explainbox}

\subsection{怎么用教师 metameta}

\begin{enumerate}[leftmargin=2em]
\item 教师观察不要直接当 ground truth，先乘 $\mathrm{discount}(\mathrm{AUC}_O)$；
\item 教师信心不是教师能力，它本身也要被校准；
\item 教师在场会产生观测者效应，所以还要记录“教师在什么观察方式下更准”；
\item 多教师实验至少 5 位，最好 10--20 位。
\end{enumerate}

\newpage
% ============================================================


\part{群体、平台与供应链}

\section{实验一：平台逆向实验}
% ============================================================

\subsection{要解决什么问题}

内容创作者想知道：\hl{同一条内容，发到不同平台，效果会不会不一样？怎么改能更好？}

这个问题不能靠“平台内部算法”回答（那是拿不到的），
但可以靠\textbf{公开能看到的数据}建立一个近似模型。

\begin{explainbox}{本实验的界限}
\textbf{做什么}：用公开可观测的互动数据，建立一个可校准的排序先验。\\
\textbf{不做什么}：不读取平台私有算法、不刷量、不控制账号、不绕过风控。\\
\textbf{所有平台权重都标记为 \sym{prior\_only}}，即“未经真实数据校准的猜测值”。
\end{explainbox}

\subsection{内容方案：把一条视频变成一组数字}

系统把一条内容抽象成一组 0 到 1 之间的分数：

\begin{center}
\begin{tabular}{p{4.4cm}p{9.4cm}}
\toprule
\textbf{字段} & \textbf{含义} \\
\midrule
\sym{hook\_strength} & 开头吸引力（前两秒能不能抓住人） \\
\sym{cognitive\_load} & 认知负荷（理解起来累不累） \\
\sym{novelty} & 新颖度（跟别人像不像） \\
\sym{authority} & 权威感（说话的人可信不可信） \\
\sym{social\_proof} & 社会证明（有没有案例、别人怎么说） \\
\sym{search\_intent} & 搜索意图（标题有没有覆盖真实搜索词） \\
\sym{emotional\_intensity} & 情绪强度 \\
\sym{completion\_prob} & 完播概率 \\
\sym{interaction\_prob} & 互动概率 \\
\sym{share\_prob} & 转发概率 \\
\sym{follow\_prob} & 关注概率 \\
\sym{content\_quality} & 内容质量 \\
\sym{compliance\_risk} & 合规风险（越高越差） \\
\sym{negative\_feedback} & 负反馈风险 \\
\bottomrule
\end{tabular}
\end{center}

\subsection{排序评分公式}

\begin{equation}
\text{得分} = \sigma\Big( \big(\text{调整后正分} - 0.42\big) \times 8 \Big)
\label{p1:eq:platform-score}
\end{equation}

其中“调整后正分”为：

\begin{equation}
\text{调整后正分} = \text{原始正分} \times \Big(0.75 + 0.25 \times \text{群体修正}\Big) \times e^{-\text{惩罚项}}
\label{p1:eq:adjusted}
\end{equation}

\begin{explainbox}{逐符号解释}
\begin{longtable}{@{}p{3.6cm}p{10.2cm}@{}}
$\sigma(\cdot)$ & sigmoid 函数（见 2.4 节），把任意实数压到 $(0,1)$。 \\
$0.42$ & 中心偏移量。让“没有优势的内容”得分低于 0.5。 \\
$8$ & 陡峭度。值越大，得分曲线越陡（区分度越高）。 \\
原始正分 & 各正向信号乘以各自权重后的总和。 \\
群体修正 & 由观众群体的对齐度、多样性等算出的修正系数。 \\
$e^{-\text{惩罚项}}$ & 指数衰减。惩罚项越大，整体得分被压得越低。 \\
\end{longtable}
\end{explainbox}

\textbf{通俗理解}：这个公式做了三件事——
\begin{enumerate}[leftmargin=1.6em]
\item 把各个正向信号加权加起来（原始正分）；
\item 乘上一个群体修正（观众群体配合得好不好）；
\item 乘上一个衰减项（有风险就扣分）。
最后用 sigmoid 压到 0 到 1，方便比较。
\end{enumerate}

\subsection{群体状态}

系统还跟踪观众群体的六个指标：

\begin{center}
\begin{tabular}{p{4.6cm}p{9.2cm}}
\toprule
\textbf{指标} & \textbf{含义} \\
\midrule
\sym{alignment} & 对齐度：大家的看法一致不一致 \\
\sym{diversity} & 多样性：观点是不是单一 \\
\sym{reciprocity} & 互惠度：有没有来回互动 \\
\sym{network\_density} & 网络密度：人群之间连得紧不紧 \\
\sym{common\_ground} & 共同基础：有没有共享背景 \\
\sym{mismatch} & 错位度：有没有理解偏差 \\
\bottomrule
\end{tabular}
\end{center}

\begin{explainbox}{为什么要跟踪多样性}
如果内容只吸引一种人，短期数据好看，长期会把群体压成“回音室”——
大家只能说同样的话。多样性指标就是为了防止这件事。
\end{explainbox}

\subsection{操作过程：从 PDF 到内容库}

具体做了什么：

\begin{enumerate}[leftmargin=1.8em]
\item \textbf{读入原始文案}：一份 22 页的 PDF，里面有 58 条口播文案。
\item \textbf{解析}：把每条文案拆成“标题 + 正文”，并识别它属于哪个栏目
      （K12 学习方法、家庭教育、音乐、体育、书法、编程等）。
\item \textbf{规则标注}：用关键词规则给每条打标签：
  \begin{itemize}
  \item 结尾动作（CTA）：是在引导评论、引导领取资料，还是引导转发；
  \item 钩子类型：是问句、是反常识，还是负面警示；
  \item 情绪驱动：焦虑、损失厌恶、亲情内疚等；
  \item 合规风险：有没有绝对化承诺、有没有制造教育焦虑。
  \end{itemize}
\item \textbf{映射成内容方案}：把标注结果换算成 4.2 节的那组 0–1 分数。
\item \textbf{分平台排序}：用不同平台的权重算出得分，排出名次。
\item \textbf{生成排期}：按“平台轮换 + 栏目去重”的规则排出 30 天日历。
\end{enumerate}

\begin{explainbox}{规则标注是什么}
“规则”就是“如果出现某个词，就打某个标签”。例如：\\
如果结尾出现“评论区”，就标记为引导互动；\\
如果出现“一定”、“必须”、“100\%”，就加合规风险分。\\
这样做的好处是\textbf{完全可解释}——每一条标签都能说出为什么。\\
缺点是\textbf{不理解语义}，只能靠词面。
\end{explainbox}

\subsection{双口径对比：两套权重谁对}

实验里做了一件很重要的事：同一批文案，用\textbf{两套权重}各排一次名。

\begin{center}
\begin{tabular}{p{3.2cm}p{10.6cm}}
\toprule
\textbf{口径} & \textbf{来源} \\
\midrule
实验口径 & 本实验自己设定的一套权重，强调完播和互动 \\
文档口径 & 一份招生素材里声称的平台权重，例如“收藏率 35\%、复访率 25\%” \\
\bottomrule
\end{tabular}
\end{center}

\textbf{为什么要对比}：因为在没有真实数据的情况下，\hl{两套权重谁对是不知道的}。
对比的价值在于找出“排名变化最大的那些选题”——那些就是最需要做实验验证的。

结果示例（抖音）：

\begin{center}
\begin{tabular}{p{1.6cm}p{5.2cm}p{2.0cm}p{2.0cm}p{1.8cm}}
\toprule
\textbf{编号} & \textbf{选题} & \textbf{实验排名} & \textbf{文档排名} & \textbf{变化} \\
\midrule
WX-011 & 孩子“一眨眼”就长大了 & 4 & 27 & $-23$ \\
WX-003 & 数学考 90 分以上的五个关键习惯 & 22 & 2 & $+20$ \\
WX-004 & 小学语文预习五步法 & 24 & 4 & $+20$ \\
\bottomrule
\end{tabular}
\end{center}

\begin{explainbox}{怎么读这张表}
“变化 = 实验排名 $-$ 文档排名”。负数表示在两套口径里都排在前面（稳定）；变化大的表示\hl{口径敏感}。\\
比如《数学考 90 分以上的五个关键习惯》，在实验口径里排第 22，
在文档口径里排第 2——差了 20 名。说明这条内容的好坏高度依赖“收藏重要还是完播重要”。
\end{explainbox}

\subsection{招生漏斗}

把“看到内容”到“报名”分成六级：

\begin{equation}
\text{曝光} \to \text{有效停留} \to \text{互动} \to \text{留资} \to \text{咨询} \to \text{报名}
\end{equation}

每一级的转化率用乘法串联：

\begin{equation}
\text{报名数} = \text{曝光} \times r_1 \times r_2 \times r_3 \times r_4 \times r_5
\end{equation}

\begin{explainbox}{逐符号解释}
\begin{longtable}{@{}p{2.4cm}p{11.4cm}@{}}
$r_1$ & 曝光 $\to$ 有效停留 的比例（先验值 0.30） \\
$r_2$ & 有效停留 $\to$ 互动 的比例 \\
$r_3$ & 互动 $\to$ 留资 的比例（先验值 0.06） \\
$r_4$ & 留资 $\to$ 咨询 的比例（先验值 0.45） \\
$r_5$ & 咨询 $\to$ 报名 的比例（先验值 0.18） \\
\end{longtable}
\end{explainbox}

\begin{explainbox}{这些数字从哪来}
有的来自那份招生素材里写的行业基准（例如“5 秒完播率 30\% 是基础标准”），
有的是直接假设的。\warn{全部标记为“未校准先验”，不能当预测用。}
\end{explainbox}

\textbf{一句重要的话}：一个内容的“曝光到报名”转化率通常只有\textbf{万分之几}。
所以不要指望一条视频带来很多报名——它是长期积累的事。

\newpage
% ============================================================

\section{用量化交易的眼光看平台算法}
% ============================================================

\subsection{一个自然的类比}

量化交易（用程序买卖股票）和做内容，在结构上很像：

\begin{center}
\begin{tabular}{p{4.2cm}p{9.6cm}}
\toprule
\textbf{量化交易} & \textbf{做内容} \\
\midrule
市场 & 平台的流量池 \\
标的（股票） & 一条内容 \\
价格 & 互动指标（完播、点赞、收藏） \\
成交量 & 曝光量 \\
撮合引擎 & 推荐系统 \\
交易员 & 创作者 \\
杠杆 & 投流（花钱买曝光） \\
回撤 & 流量下滑 \\
制度变化 & 平台改算法 \\
\bottomrule
\end{tabular}
\end{center}

\subsection{但有五个致命差异}

如果忽略下面这五条，量化框架就会被误用。

\begin{longtable}{@{}p{0.8cm}p{3.4cm}p{9.6cm}@{}}
\toprule
\textbf{序号} & \textbf{差异} & \textbf{说明} \\
\midrule
1 & 没有价格 & 你拿到的“曝光量”是平台单方面给的，不是所有参与者同意的数字 \\
2 & 看不到订单簿 & 你只能看到自己这一条内容，看不到同赛道有多少人在抢观众 \\
3 & 平台是裁判又是对手盘 & 它希望你一直生产，但不保证给你回报 \\
4 & 样本量差 6 个数量级 & 股票 5 年有 50 万根 K 线；一个创作者一年只有几百条内容 \\
5 & 不可做空、不可对冲 & 你只能“做多”，数据差时唯一的选择是不发 \\
\bottomrule
\end{longtable}

\begin{explainbox}{第 2 条为什么最要命}
在股票市场，你能看到\textbf{订单簿}——知道有多少人想买、多少人想卖。\\
在平台上，你\textbf{看不到同赛道有多少人在发同类内容}。\\
所以“这条内容是不是太拥挤了”这个判断，你\textbf{根本没有数据}。
\end{explainbox}

\subsection{民间算法 = 市场传言}

“民间算法”指创作者圈子里流传的平台规则解读（比如各种“秘籍”“白皮书”）。

\textbf{在量化里，它对应的是“技术分析流派”和“市场传言”。}

用因子的标准去筛：

\begin{center}
\begin{tabular}{p{4.6cm}p{2.2cm}p{2.0cm}p{4.4cm}}
\toprule
\textbf{说法} & \textbf{可观测吗} & \textbf{容量} & \textbf{判定} \\
\midrule
“收藏率权重 35\%、复访 25\%” & 否（拿不到权重） & — & \warn{不可检验} \\
“5 秒完播率 30\% 是基础线” & 是 & 大 & \hl{可检验} \\
“7 天长效赛马” & 是 & 大 & \hl{可检验} \\
“前 3 秒放钩子” & 是 & \warn{小} & \warn{会衰减} \\
\bottomrule
\end{tabular}
\end{center}

\begin{explainbox}{最重要的一个概念：因子衰减}
在量化里，一个因子一旦被广泛使用，超额收益会被套利掉。\\
在平台上，\textbf{民间算法的传播本身就在加速这个衰减过程}：

\begin{quote}
有人发现“前 3 秒放钩子”有效
$\to$ 写进白皮书、做成课程卖钱
$\to$ 所有人都在用
$\to$ 观众对钩子脱敏
$\to$ \textbf{钩子失效}
\end{quote}

\textbf{推论}：公开流传的平台“秘籍”，价值已经在传播过程中被消耗掉了。
\end{explainbox}

\newpage

% ============================================================

\section{群体拓扑：拆闭环、留桥接}
% ============================================================

\subsection{把班级看成一张图}

把班级互动看成图 $G=(V,E)$：

\begin{itemize}[leftmargin=2em]
\item $V$：学生节点；
\item $E$：互动边（提问、互评、讨论、求助）；
\item $C_i$：节点 $i$ 的局部聚类系数，衡量“我的朋友是否也互相认识”；
\item $Q$：模块度，衡量班级是否分裂成若干封闭社区；
\item $B$：桥接比例，跨社区边数 $/$ 总边数；
\item $D_{\text{div}}$：群体多样性（不同观点、强项、节奏的接触面）。
\end{itemize}

设计目标可以写成：

\begin{equation}
\max J=
\alpha B
-\beta \overline{C}
-\gamma Q
+\delta D_{\text{div}}
+\eta \mathrm{Var}(\text{individual growth}).
\end{equation}

符号说明：

\begin{itemize}[leftmargin=2em]
\item $\overline{C}$：平均聚类系数；越高，越容易形成封闭小圈子；
\item $Q$：模块度；越高，班级越分裂成抱团社区；
\item $B$：桥接比例；越高，跨圈信息流越强；
\item $D_{\text{div}}$：多样性；
\item $\mathrm{Var}(\text{individual growth})$：个体成长差异，防止“为了群体和谐牺牲个体发展”；
\item $\alpha,\beta,\gamma,\delta,\eta$：权重，需要按教学目标标定。
\end{itemize}

\subsection{E4 真实结果：小圈子 vs 桥接}

数据：Stanford MOOCPosts。

\begin{itemize}[leftmargin=2em]
\item 29{,}585 条帖子；
\item 11{,}042 名作者；
\item 62{,}857 条共帖边；
\item 有网络指标的作者 7{,}148。
\end{itemize}

\subsubsection{帖子触达（mean reads）模型}

控制 degree 和课程固定效应：

\begin{center}
\begin{longtable}{@{}lrr@{}}
\toprule
\textbf{预测变量} & \textbf{标准化 $\beta$} & \textbf{$p$} \\
\midrule
桥接比例 $B_i$ & $-0.096$ & $8.7\times10^{-58}$ \\
聚类系数 $C_i$ & $+0.188$ & $1.6\times10^{-98}$ \\
$\log(\text{degree})$ & $+0.661$ & $\approx0$ \\
\bottomrule
\end{longtable}
\end{center}

\subsubsection{帖子点赞（mean up）模型}

\begin{center}
\begin{longtable}{@{}lrr@{}}
\toprule
\textbf{预测变量} & \textbf{标准化 $\beta$} & \textbf{$p$} \\
\midrule
桥接比例 $B_i$ & $+0.043$ & $0.0008$ \\
聚类系数 $C_i$ & $-0.012$ & $0.400$ \\
$\log(\text{degree})$ & $+0.038$ & $0.0059$ \\
\bottomrule
\end{longtable}
\end{center}

\begin{explainbox}{关键解释}
结果不是简单的“抱团坏、桥接好”：\\
对\textbf{触达（reads）}，小圈子作者更高；\\
对\textbf{认可（upvotes）}，桥接作者略高。\\
所以“拆闭环”要按目标选：要传播，还是要质量反馈。
\end{explainbox}

\subsection{边界}

\begin{itemize}[leftmargin=2em]
\item 共帖网络不是真实社会关系；
\item reads/up 不是学习结果；
\item 这是观察性网络分析，不能直接说拆圈子导致成长；
\item 19 人班级网络只有 $19\times18/2=171$ 条可能边，只能当探索性试点。
\end{itemize}

\subsection{课堂 ABAB 设计}

\begin{center}
\begin{longtable}{@{}p{2.0cm}p{4.0cm}p{7.0cm}@{}}
\toprule
\textbf{周} & \textbf{分组} & \textbf{任务} \\
\midrule
W1--W2 & 固定分组 & 基线 \\
W3--W4 & 每周轮换分组 & 加跨组任务和桥接角色 \\
W5--W6 & 再固定 & 看是否回到闭环 \\
每周五 & 互动调查 & 写本周问过/讨论过的 3 个同学，其中至少 1 个跨组 \\
\bottomrule
\end{longtable}
\end{center}

\subsection{“分而治之”的教育翻译}

\begin{enumerate}[leftmargin=2em]
\item 定期轮换分组，打断固定小团体的路径依赖；
\item 设置跨组任务和桥接角色；
\item 小组任务必须能拆出个人贡献，避免搭便车；
\item 每个学生有自己的成长曲线，不拿圈子平均水平替代个体；
\item 保留正当合作，只禁止\textbf{封闭合作}。
\end{enumerate}

\newpage

% ============================================================

\section{E1--E6 实验总表与运行方法}
% ============================================================

\subsection{总表}

\begin{center}
\begin{longtable}{@{}p{1.7cm}p{2.6cm}p{3.1cm}p{4.6cm}p{2.4cm}@{}}
\toprule
\textbf{实验} & \textbf{状态} & \textbf{数据} & \textbf{最重要结果} & \textbf{产物目录} \\
\midrule
E1 教师--学生--观察者 & 设计+合成+功效；课堂数据待录入 & 实验 B 教师 AUC 先验 & AUC $=0.5625$，折后可信度 $\approx0.409$；MAE 优势 $0.08$ 功效 $0.891$ & E1\_observer \\
E2 节律--安排匹配 & 真实公开数据完成 & EdNet：340 人、16{,}335 会话 & 参与度 $\beta=0.164$（$p=8.3\times10^{-10}$）；退出率 $\beta=-0.0034$；$d_z=0.572$ & E2\_arrangement\_fit \\
E3 习惯改变成本 & 真实公开数据完成 & OULAD：3{,}829 单元、84{,}095 学生--周 & 参与度 $\beta=-1.134$；控制当前点击后 $\beta=-0.431$；退课 $\beta=+0.233$ & E3\_habit\_change \\
E4 群体拓扑 & 真实公开数据完成 & Stanford MOOCPosts：11{,}042 作者、62{,}857 边 & 触达：聚类 $+0.188$、桥接 $-0.096$；点赞：桥接 $+0.043$ & E4\_group\_topology \\
E5 多教师 metameta & 合成功效完成；真实多教师待采集 & 模拟 5--50 位教师 & 5 位教师、$sd_\tau=0.3$ 时 ANOVA 功效 $0.903$ & E5\_multi\_teacher \\
E6 天性/安排先验 & 合成+OULAD 先验完成；真实量表待采集 & 19 人 $\times$ 10 次模拟；OULAD 效应量 & MDES $\approx0.35$；中位效应 $0.053$ 需约 2{,}789 人，最大 $0.289$ 需约 92 人 & E6\_trait\_prior \\
\bottomrule
\end{longtable}
\end{center}

\subsection{运行方法}

在 PowerShell 里，先设好 Python 和 PYTHONPATH：

\begin{verbatim}
$b = "D:\Tender"
$py = "$b\Tender-Demo\backend\.venv\Scripts\python.exe"
$env:PYTHONPATH = "$b\experiments\platform_reverse_lab"
\end{verbatim}

然后逐个运行：

\begin{verbatim}
& $py run_teacher_student_observer.py
& $py run_arrangement_fit.py
& $py run_habit_change_cost.py
& $py run_group_topology.py
& $py run_multi_teacher_trait.py
\end{verbatim}

产物都在：

\begin{verbatim}
D:\Tender\experiments\platform_reverse_lab\outputs\related_experiments\
\end{verbatim}

\subsection{每个实验最少需要什么数据}

\begin{center}
\begin{longtable}{@{}p{2.2cm}p{4.6cm}p{6.8cm}@{}}
\toprule
\textbf{实验} & \textbf{最少数据} & \textbf{注意} \\
\midrule
E1 & 19 学生 $\times$ 5 次：学生预测、教师预测、教师信心、实际 & 三方不要互相看 \\
E2 & 每人 1--2 周行为日志 + 节律峰型 & 公开 EdNet 可先做先验 \\
E3 & 每人 10 周学习节奏 & OULAD 可先做先验 \\
E4 & 6 周互动调查，每人每周写 3 个联系人 & 19 人只能当试点 \\
E5 & 至少 5 位教师，每位 19 学生 $\times$ 5 次 & 每位教师 95 个观测 \\
E6 & MBTI 连续分数 + 节律峰型 + 任务结构偏好 & 不要填 16 类型 \\
\bottomrule
\end{longtable}
\end{center}

\newpage

% ============================================================


\part{Agent 交互层与 GitHub 教育 Agent 集成}

这一部分把产品补成完整的三层技术栈：

\begin{center}
\begin{longtable}{@{}p{2.4cm}p{4.0cm}p{7.4cm}@{}}
\toprule
\textbf{层} & \textbf{角色} & \textbf{回答什么} \\
\midrule
DAOQL 层（底层） & 事实与因果 & 发生了什么？因为什么？证据是什么？ \\
Tender 层（中层） & 心智化建模 & 学生现在是什么状态？不确定性多大？ \\
Agent 层（上层） & 交互与执行 & 现在该对谁说什么、做什么、怎么跟进？ \\
\bottomrule
\end{longtable}
\end{center}

\section{为什么需要 Agent 层}

DAOQL 负责“记事实”，Tender 负责“懂状态”，但两者都不会主动做事。
Agent 层负责把 DAOQL 的事实和 Tender 的状态，翻译成对老师、学生、家长、运营可执行的动作。

\begin{explainbox}{一句话}
\textbf{DAOQL 是因果事实层，Tender 是心智层，Agent 是交互层。}\\
没有 Agent，前两层只是数据仓库和模型；有了 Agent，才成为产品。
\end{explainbox}

\subsection{三层的数据契约}

\subsubsection{DAOQL 事件}

DAOQL 只存\textbf{发生过的、可追溯的事实}，不存模型的猜测。

\begin{verbatim}
{
  "event_id": "evt_...",
  "event_type": "metacognition_jol",
  "student_id": "S01",
  "session_id": "ses_...",
  "timestamp": "2026-10-16T20:00:00+08:00",
  "payload": {"predicted_correct": 7, "n_items": 10},
  "source": "browser_plugin|teacher_card|platform|agent",
  "provenance": {"device": "...", "course": "...", "item_ids": [...]}
}
\end{verbatim}

\subsubsection{Tender 状态}

Tender 把 DAOQL 事件合成为带不确定性的学生/教师/群组状态。

\begin{verbatim}
{
  "student_id": "S01",
  "timestamp": "2026-10-16T20:05:00+08:00",
  "cognitive": {"theta": 0.72, "phi": 1.20, "uncertainty": 0.08},
  "emotion": {"valence": 0.30, "arousal": 0.60},
  "metacognition": {"jol_bias": 0.08, "auc": 0.63, "hesitation": 0.42},
  "baseline": {"arrangement_fit": 0.81, "habit_change": 0.12},
  "knowledge": {"concept_id": "grammar.tense", "mastery": 0.67},
  "group": {"clustering": 0.55, "bridge_ratio": 0.21}
}
\end{verbatim}

\subsubsection{Agent 消息}

Agent 之间的通信也走统一消息格式，保证可审计、可回放。

\begin{verbatim}
{
  "message_id": "msg_...",
  "from_agent": "metacognition_coach",
  "to_agent": "student_tutor",
  "student_id": "S01",
  "intent": "scaffold_hesitation",
  "evidence": {"daoql_event_ids": [...], "tender_state_version": "..."},
  "action": {"type": "prompt", "text": "你卡在哪一步？"},
  "policy": {"requires_teacher_approval": false, "max_turns": 3},
  "status": "proposed|approved|executed|rejected"
}
\end{verbatim}

\section{Agent 角色清单}

\begin{center}
\begin{longtable}{@{}p{2.8cm}p{3.6cm}p{4.6cm}p{3.0cm}@{}}
\toprule
\textbf{Agent} & \textbf{对谁} & \textbf{主要职责} & \textbf{依赖层} \\
\midrule
Student Tutor Agent & 学生 & 苏格拉底式提问、分步提示、变式练习 & Tender + DAOQL + RAG \\
Metacognition Coach & 学生 & JOL、信心校准、犹豫四象限干预 & Tender 元认知 + DAOQL \\
Assessment Agent & 学生/老师 & 题目生成、难度标定、知识追踪更新 & Tender 知识追踪 + DAOQL \\
Teacher Copilot & 教师 & 班级状态摘要、干预建议、教师 metameta 反馈 & Tender + DAOQL \\
Parent Communication Agent & 家长 & 周报、成长曲线、隐私过滤后的沟通 & Tender 聚合状态 \\
Group Facilitator Agent & 教师/学生 & 分组轮换、桥接角色、闭环预警 & Tender 群体拓扑 + DAOQL \\
Content Curator Agent & 教研 & 课标/讲义/题库入库、RAG 切片、版本管理 & DAOQL + RAG \\
Sales/Operations Agent & 机构 & 线索评分、 onboarding、续费信号 & 聚合 DAOQL + 运营数据 \\
Safety/Guardrail Agent & 全部 & 内容安全、年龄适配、隐私、升级人工 & 全局策略 \\
Data Steward Agent & 工程/合规 & 数据质量、血缘、同意、保留与删除 & DAOQL + 治理规则 \\
\bottomrule
\end{longtable}
\end{center}

\section{GitHub 教育 Agent 与平台选型}

\subsection{教育专用层}

\begin{center}
\begin{longtable}{@{}p{3.2cm}p{1.4cm}p{2.0cm}p{4.8cm}p{3.0cm}@{}}
\toprule
\textbf{仓库} & \textbf{Stars} & \textbf{License} & \textbf{能提供什么} & \textbf{集成方式} \\
\midrule
ECNU-ICALK/EduChat & 973 & 无 License & 中文教育对话大模型、数据清理、教育语料 & 参考/替换为有 License 的模型 \\
ECNU-ICALK/EduChat-Math & 59 & 无 License & 中文多模态数学数据集与评测 & 作为评测集参考 \\
classroomio/classroomio & 1{,}719 & AGPL-3.0 & 课程、作业、班级管理 LMS & 外部 LMS，通过 API 接 \\
moodle/moodle & 7{,}488 & GPL-3.0 & 全球最成熟开源 LMS & 外部 LMS，通过插件/API 接 \\
openedx/openedx-platform & 8{,}203 & AGPL-3.0 & 大规模在线课程平台 & 外部 LMS，通过 API 接 \\
pykt-team/pykt-toolkit & 450 & MIT & 深度知识追踪基准与模型库 & 直接用于 Tender 知识追踪 \\
bigdata-ustc/EduKTM & 266 & Apache-2.0 & 知识追踪模型动物园 & 直接用于 Tender 知识追踪 \\
\bottomrule
\end{longtable}
\end{center}

\subsection{Agent/RAG 编排层}

\begin{center}
\begin{longtable}{@{}p{3.2cm}p{1.6cm}p{2.2cm}p{4.4cm}p{3.0cm}@{}}
\toprule
\textbf{仓库} & \textbf{Stars} & \textbf{License} & \textbf{能提供什么} & \textbf{集成方式} \\
\midrule
langchain-ai/langgraph & 43{,}347 & MIT & 有状态 Agent 工作流编排 & 核心编排引擎 \\
crewAIInc/crewAI & 59{,}759 & MIT & 角色化多 Agent 协作 & 教研/运营多 Agent \\
microsoft/autogen & 61{,}440 & CC-BY-4.0 & 多 Agent 研究与实验 & 研究实验，注意 License \\
langgenius/dify & 158{,}627 & Other & 低代码 Agent + RAG + 工具 & 快速原型/运营后台 \\
labring/FastGPT & 29{,}840 & Other & 知识库问答 + 工作流 & 国内部署候选 \\
infiniflow/ragflow & 92{,}317 & Apache-2.0 & 文档解析 + 深度 RAG & 教材/讲义/题库 RAG \\
open-webui/open-webui & 155{,}868 & Other & 聊天界面、多模型接入 & 学生/教师前端 \\
n8n-io/n8n & 207{,}606 & Fair-code & 400+ 集成与自动化 & CRM/通知/运营自动化 \\
\bottomrule
\end{longtable}
\end{center}

\begin{explainbox}{License 风险必须先说}
\textbf{可直接商用：}RAGFlow（Apache-2.0）、LangGraph（MIT）、CrewAI（MIT）、pyKT（MIT）、EduKTM（Apache-2.0）。\\
\textbf{Copyleft：}Moodle（GPL-3.0）、Open edX（AGPL-3.0）、ClassroomIO（AGPL-3.0）——建议当外部 LMS 通过 API 接，不要直接改代码闭源分发。\\
\textbf{需要法务确认：}EduChat（无 License）、Dify/FastGPT/Open WebUI（Other）、n8n（fair-code）、AutoGen（CC-BY-4.0）。\\
\textbf{原则：核心链路优先用 MIT/Apache；无 License 的模型只能参考，不能直接商用。}
\end{explainbox}

\section{集成架构}

\subsection{三层部署图}

\begin{center}
\begin{longtable}{@{}c@{}}
\toprule
\textbf{Agent 交互层} \\
Student Tutor · Metacognition Coach · Assessment · Teacher Copilot · Parent · Group Facilitator · Sales \\
$\downarrow$ \\
\textbf{Tender 心智层} \\
$\theta$ · Bloch 球 · 情绪 · 元认知 · 知识追踪 · 节律 · 基线 · 群体拓扑 \\
$\downarrow$ \\
\textbf{DAOQL 因果层} \\
事件图 · 因果关系 · 证据 · 血缘 · 时间线 · 版本 \\
\bottomrule
\end{longtable}
\end{center}

\subsection{推荐技术栈}

\begin{center}
\begin{longtable}{@{}p{3.4cm}p{4.4cm}p{6.0cm}@{}}
\toprule
\textbf{能力} & \textbf{推荐组件} & \textbf{理由} \\
\midrule
Agent 编排 & LangGraph & MIT、状态机清晰、可审计 \\
多 Agent 协作 & CrewAI & MIT、角色化、适合教研/运营 \\
RAG & RAGFlow & Apache-2.0、文档解析强、国内可部署 \\
低代码后台 & Dify / FastGPT & 快速搭运营和教研 Agent，注意 License \\
聊天前端 & Open WebUI & 多模型、可自托管，注意 License \\
工作流自动化 & n8n & CRM/通知/运营集成，注意 fair-code \\
知识追踪 & pyKT / EduKTM & MIT / Apache-2.0，直接进 Tender \\
教育对话 & EduChat（参考） & 中文教育语料强，但无 License，需替换或获授权 \\
LMS & ClassroomIO / Moodle / Open edX & 外部系统，通过 API 接，避免 AGPL/GPL 污染 \\
\bottomrule
\end{longtable}
\end{center}

\section{五条核心工作流}

\subsection{学生辅导工作流}

\begin{enumerate}[leftmargin=2em]
\item DAOQL 收到答题事件、响应时间、修改次数、JOL；
\item Tender 更新 $\theta$、元认知偏差、犹豫 AUC、知识掌握度；
\item Metacognition Coach 判断“高犹豫/低犹豫 × 对错”；
\item Student Tutor 选择提示层级：最小提示 $\to$ 变式 $\to$ 讲解；
\item 学生行动再次写回 DAOQL，形成闭环。
\end{enumerate}

\subsection{教师副驾驶工作流}

\begin{enumerate}[leftmargin=2em]
\item Teacher Copilot 读取班级 Tender 状态；
\item 标出“需要今天干预的 3 个学生”；
\item 给出可执行建议，并附证据链接到 DAOQL 事件；
\item 教师执行后记录结果，作为教师 metameta 的训练数据；
\item 每周给教师一份“你的判断 vs 实际结果”的校准报告。
\end{enumerate}

\subsection{群体促进工作流}

\begin{enumerate}[leftmargin=2em]
\item Group Facilitator 读取互动调查，构建班级图；
\item 计算 $C_i,Q,B$，识别封闭小圈子；
\item 生成轮换分组和桥接角色建议；
\item 教师确认后执行，下一周重新测量；
\item 目标不是拆散合作，而是拆闭环、留桥接。
\end{enumerate}

\subsection{家长沟通工作流}

\begin{enumerate}[leftmargin=2em]
\item Parent Agent 只读取聚合后的 Tender 状态；
\item 输出“本周进步、下周建议、需要家庭配合的一件事”；
\item 禁止输出诊断标签、排名、原始情绪分数；
\item 家长反馈写回 DAOQL，作为环境变量 $e_t$。
\end{enumerate}

\subsection{销售/运营工作流}

\begin{enumerate}[leftmargin=2em]
\item Sales Agent 读取机构级聚合数据，不接触学生 PII；
\item 识别续费风险、招生线索、课程匹配；
\item n8n 自动触发 CRM、企微、短信、日程；
\item 所有动作留审计日志，遵守教育数据合规。
\end{enumerate}

\section{安全、隐私与治理}

\subsection{十条硬红线}

\begin{enumerate}[leftmargin=2em]
\item 不把原始学生 PII 交给第三方模型；
\item 不把诊断标签展示给学生或家长；
\item 不用情绪分数做惩罚或排名；
\item Agent 不得自行修改课程标准；
\item 所有自动干预必须可回滚、可审计；
\item 未成年人数据默认最小化采集；
\item 数据出境必须走白名单和脱敏；
\item 无 License 的模型只能参考，不得直接商用；
\item AGPL/GPL 组件当外部服务接，避免污染闭源代码；
\item 任何“因果”结论必须有 DAOQL 证据链。
\end{enumerate}

\subsection{Human-in-the-loop}

\begin{itemize}[leftmargin=2em]
\item 高风险动作（通知家长、调整分组、生成正式评语）必须教师确认；
\item 低风险动作（提示、练习推荐）可以自动执行，但要记录；
\item 所有 Agent 建议都带证据链接和置信度；
\item 教师可以一键否决，否决本身也是训练数据。
\end{itemize}

\section{产品完成路线图}

\subsection{MVP：8 周}

\begin{center}
\begin{longtable}{@{}p{1.6cm}p{4.6cm}p{7.4cm}@{}}
\toprule
\textbf{周} & \textbf{目标} & \textbf{交付} \\
\midrule
W1--W2 & DAOQL 事件模型 + Tender 状态服务 & 事件 schema、状态 schema、数据接入 \\
W3--W4 & Student Tutor + Metacognition Coach & RAGFlow + LangGraph + 合规 LLM \\
W5 & Teacher Copilot & 班级仪表盘、metameta 报告 \\
W6 & Group Facilitator + 基线耦合 & 互动调查、分组建议、节律匹配 \\
W7 & Safety + Data Steward & 红线策略、审计日志、数据血缘 \\
W8 & ClassroomIO/Moodle 接入 + 19 人试点 & 端到端闭环、课堂报告 \\
\bottomrule
\end{longtable}
\end{center}

\subsection{Phase 2：3 个月}

\begin{itemize}[leftmargin=2em]
\item 多教师 metameta（至少 5 位教师）；
\item Parent Communication Agent；
\item Sales/Operations Agent（对接万象学院招生白皮书）；
\item 群体拓扑 ABAB 实验；
\item pyKT/EduKTM 知识追踪上线；
\item 浏览器插件/移动端数据采集。
\end{itemize}

\subsection{Phase 3：6 个月}

\begin{itemize}[leftmargin=2em]
\item 多租户 SaaS；
\item 学生成长联盟链/数字供应链；
\item 平台逆向与推荐算法；
\item 教育 Agent 应用商店；
\item 合规认证与等保。
\end{itemize}

\section{Agent 交互层的完整工程实现}

这一节把前面的架构落成可运行工程。工程目录：\texttt{D:\textbackslash Tender\textbackslash agent\_layer}。

\subsection{DAOQL 因果层}

已实现：

\begin{itemize}[leftmargin=2em]
\item 事件表 \texttt{events}：事件编号、类型、学生、会话、时间、载荷、来源、血缘；
\item 因果边表 \texttt{causal\_edges}：源事件、目标事件、关系、置信度、证据；
\item 审计表 \texttt{audit\_log}：谁、做了什么、目标、载荷、时间；
\item \texttt{causal\_path}：在因果图上做 BFS，查找“原因 $\to$ 结果”路径；
\item \texttt{event\_graph}：用 NetworkX 输出学生事件图。
\end{itemize}

\subsection{Tender 心智层}

已实现并复用项目已有模型：

\begin{itemize}[leftmargin=2em]
\item 掌握度与 $\theta$；
\item JOL 偏差、犹豫程度；
\item 元认知敏感性 AUC、偏差、$d'$、L2 折扣（复用 \texttt{platform\_reverse\_lab.metacognition}）；
\item 节律 $R_{24}$、$R_{168}$、平衡 $B$（复用 \texttt{rhythm\_periodogram}）；
\item 基线偏离（复用 \texttt{rhythm\_baseline}）；
\item 教师 metameta（复用 \texttt{metameta}）；
\item 群体聚类系数与桥接比例。
\end{itemize}

\subsection{Agent 运行时}

\begin{itemize}[leftmargin=2em]
\item 10 个 Agent YAML 配置；
\item RAG：本地 TF-IDF 字符 $n$-gram，可选 RAGFlow；
\item LLM：OpenAI-compatible API，默认 mock，配置 \texttt{OPENAI\_API\_KEY} 后自动切换；
\item 护栏：敏感词、PII、年龄适配、人工升级；
\item 审计：每次 Agent 调用写入 \texttt{audit\_log}；
\item 报告：学生、教师、家长、机构四类报告。
\end{itemize}

\subsection{已接入的真实数据}

\begin{itemize}[leftmargin=2em]
\item 完整版手册：91{,}096 字，切成 115 个 RAG 块；
\item 语法填空题库：1{,}148 行 CSV 已写成 DAOQL 事件；
\item 19 名学生：种子演示生成 380 条事件；
\item 已有 Tender 实验模块：元认知、节律、基线、metameta 已接入。
\end{itemize}

\subsection{API 端点}

\begin{longtable}{@{}p{6.4cm}p{7.0cm}@{}}
\toprule
\textbf{端点} & \textbf{作用} \\
\midrule
\texttt{GET /health} & 健康检查 + 事件总数 \\
\texttt{GET /agents} & 列出 10 个 Agent \\
\texttt{POST /daoql/events} & 写入 DAOQL 事件 \\
\texttt{POST /daoql/edges} & 写入因果关系 \\
\texttt{GET /daoql/graph/\{student\_id\}} & 学生事件图 \\
\texttt{GET /daoql/causal\_path} & 因果路径查询 \\
\texttt{GET /tender/state/\{student\_id\}} & Tender 状态 \\
\texttt{GET /reports/student/\{id\}} & 学生报告 \\
\texttt{GET /reports/teacher/\{class\}} & 教师报告 \\
\texttt{GET /reports/parent/\{id\}} & 家长报告（隐私过滤） \\
\texttt{POST /agents/\{agent\_id\}/invoke} & 调用 Agent \\
\texttt{GET /audit} & 审计日志 \\
\texttt{POST /rag/search} & RAG 检索 \\
\texttt{POST /rag/ingest} & RAG 入库 \\
\bottomrule
\end{longtable}

\subsection{部署与验证}

\begin{verbatim}
cd D:\Tender\agent_layer
python run_demo.py
python tests/run_all.py
python scripts/ingest_manual.py
python scripts/ingest_grammar_bank.py
python scripts/run_pilot.py
uvicorn app.main:app --port 8080
\end{verbatim}

已验证：

\begin{itemize}[leftmargin=2em]
\item 4 个测试模块全部通过；
\item 端到端试点通过：19 名学生、380 条事件、Tender 状态、元认知教练、学生报告；
\item 手册 RAG 115 块、语法填空题库 1{,}148 行已入库；
\item LLM 默认 mock，配置真实 Key 后切换为 OpenAI-compatible 模型。
\end{itemize}

\subsection{还差什么}

\begin{enumerate}[leftmargin=2em]
\item 真实 LLM Key 与模型选型；
\item 真实 LMS API（ClassroomIO/Moodle/Open edX）；
\item 19 人课堂试点；
\item AGPL/GPL/Other License 法务确认；
\item 浏览器插件/移动端采集；
\item 等保与未成年人数据合规。
\end{enumerate}

\begin{explainbox}{完成度}
\textbf{架构、手册、DAOQL、Tender、Agent、RAG、报告、护栏、审计、测试、部署骨架全部完成。}\\
剩下的是“接真实账号、接真实学校、跑试点”，不是从零开始。
\end{explainbox}


\section{产品验收清单}

\begin{center}
\begin{longtable}{@{}p{3.4cm}p{9.8cm}@{}}
\toprule
\textbf{维度} & \textbf{验收标准} \\
\midrule
数据 & 每次事件可追溯到来源；同意、保留、导出、删除全流程可用 \\
模型 & Tender 状态有版本号和不确定性；离线校准达标 \\
Agent & 每个 Agent 有工具白名单、护栏、审计日志 \\
集成 & LMS、插件、教师卡、家长端、CRM 至少各打通一条链路 \\
评估 & 离线指标 + 在线 A/B + 安全评估 + 人工抽检 \\
商业 & 定价、打包、ROI、续费信号、机构隔离 \\
合规 & 十条红线全部通过；无 License 风险 \\
\bottomrule
\end{longtable}
\end{center}

\section{一句话总结}

\begin{quote}
\large
\textbf{DAOQL 记事实，Tender 懂心智，Agent 做执行。}\\
三层各司其职，才能把“教育数据”变成“教育产品”。
\end{quote}

\part{实验复现与课堂操作}

\section{实验二：知识追踪基准实验（重点）}
% ============================================================

这是本文篇幅最长、也最重要的一章。它回答一个很实际的问题：

\begin{quote}
\large \textbf{本系统的 $\theta$ 表示，和学术界成熟的方法比，到底行不行？}
\end{quote}

\subsection{什么是知识追踪}

\term{知识追踪}（Knowledge Tracing）要预测的是：

\begin{quote}
\textbf{给定一个学生过去做过的所有题，预测他下一题做对的概率。}
\end{quote}

举例：小明做了 5 道关于“一元一次方程”的题，对错分别是
“对、对、错、对、错”。知识追踪要回答：第 6 道同类题，他做对的概率有多大？

\begin{explainbox}{为什么这件事重要}
如果能准确预测“哪道题他做不对”，系统就能：
\begin{itemize}[leftmargin=1.4em,itemsep=1pt]
\item 只推他不会的题（不浪费时间）；
\item 判断他是否真的掌握了（而不是运气好）；
\item 给老师提示“这个知识点全班都没掌握”。
\end{itemize}
\end{explainbox}

\subsection{实验用的数据}

\begin{center}
\begin{tabular}{p{4.2cm}p{9.6cm}}
\toprule
\textbf{项目} & \textbf{数值} \\
\midrule
数据集名称 & ASSISTments 2009（美国中小学数学辅导系统的公开数据） \\
原始记录数 & 346,860 条 \\
原始学生数 & 4,217 人 \\
原始知识点数 & 101 个 \\
清洗后记录数 & \textbf{273,364 条} \\
清洗后学生数 & \textbf{4,128 人} \\
清洗后知识点数 & \textbf{81 个} \\
训练集 & 218,945 条 / 3,303 人 \\
测试集 & 54,419 条 / 825 人 \\
\bottomrule
\end{tabular}
\end{center}

\subsubsection{数据清洗规则}

\begin{enumerate}[leftmargin=1.8em]
\item \textbf{删掉没有知识点标签的记录}：101 个知识点里有 20 个样本太少，保留会造成误判。
\item \textbf{去重}：同一个人对同一道题的多条重复记录只保留一条。
\item \textbf{按时间和学生排序}：必须还原出“谁先做了什么、后做了什么”。
\end{enumerate}

\subsubsection{为什么要“按学生”划分训练集和测试集}

这一点很重要，是一个常见的错误来源：

\begin{center}
\begin{tabular}{p{5.0cm}p{8.8cm}}
\toprule
\textbf{划分方式} & \textbf{后果} \\
\midrule
按记录随机划分 & \warn{成绩虚高}。同一个学生的历史同时出现在训练和测试里，
模型等于“见过答案”。\\
\textbf{按学生划分} & 测试集里的学生，训练时\textbf{完全没见过}。这才反映真实使用场景。 \\
\bottomrule
\end{tabular}
\end{center}

本实验采用\hl{按学生划分}（80\% 学生训练，20\% 学生测试）。

\subsection{模型一：BKT（标准基线）}

BKT 全称 \term{贝叶斯知识追踪}（Bayesian Knowledge Tracing），是学术界用了三十多年的经典方法。

\subsubsection{核心思想}

对每一个知识点，用一个“是否已掌握”的开关来建模：

\begin{center}
\begin{tabular}{p{3.0cm}p{10.8cm}}
\toprule
\textbf{状态} & \textbf{含义} \\
\midrule
未掌握 & 该知识点还没学会 \\
已掌握 & 该知识点已经学会（而且不会忘记） \\
\bottomrule
\end{tabular}
\end{center}

\subsubsection{四个参数}

每个知识点有四个参数，全部是概率：

\begin{longtable}{@{}p{2.2cm}p{3.0cm}p{8.6cm}@{}}
\toprule
\textbf{符号} & \textbf{名称} & \textbf{含义} \\
\midrule
$P(L_0)$ & 先验掌握概率 & 开始做题之前，他已经掌握这个知识点的概率 \\
$P(T)$ & 学习转移概率 & 做一道题之后，从“未掌握”变成“已掌握”的概率 \\
$P(G)$ & 猜测概率 & \textbf{没掌握}却答对的概率（蒙对） \\
$P(S)$ & 失误概率 & \textbf{已掌握}却答错的概率（手滑） \\
\bottomrule
\end{longtable}

\begin{explainbox}{为什么要区分“猜对”和“真会”}
这是 BKT 的关键洞察。如果一个学生答对了：
\begin{itemize}[leftmargin=1.4em,itemsep=1pt]
\item 可能是真会（掌握状态 + 正常发挥）；
\item 也可能只是蒙对（未掌握状态 + 猜对）。
\end{itemize}
BKT 用概率的方式同时考虑这两种可能，所以比“答对就是会”要准得多。
\end{explainbox}

\subsubsection{预测公式}

\begin{equation}
P(\text{答对})_t = P(L_t)\big(1 - P(S)\big) + \big(1 - P(L_t)\big)P(G)
\label{p1:eq:bkt-predict}
\end{equation}

\begin{explainbox}{逐项解释}
\begin{itemize}[leftmargin=1.4em,itemsep=1pt]
\item $P(L_t)$：第 $t$ 步时“已掌握”的概率。
\item $P(L_t)(1-P(S))$：如果他已掌握，那么“正常答对”的概率。
\item $(1-P(L_t))P(G)$：如果他未掌握，那么“蒙对”的概率。
\item 两项相加 = 总的“答对”概率。这是\hl{全概率公式}的应用：把所有可能的情况加总。
\end{itemize}
\end{explainbox}

\subsubsection{状态更新（看到答案之后）}

答对之后，掌握概率要更新：

\begin{equation}
P(L_t \mid \text{答对}) = \frac{P(L_t)\big(1-P(S)\big)}{P(\text{答对})_t}
\label{p1:eq:bkt-update}
\end{equation}

看到答案之后，还要考虑“学了一道题会不会掌握”：

\begin{equation}
P(L_t^{+}) = P(L_t \mid \text{答对}) + \Big(1 - P(L_t \mid \text{答对})\Big)P(T)
\label{p1:eq:bkt-transition}
\end{equation}

\begin{explainbox}{$(1-P(S))$ 是什么}
$P(S)$ 是“已掌握却答错”的概率。那么 $1-P(S)$ 就是“已掌握且正常答对”的概率。
两个数加起来是 1，这是概率的基本要求。
\end{explainbox}

\begin{explainbox}{式 \eqref{p1:eq:bkt-update} 的直觉}
这是一个\hl{贝叶斯更新}。分子是“已掌握并且答对”的概率，
分母是“答对”的总概率。相除得到“在答对的条件下，他确实掌握了”的概率。\\
直观理解：答对之后，“已掌握”的可信度应该上升，但不会变成 100\%——因为可能只是蒙对。
\end{explainbox}

\subsubsection{怎么求出这四个参数：EM 算法}

四个参数不是人给的，是从数据里“学”出来的。用的是一个叫 \term{EM 算法} 的方法。

\begin{explainbox}{EM 算法用生活例子解释}
假设你想知道“一个班有多少人真正会做这道题”，但你不能直接问。
你只能看到“谁答对了”。\\
问题：答对的人里，哪些是真会，哪些是蒙的？\\
EM 算法的做法是“猜—修正”循环：
\begin{enumerate}[leftmargin=1.4em,itemsep=1pt]
\item \textbf{E 步（估计）}：先随便猜一组参数，然后算出“每个答对的人有多大概率是真会”；
\item \textbf{M 步（最大化）}：用上一步的结果，重新算一组更合理的参数；
\item 重复 1、2，直到参数不再明显变化。
\end{enumerate}
数学上可以证明，这个循环会让“数据出现的可能性”（似然）单调上升，最终收敛到一个较好的解。
\end{explainbox}

\begin{explainbox}{为什么要跑多次（多次重启）}
EM 算法有一个缺点：它可能停在“局部最优”——也就是在一个小山头上停住，
而没找到最高峰。\\
解决办法：\textbf{用不同的随机初始值跑很多次，取最好的那个}。\\
本实验对每个知识点跑了 3 次随机重启，每次最多迭代 50 轮。
\end{explainbox}

\subsection{模型二：Tender $\theta$\_seq}

这是本项目的模型。它的思路和 BKT 完全不同：

\begin{center}
\begin{tabular}{p{3.4cm}p{5.6cm}p{4.8cm}}
\toprule
\textbf{对比项} & \textbf{BKT} & \textbf{Tender $\theta$\_seq} \\
\midrule
状态是什么 & 一个概率（0 到 1） & 一个角度（0 到 $\pi$） \\
状态怎么变 & 学会就永久记住 & 可以上升也可以下降 \\
答对的影响 & 掌握概率上升 & 角度下降（更清醒） \\
答错的影响 & 掌握概率下降 & 角度上升（更糊涂） \\
\bottomrule
\end{tabular}
\end{center}

\subsubsection{用到的公式}

\hl{（1）状态递推}——公式 \eqref{p1:eq:theta-seq}：

\begin{equation}
\theta_t = \theta_{t-1} + \eta \cdot \text{趋势} + \xi \cdot \text{场景切换惩罚}
\end{equation}

\hl{（2）预测}——公式 \eqref{p1:eq:bloch-map}：

\begin{equation}
P(\text{答对}) = \cos^2\!\left(\frac{\theta}{2}\right)
\end{equation}

\hl{（3）默认参数}（来自项目文档）：

\begin{center}
\begin{tabular}{p{2.6cm}p{2.4cm}p{9.0cm}}
\toprule
\textbf{参数} & \textbf{文档默认值} & \textbf{含义} \\
\midrule
$\eta$ & 0.30 & 每步 $\theta$ 的更新幅度 \\
$\xi$ & 0.15 & 切换到新知识点时的额外惩罚 \\
$\theta$ 初始值 & 由 0.65 反解 & 假设该知识点初始正确率约 65\% \\
\bottomrule
\end{tabular}
\end{center}

\begin{explainbox}{“由 0.65 反解”是什么意思}
如果我们假设某个知识点的初始正确率是 65\%（即 0.65），
那就可以用公式 \eqref{p1:eq:bloch-map} 反过来求出对应的 $\theta$：\\
\begin{equation*}
0.65 = \cos^2\!\left(\frac{\theta}{2}\right)
\;\Longrightarrow\;
\theta = 2\arccos(\sqrt{0.65}) \approx 1.26
\end{equation*}
这一步叫“反解”，是为了让模型的出发点符合该知识点的实际难度。
\end{explainbox}

\subsubsection{本实验实现的三个变体}

为了弄清楚“是哪一部分在起作用”，实验做了三个变体：

\begin{enumerate}[leftmargin=1.8em]
\item \textbf{全局初始化}：所有知识点共用一个初始 $\theta$，搜索最优的 $\theta_0$、$\eta$、$\xi$。
\item \textbf{分技能初始化}：每个知识点的初始 $\theta$ 由该知识点的实际正确率反解，只搜索 $\eta$、$\xi$。
\item \textbf{文档默认参数}：完全采用文档给的值（$\eta=0.3$、$\xi=0.15$），不做任何调优。
\end{enumerate}

另外还做了一个\textbf{消融实验}：把“按知识点分状态”这个设计去掉（只维护一个 $\theta$），
看性能掉多少。消融实验的作用是验证“某个设计到底有没有用”。

\subsection{评测指标}

\subsubsection{指标一：AUC（看排序能力）}

\begin{explainbox}{AUC 用生活例子解释}
想象你把这批测试题分成两组：一组是学生实际答对的，一组是实际答错的。\\
现在\textbf{随机取一道答对的题和一道答错的题}，比较模型给它们打的“答对概率”：
\begin{itemize}[leftmargin=1.4em,itemsep=1pt]
\item 如果模型给答对的那道打了更高的分，算赢；
\item 如果给答错的那道打了更高的分，算输；
\item 平局算一半。
\end{itemize}
\textbf{AUC 就是“赢”的比例。}
\end{explainbox}

\begin{center}
\begin{tabular}{p{3.0cm}p{10.8cm}}
\toprule
\textbf{AUC 值} & \textbf{含义} \\
\midrule
0.50 & 等于瞎猜（模型没有区分能力） \\
0.60 & 有一点用 \\
0.70 & 不错 \\
0.80 & 很好 \\
1.00 & 完美（现实中几乎不可能） \\
\bottomrule
\end{tabular}
\end{center}

\textbf{注意}：AUC 只看“谁排在前面”，\hl{不看概率数值是否准确}。
一个模型可以说“这题 99\% 会做”而实际只有 60\%，只要排序对，AUC 仍然高。

\subsubsection{指标二：RMSE（看数值准不准）}

RMSE 全称“均方根误差”。公式：

\begin{equation}
\text{RMSE} = \sqrt{\frac{1}{N}\sum_{i=1}^{N}\big(y_i - \hat{p}_i\big)^2}
\label{p1:eq:rmse}
\end{equation}

\begin{explainbox}{逐符号解释}
\begin{longtable}{@{}p{2.2cm}p{11.6cm}@{}}
$N$ & 测试样本总数（本实验 54,419） \\
$y_i$ & 第 $i$ 条的实际结果。答对 = 1，答错 = 0 \\
$\hat{p}_i$ & 模型预测第 $i$ 条答对的概率 \\
$y_i - \hat{p}_i$ & 误差。比如实际答对（1）但模型只给 0.3，误差就是 0.7 \\
$\sum$ & 求和符号，把所有样本的误差平方加起来 \\
$\frac{1}{N}$ & 除以样本数，得到平均值 \\
$\sqrt{\;\cdot\;}$ & 开平方。这样单位回到原来的量纲 \\
\end{longtable}
\end{explainbox}

\textbf{通俗理解}：RMSE 就是“平均差多少”。值越小越好。

\subsubsection{指标三：LogLoss（看自信程度是否合理）}

LogLoss 是本文最重要的指标，因为它能抓出一种很危险的问题：\hl{过度自信}。

\begin{equation}
\text{LogLoss} = -\frac{1}{N}\sum_{i=1}^{N}\Big[y_i \log \hat{p}_i \;+\; (1-y_i)\log(1-\hat{p}_i)\Big]
\label{p1:eq:logloss}
\end{equation}

\begin{explainbox}{为什么这个公式能抓出“过度自信”}
关键在 $\log$（对数）的形状：$\log$ 在接近 0 的地方会急剧下降。

\begin{center}
\begin{tabular}{p{3.4cm}p{3.4cm}p{6.6cm}}
\toprule
\textbf{实际情况} & \textbf{模型预测} & \textbf{惩罚} \\
\midrule
答对 & 0.99 & 很小 \\
答对 & 0.60 & 中等 \\
答对 & 0.10 & 很大 \\
\textbf{答对} & \textbf{0.01} & \warn{极大} \\
\bottomrule
\end{tabular}
\end{center}

所以如果模型用很肯定的语气说“这题他肯定不会做”，结果他做对了，LogLoss 会飙升。
\end{explainbox}

\begin{explainbox}{数值参考}
以一个只有 66\% 正确率的总体为参照，如果模型只会说“永远是 66\%”，
LogLoss 大约 0.639。\\
\textbf{如果模型的 LogLoss 比 0.639 还差，说明它的预测比“什么都不说”更有害。}
本实验里就出现了这种情况——见 5.8 节。
\end{explainbox}

\subsection{实验流程}

\begin{enumerate}[leftmargin=1.8em]
\item 读入 346,860 条原始记录；
\item 清洗，得到 273,364 条有效记录；
\item 按学生分成训练集（3,303 人）和测试集（825 人）；
\item 在训练集上：
  \begin{itemize}
  \item 对 81 个知识点分别跑 BKT 的 EM 估计；
  \item 用网格搜索找出 Tender 的最优 $\eta$、$\xi$（以及初始 $\theta$）；
  \end{itemize}
\item 在测试集上预测每一题的答对概率；
\item 计算 AUC、RMSE、LogLoss；
\item 做 $\eta$ 敏感性扫描、分知识点对比、校准曲线。
\end{enumerate}

\begin{explainbox}{什么是“网格搜索”}
如果要找最优的 $\eta$，它可能是 0.01、0.02、0.05…… 任何一个值。
网格搜索的做法是：\textbf{列出一串候选值，一个个试，取最好的}。\\
比如本实验试了 $\eta \in \{0.01, 0.02, 0.05, 0.10, 0.15, 0.20, 0.30, 0.45\}$，
再和 $\xi \in \{0, 0.1, 0.25, 0.4\}$ 组合，总共试 32 种搭配。\\
这是最笨但在小参数空间上最可靠的方法。
\end{explainbox}

\subsection{实验结果}

\begin{longtable}{@{}p{6.0cm}rrr@{}}
\toprule
\textbf{模型} & \textbf{AUC $\uparrow$} & \textbf{RMSE $\downarrow$} & \textbf{LogLoss $\downarrow$} \\
\midrule
全体均值（下限参照） & 0.5000 & 0.4726 & 0.6388 \\
学生历史平均分 & 0.6797 & 0.4532 & 0.7436 \\
知识点平均分 & 0.6135 & 0.4634 & 0.6196 \\
\textbf{BKT（标准基线）} & \textbf{0.6794} & \textbf{0.4487} & \textbf{0.5910} \\
\midrule
Tender $\theta$（全局初始化） & \textbf{0.6853} & 0.4608 & 0.6164 \\
Tender $\theta$（分技能初始化） & 0.6661 & 0.4534 & 0.6010 \\
Tender $\theta$（\textbf{文档默认} $\eta=0.3$） & \textbf{0.6895} & 0.4740 & \warn{1.1155} \\
Tender $\theta$（单状态消融） & 0.6444 & 0.4864 & \warn{1.3291} \\
Tender $\theta$（分技能 + 校准） & 0.6661 & 0.4533 & 0.6005 \\
\midrule
\textbf{逻辑回归（控制组）} & \textbf{0.7058} & \textbf{0.4416} & \textbf{0.5752} \\
\bottomrule
\caption{模型对比。箭头表示“越大越好”或“越小越好”。}
\end{longtable}

\begin{explainbox}{“逻辑回归控制组”是什么，为什么要有它}
逻辑回归是一种最基础、最简单的统计模型。\\
实验给它喂了几个朴素的特征：学生历史平均分、知识点平均分、做题位置、在该知识点上的尝试次数。\\
\textbf{为什么要加它}：如果 Tender 连这么简单的方法都打不过，
那就说明“$\theta$ 这个表示没有带来额外价值”。\\
结果是：\warn{逻辑回归全面最好}（AUC 0.7058）。
\end{explainbox}

\subsubsection{怎么读这张表}

\begin{enumerate}[leftmargin=1.8em]
\item \textbf{按 AUC 看}：Tender（0.666–0.690）和 BKT（0.6794）\hl{基本打平}，
      其中“文档默认参数”那一档（0.6895）还略高。
\item \textbf{按 RMSE / LogLoss 看}：Tender \warn{明显不如 BKT}。
\item \textbf{和逻辑回归比}：Tender 和 BKT \warn{都不如}那个简单控制组。
\item \textbf{BKT 的 0.6794 落在文献常见区间（0.67–0.70）}，
      说明实现是可信的，不是因为代码写错才输。
\end{enumerate}

\subsection{关键发现一：默认参数 $\eta$ 太大了}

这是整个实验最有价值的发现。

\subsubsection{敏感性扫描结果}

固定 $\xi = 0$，让 $\eta$ 从小到大变化，观察两个指标：

\begin{center}
\begin{tabular}{rrrr}
\toprule
$\eta$ & \textbf{测试 LogLoss $\downarrow$} & \textbf{测试 AUC $\uparrow$} & \textbf{说明} \\
\midrule
0.01 & 0.626 & 0.6853 & \\
\textbf{0.02} & \textbf{0.616} & 0.6853 & \hl{校准最优点} \\
0.05 & 0.638 & 0.6854 & \\
0.10 & 0.722 & 0.6855 & \\
0.20 & 0.953 & 0.6885 & \\
\textbf{0.30} & \warn{1.148} & 0.6918 & \warn{文档默认值} \\
0.45 & 1.421 & \textbf{0.6957} & 排序最优点 \\
0.60 & 1.549 & 0.6948 & \\
0.80 & 1.759 & 0.6923 & \\
\bottomrule
\end{tabular}
\end{center}

\begin{explainbox}{这张表里最重要的一件事}
请注意这两个指标的方向是\textbf{相反}的：
\begin{itemize}[leftmargin=1.4em,itemsep=1pt]
\item 按 \textbf{LogLoss}（概率准不准），$\eta$ 越小越好，最优点在 \textbf{0.02}；
\item 按 \textbf{AUC}（排序对不对），$\eta$ 越大越好，最优点在 \textbf{0.45}。
\end{itemize}
\end{explainbox}

\subsubsection{为什么会这样}

\begin{explainbox}{用一句话解释这个矛盾}
$\eta$ 越大，$\theta$ 摆动越剧烈，预测概率越极端（越来越接近 0 或 1）。\\[3pt]
\textbf{极端预测不改变“谁排在前面”}——所以排序指标（AUC）反而略有改善；\\
\textbf{但极端预测会让概率严重偏离真实频率}——所以校准指标（LogLoss）崩掉。
\end{explainbox}

\textbf{一个具体的例子}：

\begin{center}
\begin{tabular}{p{2.4cm}p{3.4cm}p{3.4cm}p{5.0cm}}
\toprule
$\eta$ & 模型说的话 & 实际正确率 & 后果 \\
\midrule
0.02 & “这题 70\% 会做” & 70\% & 正确，老师可以信任 \\
0.30 & “这题 95\% 会做” & 70\% & \warn{误导，老师会误判} \\
\bottomrule
\end{tabular}
\end{center}

\subsubsection{这意味着什么}

\begin{explainbox}{结论}
\textbf{文档里的 $\eta = 0.3$ 不算错，但它优化的目标不是“概率校准”。}\\
\begin{itemize}[leftmargin=1.4em,itemsep=1pt]
\item 如果 Tender 只用于\textbf{排序}（比如决定今天先讲哪道题），$\eta=0.3$ 可以接受；
\item 如果 Tender 要\textbf{显示概率}给老师看，必须用 $\eta \approx 0.02$ 并做校准。
\end{itemize}
\hl{这个区别必须在产品里写清楚。}
\end{explainbox}

\subsection{关键发现二：场景切换惩罚没有用}

网格搜索的结果显示，$\xi$ 的最优值落在边界 \textbf{0}。
也就是说，“学生切换到新知识点时加一点不确定性”这条设定，
在这个数据上\warn{没有带来任何增益}。

\begin{explainbox}{可能的原因}
因为模型本来就是\textbf{按知识点分别维护 $\theta$} 的。
学生切换到新知识点时，这个知识点的 $\theta$ 从初始值开始——
\textbf{不确定性已经通过“初始化”体现了}，再加惩罚是重复计算。
\end{explainbox}

\subsection{关键发现三：按知识点分状态是必要的}

\begin{center}
\begin{tabular}{p{5.4cm}rr}
\toprule
\textbf{变体} & \textbf{AUC} & \textbf{LogLoss} \\
\midrule
按知识点分别维护 $\theta$ & 0.666–0.685 & 0.601–0.616 \\
\textbf{只维护一个 $\theta$（消融）} & \warn{0.6444} & \warn{1.3291} \\
\bottomrule
\end{tabular}
\end{center}

去掉技能维度之后，AUC 掉了 0.03–0.04，LogLoss 翻了一倍多。

\textbf{结论}：Tender“每个知识点各自维护 $\theta$”的设计\hl{是需要的，不是多余的复杂度}。

\subsection{分知识点表现与校准曲线}

\subsubsection{分知识点对比}

在 77 个可比知识点上，逐一对 BKT 和 Tender 算 AUC：

\begin{itemize}[leftmargin=1.6em]
\item Tender 更好的知识点：\textbf{42 个}
\item BKT 更好的知识点：\textbf{35 个}
\item 中位差：\textbf{+0.0050}（Tender 略占优）
\end{itemize}

\begin{explainbox}{怎么解读“互有胜负、中位持平”}
这说明\textbf{两者没有谁系统性碾压谁}。
不同的知识点，两种方法各有优势——
可能有的知识点适合“掌握/未掌握”的二元模型，
有的更适合“连续变化的角度”模型。
这种结果在研究中很常见，不是失败。
\end{explainbox}

\subsubsection{校准曲线}

校准曲线的读法：
\begin{itemize}[leftmargin=1.6em]
\item 横轴：模型说的答对概率；
\item 纵轴：实际答对频率；
\item 如果模型是可信的，这条线应该\hl{贴近 45 度对角线}。
\end{itemize}

实验结果：
\begin{itemize}[leftmargin=1.6em]
\item \textbf{BKT}：基本贴合对角线；
\item \textbf{Tender（$\eta=0.02$）}：也接近，两端略有偏离；
\item \textbf{Tender（$\eta=0.3$）}：\warn{严重偏离}。在预测 0.9 以上的样本里，
      实际正确率远低于 0.9。
\end{itemize}

\begin{explainbox}{这就是“过度自信”的样子}
把不确定的状态说成确定的。\\
在教学场景里，这种错误很危险：老师会以为学生掌握了，从而跳过必要的讲解。
\end{explainbox}

\subsection{本实验的局限（必须说明）}

\begin{enumerate}[leftmargin=1.8em]
\item \textbf{只测了一个来源}。完整 Tender 有四个来源（颜色、文本、序列、生理），
      本实验的数据只有“作答序列”，所以只测了 $\theta^{\text{seq}}$。
      \hl{这是下界，不是完整 Tender 的能力上限。}
\item \textbf{数据集不匹配}。ASSISTments 是\textbf{美国中学生的数学题}，
      不是中文语法填空。结论对数学知识点更可靠，对语言类要重新验证。
\item \textbf{没有和深度模型比}。像 DKT、SAKT 这类深度知识追踪方法，
      在这个数据集上 AUC 通常能到 0.72–0.75。本实验没有跑这些，
      所以\warn{不能宣称“$\theta$ 达到了前沿水平”}。
\item \textbf{只做了一次划分}。种子固定，没有做多种子重复实验。
      差异在 $\pm 0.01$ 以内的比较不应过度解读。
\end{enumerate}

\subsection{实验得出的四条修正}

\begin{longtable}{@{}p{1.0cm}p{6.4cm}p{6.2cm}@{}}
\toprule
\textbf{序号} & \textbf{改什么} & \textbf{依据} \\
\midrule
1 & 把 $\eta$ 从 0.30 改成 0.02 附近，并加概率校准 & LogLoss 从 1.15 降到 0.62 \\
2 & $\xi$ 默认改成 0，文档注明与分技能状态重复 & 网格搜索最优在边界 \\
3 & 保留“按知识点分状态” & 消融实验显著变差 \\
4 & 教师端不显示概率，改显示等级（未掌握 / 接近 / 已掌握） & 规避校准误差误导 \\
\bottomrule
\end{longtable}

\begin{explainbox}{第 4 条值得特别说明}
即使调整了 $\eta$ 并做了校准，模型输出的仍然是概率。
但老师看到“73\%”这样的数字，其实很难判断该怎么做。\\
\textbf{改成三个等级}之后：
\begin{itemize}[leftmargin=1.4em,itemsep=1pt]
\item 老师一眼就能判断；
\item 校准误差被“分级”吸收了，不会造成严重误导。
\end{itemize}
这在产品设计上叫“用合适的粒度呈现不确定性”。
\end{explainbox}

\newpage
% ============================================================

\section{附录：第一期复现指南}
% ============================================================

\subsection{平台逆向实验}

\begin{verbatim}
$env:PYTHONPATH='D:\Tender\experiments\platform_reverse_lab'
$py = 'D:\Tender\Tender-Demo\backend\.venv\Scripts\python.exe'

# 1) 重建内容库（从原始文案文本）
& $py "$base\tools\build_content_library.py"

# 2) 内容库排序 + 30 天日历 + 图表
& $py "$base\run_content_library.py"

# 3) 实验口径 vs 文档口径 排序对比
& $py -m platform_reverse_lab.dual_prior

# 4) 整合报告（含招生漏斗）
& $py "$base\run_integration.py"
\end{verbatim}

\subsection{知识追踪基准实验}

\begin{verbatim}
$env:PYTHONPATH='D:\Tender\experiments\knowledge_tracing_bench'
$py = 'D:\Tender\Tender-Demo\backend\.venv\Scripts\python.exe'

# 跑基准（约 5 分钟）
& $py D:\Tender\experiments\knowledge_tracing_bench\run_benchmark.py

# 出图表与分析（约 30 秒）
& $py D:\Tender\experiments\knowledge_tracing_bench\analyze.py
\end{verbatim}

\subsection{产物位置}

\begin{longtable}{@{}p{8.0cm}p{5.6cm}@{}}
\toprule
\textbf{文件} & \textbf{内容} \\
\midrule
\texttt{outputs/metrics.csv} & 各模型指标 \\
\texttt{outputs/metrics\_with\_lr.csv} & 含逻辑回归控制组 \\
\texttt{outputs/per\_skill\_auc.csv} & 分知识点 AUC \\
\texttt{outputs/eta\_sensitivity.csv} & $\eta$ 敏感性扫描 \\
\texttt{outputs/test\_predictions.parquet} & 逐条预测（可做误差分析） \\
\texttt{outputs/fig\_*.png} & 五张图 \\
\texttt{REPORT.md} & 完整实验报告 \\
\bottomrule
\end{longtable}

\newpage

% ============================================================

\section{附录：第一期实验总结}
% ============================================================

按时间顺序，本阶段完成的工作：

\begin{longtable}{@{}p{1.0cm}p{6.0cm}p{6.6cm}@{}}
\toprule
\textbf{序号} & \textbf{做了什么} & \textbf{产出} \\
\midrule
1 & 用 OCR 识别了 67 页的图片版 PDF & 3.6 万字文本 \\
2 & 解析了 22 页的文案 PDF & 58 条结构化文案 \\
3 & 给每条文案做规则标注 & 栏目、钩子、CTA、情绪、合规风险 \\
4 & 建立平台排序模型 & 4 个平台的评分公式 \\
5 & 做了双口径对比 & 找出 12 条“口径敏感”选题 \\
6 & 建立招生漏斗 & 六级转化 + 系数依据 \\
7 & 检索了相关学术文献 & 13 个方向、带 DOI \\
8 & \textbf{做知识追踪基准实验} & \textbf{8 个模型对比、5 张图} \\
9 & 验证了 $\eta$ 参数 & 发现默认值偏大 15 倍 \\
10 & 做了消融实验 & 验证“分知识点状态”的必要性 \\
11 & 建立三端架构设计 & 教师端 / 学生端 / 销售端 \\
12 & 设计了数据契约与合规清单 & 10 条红线 \\
13 & 设计了平台化路径 & 链上链下边界、加盐承诺 \\
\bottomrule
\end{longtable}

\subsection{最重要的三个数字}

\begin{enumerate}[leftmargin=1.8em]
\item \textbf{0.02} —— $\eta$ 应该改成这个值（原为 0.30）。
\item \textbf{1.15 $\to$ 0.62} —— 改完之后 LogLoss 的改善幅度。
\item \textbf{42 : 35} —— Tender 与 BKT 在 77 个知识点上的胜负比。
\end{enumerate}

\subsection{最重要的一句话}

\begin{quote}
\large
\textbf{$\theta$ 这个表示可用但不占优：排序能力和标准方法打平，概率校准还不够好。\\
真正的问题不在表示，而在默认参数选错了优化目标——\\
文档调的是“排序好看”，产品需要的是“概率可信”。}
\end{quote}

\vspace{12pt}
\hrule
\vspace{6pt}
\begin{center}
\small 本文档由项目实际执行结果生成 ｜ 所有数值均可在 \texttt{outputs/} 中复核
\end{center}




\section{附录：第二期实验复现}
% ============================================================

按时间顺序：

\begin{longtable}{@{}p{0.8cm}p{6.2cm}p{6.8cm}@{}}
\toprule
\textbf{序号} & \textbf{做了什么} & \textbf{产出} \\
\midrule
1 & 把量化交易框架映射到平台算法 & 结构映射表、五个致命差异、因子衰减概念 \\
2 & 复用项目已有的 Hawkes 强度函数 & \texttt{hawkes.py} 模块 \\
3 & 合成回环验证 & 优化器能用，但有系统性低估 \\
4 & EdNet 真实事件流拟合 & 分支比 0.94（看似成功） \\
5 & \textbf{三路安慰剂检验} & \textbf{发现 0.93 是假象} \\
6 & \textbf{时间重标检验} & 引入 Weibull 更新过程做竞争模型 \\
7 & 模型比较 & Hawkes 优于 Weibull：43/45 \\
8 & \textbf{修正结论} & Hawkes 方向对、设定有误（38\% 通过） \\
9 & 尝试多尺度核 & 合成验证通过；EdNet 上 AIC 略好但重标无改善 \\
10 & 验证昼夜周期 & $\chi^2=23{,}155$，峰谷比 34.3 倍 \\
11 & 实现昼夜调制 Hawkes & AIC 30/30 改善，但重标无改善 \\
12 & \textbf{换成会话为单位} & \textbf{分支比 0.93$\to$0.54，KS 通过率 10\%$\to$100\%} \\
13 & 阈值敏感性 & 6 分钟到 2 小时范围内结论稳定 \\
\bottomrule
\end{longtable}

\subsection{最重要的四个数字}

\begin{enumerate}[leftmargin=2em]
\item \textbf{0.9318 / 0.9318} --- 真实的与打乱的分支比完全一样，第一处警报。
\item \textbf{38\%} --- 单尺度 Hawkes 的时间重标通过率，三次改进都没能提高。
\item \textbf{23,155} --- 昼夜结构的卡方值（均匀应为 23）。
\item \textbf{100\%} --- 换成会话为单位后的通过率。
\end{enumerate}

\subsection{最重要的一句话}

\begin{quote}
\large
\textbf{问题从来不是模型不够强，而是基本单位选错了。}\\[6pt]
前三步都在把模型改复杂，第四步是把单位换对---
而能发现“该换单位”的，是那两个检验工具。
\end{quote}

\subsection{两期手册的关系}

\begin{center}
\begin{tabular}{p{4.2cm}p{5.4cm}p{5.0cm}}
\toprule
& \textbf{第一期} & \textbf{第二期} \\
\midrule
主题 & 建模与基准实验 & 检验与换单位 \\
核心工具 & AUC / RMSE / LogLoss & 安慰剂 / 时间重标 \\
数据 & ASSISTments（数学题） & EdNet（学习事件流） \\
主要发现 & $\eta$ 默认值偏大 15 倍 & 基本单位应是会话 \\
结论 & 概率校准比排序更重要 & 换单位比加复杂度更重要 \\
\bottomrule
\end{tabular}
\end{center}

\vspace{12pt}
\hrule
\vspace{6pt}
\begin{center}
\small 本文档由项目实际执行结果生成 ｜
所有数值均可在 \texttt{experiments/platform\_reverse\_lab/outputs/} 中复核
\end{center}




\section{附录：第三期实验与数据}
% ============================================================

\subsection{主要数据实验}

\begin{longtable}{@{}p{4.4cm}p{3.2cm}p{7.2cm}@{}}
\toprule
\textbf{实验} & \textbf{数据} & \textbf{结论} \\
\midrule
会话大小分布 & 653 次会话 & 幂律（$\alpha$=1.342），AIC 领先 1,140 \\
会话内节奏 & 10,000+ 间隔 & 无疲劳，有热身（42.2s $\to$ 38.5s） \\
周周期 & 全会话 & $\chi^2$=59（星期）、781（周内小时） \\
参数化调制 & 40 人 & 作息解释约 1/3，AIC 最优为 7 参数 \\
共享形状 & 21 人 & 参数减半、拟合不变；一日双峰 \\
效应量标定 & OULAD 26,074 人 & 中位 0.053、最大 0.289 \\
作息周期图 & EdNet 41 人 & $R_{24}$=0.533、$R_{168}$=0.320、$B$=0.659；26 人日内主导 \\
节律基线 & EdNet 25 人 & 合成验证 AUC=1.0；458 个学生-日 \\
元元评分 & 六种设计场景 & 理想执行 MC=0.950 / 现实规模 MC=0.738；两项致命缺陷被一票否决 \\
元认知公开数据 & ASSISTments 1,936 人 & 敏感性中位 0.585；与准确性 $r=+0.513$；求助行为不可用作信号 \\
功效分析 & 模拟 & $N$=19 只能测 $\lvert r\rvert \ge 0.605$ \\
管线自检 & 合成 95 人 & H1 $\rho$=0.627，无假阳性 \\
JOL 公开数据（实验 A） & OSF yaguz：84+80+53+41 人 & 相关 $0.46$--$0.51$；难易效应 $p<10^{-21}$；全局 JOL 不优于班平均基线 \\
JOL 课堂功效（实验 A） & 模拟 19 人 $\times$ 10 次 & 平均偏差 MDES $=0.10$；$r=0.5$ 功效 $0.61$；难易效应功效 $1.00$；$\delta=0.01$/次 功效 $0.87$ \\
根号讲义提取（实验 B） & 21 份讲义 PDF & 五/七/八/九/十讲共 165 个 $\surd$ 标记；其余各讲 0 \\
教师判断效度（实验 B） & 9 篇 56--65，6 人 & 4 标记 vs 6 未标；AUC $=0.5625$，置换 $p=0.795$；95\% CI $=[0.333,0.812]$ \\
学生自标元认知（实验 B） & 合成自检 & 10 题个体 AUC SD $0.14$--$0.18$；零模型 95\% 范围 $0.143$--$0.798$；真实学生数据等待录入 \\
难度基线 v2（实验 C） & 4 套测试、30 个项目、6--7 人/题 & Wilson 中位宽度 $0.591$；$n$ 为 $0.30/0.25/0.20$ 需 $39/58/93$ 人；10 题能力 SE $0.699$ \\
残差检验自检（实验 C） & 合成 300 次 & Wilson 覆盖率 $0.970$；假阳性率 $0.088$（偏高，故 $p<0.05$ 只当线索） \\
\bottomrule
\end{longtable}

\subsection{相关实验（E1--E6，2026-10-09）}

本节把“顺应天性、教师观察者、拆闭环”三条原则推进到可复核的实验。
详细报告见 \texttt{experiments/platform\_reverse\_lab/outputs/related\_experiments/全部相关实验执行报告.md}。

\begin{longtable}{@{}p{1.9cm}p{3.0cm}p{4.8cm}p{4.5cm}@{}}
\toprule
\textbf{实验} & \textbf{数据} & \textbf{真实结果} & \textbf{边界} \\
\midrule
E1 教师--学生--观察者 & 实验 B 教师 AUC 先验 & AUC $=0.5625$，折后可信度 $\approx0.409$ & 真实课堂数据待录入 \\
E2 节律--安排匹配 & EdNet：340 人、16{,}335 会话 & 参与度 $\beta=0.164$（$p=8.3\times10^{-10}$）；退出率 $\beta=-0.0034$（$p=0.0078$）；对齐 vs 错位 $d_z=0.572$ & 观察性数据；课堂 19 人只能测中等以上效应 \\
E3 习惯改变成本 & OULAD：3{,}829 单元、84{,}095 学生--周 & 参与度 $\beta=-1.134$（$p\approx5.9\times10^{-247}$）；控制当前点击后 $\beta=-0.431$（$p=1.9\times10^{-40}$）；退课 $\beta=+0.233$ & 观察性关联，不是随机强制改习惯 \\
E4 群体拓扑 & Stanford MOOCPosts：11{,}042 作者、62{,}857 边 & 触达：聚类 $+0.188$、桥接 $-0.096$；点赞：桥接 $+0.043$ & 共帖网络不是真实社会关系；触达不是学习结果 \\
E5 多教师 metameta & 合成 5--50 位教师 & 5 位教师、$sd_{\tau}=0.3$ 时 ANOVA 功效 $0.903$ & 真实多教师数据待采集 \\
E6 天性/安排先验 & 合成 19 人 $\times$ 10 次 + OULAD 先验 & 19 人 MDES $\approx0.35$；OULAD 中位 $0.053$ 需约 2{,}789 人，最大 $0.289$ 需约 92 人 & MBTI 只作连续先验，不作硬分类 \\
\bottomrule
\end{longtable}

\subsection{产物位置}

\begin{longtable}{@{}p{8.4cm}p{6.0cm}@{}}
\toprule
\textbf{目录} & \textbf{内容} \\
\midrule
\texttt{outputs/session\_model/} & 会话层刻画 \\
\texttt{outputs/session\_fourier/} & 参数化调制 \\
\texttt{outputs/shared\_rhythm/} & 分层共享形状 \\
\texttt{outputs/traits/} & MBTI 假设与功效分析 \\
\texttt{outputs/oulad\_traits/} & 26,074 人效应量标定 \\
\texttt{outputs/preregistration/} & 预注册研究包（含采集模板） \\
\texttt{outputs/jol/} & 实验 A：JOL 公开数据、功效分析、采集卡 \\
\texttt{outputs/root\_marks/} & 实验 B：讲义 $\surd$ 提取、教师判断效度、学生自标模板 \\
\texttt{outputs/difficulty\_baseline/} & 实验 C：题目难度、Wilson 区间、LOO 残差、v2 图 \\
\bottomrule
\end{longtable}

\subsection{最重要的十二个数字}

\begin{enumerate}[leftmargin=2em]
\item \textbf{1.342} --- 会话大小分布的幂律指数。
\item \textbf{509.7 vs 145.8} --- 自由格与参数化模型的 AIC（自由格差 364）。
\item \textbf{0.289} --- OULAD 标定出的最大真实效应。
\item \textbf{0.053} --- 同一批数据里的中位效应。
\item \textbf{92} --- 检测 0.289 所需的学生数（现在只有 19）。
\item \textbf{0.627} --- 管线自检中恢复出的真实效应（证明工具可用）。
\item \textbf{0.460--0.499} --- 公开 JOL 数据中 JOL 与实际的相关系数（Exp1）。
\item \textbf{$5.13\times10^{-22}$} --- Exp1 难易效应的配对检验 $p$：易项偏保守，难项偏过度自信。
\item \textbf{0.10} --- 19 人 $\times$ 10 次课堂设计能检测的最小平均 JOL 偏差（80\% 功效）。
\item \textbf{0.5625} --- 教师 $\surd$ 标记预测“更难”的 AUC（9 篇）；置换 $p=0.795$，没有显著证据。
\item \textbf{0.591} --- 难度基线的 Wilson 宽度中位；要压到 0.30/0.25/0.20，需要 $39/58/93$ 人/题。
\item \textbf{0.699} --- 10 题时学生能力的解析 SE（logit）；要到约 0.308，需要约 50 题。
\end{enumerate}

\subsection{最重要的一句话}

\begin{quote}
\large
\textbf{三个实验把“元认知”从一句口号变成了三组可复核的数字。}\\[6pt]
JOL 告诉我们：学生会预测，但预测不等于校准，也不自动优于“只说班平均”；
根号告诉我们：讲义上的 $\surd$ 是教师判断，5/7 篇单词甚至没有对照，学生自标必须另采；
难度基线告诉我们：6--7 人、10 题的课堂数据只能粗筛，不能精细诊断。
\textbf{真正的进步不是多了一个漂亮模型，而是知道哪些结论现在还不能下。}
\end{quote}

\vspace{12pt}
\hrule
\vspace{6pt}
\begin{center}
\small 本文档由项目实际执行结果生成 ｜
所有数值均可在 \texttt{experiments/platform\_reverse\_lab/outputs/} 中复核
\end{center}




\section{采集工具与课堂操作手册}
% ============================================================

\subsection{四个采集工具}

\begin{center}
\begin{longtable}{@{}p{3.2cm}p{5.4cm}p{5.0cm}@{}}
\toprule
\textbf{工具} & \textbf{路径} & \textbf{用途} \\
\midrule
E1 采集表 & \path{outputs/related_experiments/E1_observer/E1_teacher_student_observer.html} & 学生预测 + 教师预测 + 教师信心 + 实际 \\
E4 互动调查 & \path{outputs/related_experiments/E4_group_topology/E4_group_topology_survey.html} & 每周写 3 个联系人 + 1 个跨组联系人 \\
E5 多教师模板 & \path{outputs/related_experiments/E5_multi_teacher/multi_teacher_metameta_template.csv} & 5 位教师 $\times$ 19 学生 $\times$ 5 次 \\
E6 天性模板 & \path{outputs/related_experiments/E6_trait_prior/trait_prior_template.csv} & MBTI 连续分数 + 节律峰型 + 任务结构偏好 \\
\bottomrule
\end{longtable}
\end{center}

\subsection{六周课堂流程}

\begin{center}
\begin{longtable}{@{}p{1.8cm}p{4.2cm}p{7.2cm}@{}}
\toprule
\textbf{周} & \textbf{做什么} & \textbf{对应实验} \\
\midrule
W0 & 打印 E1 采集表、E4 互动调查表 & 准备 \\
W1 & 每次测试前学生写预测，教师写预测 + 信心 & E1 第 1 次 \\
W1 周五 & 填 E4 互动调查 & E4 基线 \\
W2 & 重复 E1；开始记录电子题的响应时间和修改次数 & E1、犹豫 \\
W3 & 轮换分组 + 跨组任务；继续 E1 & E4 干预 \\
W4 & 继续轮换；做三臂外部脚手架实验（可选） & E4、外部脚手架 \\
W5 & 恢复固定分组；继续 E1 & E4 回退 \\
W6 & 最后一次 E1 + E4；跑分析脚本 & 全部 \\
\bottomrule
\end{longtable}
\end{center}

\subsection{数据录入的三条硬规则}

\begin{enumerate}[leftmargin=2em]
\item \textbf{先预测，后看题}：学生先写预测，教师再独立写预测和信心；
\item \textbf{预测不许改}：写错可以划掉重写，但要保留痕迹；
\item \textbf{实际对题数由老师填}：不让学生自己填，避免事后合理化。
\end{enumerate}

\subsection{跑分析之前先看功效}

\begin{explainbox}{不要跳过功效}
19 人 $\times$ 5 次只能测中等以上效应。\\
如果结果是“不显著”，先问：是真的没有效应，还是功效不够？\\
第三期已经给过答案：$N=19$ 只能测 $|r|\ge0.605$。
\end{explainbox}

\newpage

% ============================================================

\section{附录：第四期实验与复现}
% ============================================================

\subsection{主要脚本}

\begin{longtable}{@{}p{6.2cm}p{7.0cm}@{}}
\toprule
\textbf{脚本} & \textbf{用途} \\
\midrule
run\_teacher\_student\_observer.py & E1 教师--学生--观察者 \\
run\_arrangement\_fit.py & E2 节律--安排匹配（EdNet） \\
run\_habit\_change\_cost.py & E3 习惯改变成本（OULAD） \\
run\_group\_topology.py & E4 群体拓扑（Stanford MOOCPosts） \\
run\_multi\_teacher\_trait.py & E5 多教师 metameta + E6 天性先验 \\
\bottomrule
\end{longtable}

\subsection{主要产物目录}

\begin{longtable}{@{}p{8.4cm}p{5.4cm}@{}}
\toprule
\textbf{目录} & \textbf{内容} \\
\midrule
\path{outputs/related_experiments/E1_observer/} & E1 采集表、功效、合成自检 \\
\path{outputs/related_experiments/E2_arrangement_fit/} & EdNet 会话、匹配度、图 \\
\path{outputs/related_experiments/E3_habit_change/} & OULAD 学生--周、习惯改变、图 \\
\path{outputs/related_experiments/E4_group_topology/} & 共帖边、作者网络指标、图 \\
\path{outputs/related_experiments/E5_multi_teacher/} & 多教师模板、ANOVA 功效 \\
\path{outputs/related_experiments/E6_trait_prior/} & 天性模板、功效、恢复 \\
\path{outputs/related_experiments/all_related_experiments_report.md} & E1--E6 完整报告 \\
\bottomrule
\end{longtable}

\subsection{最重要的数字}

\begin{enumerate}[leftmargin=2em]
\item \textbf{0.409}：教师观察的折后可信度（AUC=0.5625）。
\item \textbf{0.164}：EdNet 上节律匹配对参与度的效应（$p=8.3\times10^{-10}$）。
\item \textbf{0.572}：对齐 vs 错位会话的配对效应量 $d_z$。
\item \textbf{-1.134}：OULAD 上习惯改变对下一周参与度的效应。
\item \textbf{-0.431}：控制当前点击后的习惯改变效应（$p=1.9\times10^{-40}$）。
\item \textbf{+0.188 / -0.096}：Stanford 论坛上聚类/桥接对触达的标准化效应。
\item \textbf{0.903}：5 位教师、$sd_\tau=0.3$ 时 ANOVA 功效。
\item \textbf{0.35}：19 人 $\times$ 10 次的 arrangement-fit MDES。
\item \textbf{2,789 / 92}：OULAD 中位/最大效应所需人数。
\item \textbf{1.342}：会话幂律指数；不是黄金比。
\end{enumerate}

\subsection{一句话总结}

\begin{quote}
\large
\textbf{犹豫是元认知的行为证据；外部只能引出和训练，不能硬植入；\\
基线是用户与环境的耦合稳定点，不是固定性格；\\
黄金比目前只是隐喻，真正可检验的是闭环、自相似和尺度不变。}
\end{quote}

\vspace{12pt}
\hrule
\vspace{6pt}
\begin{center}
\small 本文档由项目实际执行结果生成 ｜
所有数值均可在\\
\texttt{experiments/platform\_reverse\_lab/}\\
\texttt{outputs/related\_experiments/} 中复核
\end{center}



\section{课堂六周操作清单}

\begin{center}
\begin{longtable}{@{}p{1.6cm}p{5.0cm}p{7.4cm}@{}}
\toprule
\textbf{周} & \textbf{做什么} & \textbf{对应实验/工具} \\
\midrule
W0 & 打印 E1 采集表、E4 互动调查表 & 准备 \\
W1 & 每次测试前学生写预测，教师写预测 + 信心 & E1 \\
W1 周五 & 填 E4 互动调查 & E4 基线 \\
W2 & 重复 E1；电子题记录响应时间和修改次数 & E1、犹豫 \\
W3 & 轮换分组 + 跨组任务；继续 E1 & E4 干预 \\
W4 & 继续轮换；三臂外部脚手架实验（可选） & E4、外部脚手架 \\
W5 & 恢复固定分组；继续 E1 & E4 回退 \\
W6 & 最后一次 E1 + E4；跑分析脚本 & 全部 \\
\bottomrule
\end{longtable}
\end{center}

\subsection{三条硬规则}

\begin{enumerate}[leftmargin=2em]
\item 先预测，后看题；
\item 预测不许改；
\item 实际对题数由老师填。
\end{enumerate}



\part{治理与合规}

\section{合规红线}
% ============================================================

这一章把“不能做什么”列成一张清单。这些不是法律建议，是产品设计的自我约束。

\subsection{十条红线}

\begin{longtable}{@{}p{1.0cm}p{6.6cm}p{6.0cm}@{}}
\toprule
\textbf{序号} & \textbf{红线} & \textbf{原因} \\
\midrule
R1 & 情绪 / 心智数据永不上链 & 一旦上链无法删除 \\
R2 & 不可从群体信息反推出个人状态 & 小样本极易定位到人 \\
R3 & 个体情绪状态不呈现给任何教师 & 贴标签会造成预期偏差 \\
R4 & 学生数据不得用于任何营销触达 & 未成年人保护 \\
R5 & 不采集未成年人生物特征 & 人脸 / 声纹属于敏感信息 \\
R6 & 不发币、不 ICO、不证券化 & 金融监管 \\
R7 & 学生数据不出校（离校凭证除外） & 校内模式的立身之本 \\
R8 & 销售端禁用学生个体数据 & 防止把教学数据变成营销素材 \\
R9 & 教育广告禁止承诺升学、禁止制造焦虑 & 广告法 \\
R10 & 未满 18 岁上链需监护人单独同意 & 未成年人保护 \\
\bottomrule
\end{longtable}

\subsection{三条最容易被忽略的红线}

\subsubsection{R3：为什么教师看不到学生的情绪标签}

给人贴情绪标签有两重危害：

\begin{enumerate}[leftmargin=1.8em]
\item \textbf{合规上}：这属于处理敏感个人信息，需要额外的合法性基础；
\item \textbf{教学上}：老师一旦认为“这个学生情绪不好”，
      就会不自觉地降低对他的期待——\hl{这是自我实现的预言}。
\end{enumerate}

\textbf{正确的做法}：

\begin{center}
\begin{tabular}{p{6.4cm}p{6.0cm}}
\toprule
\warn{不要这样} & \hl{应该这样} \\
\midrule
“张三情绪低落，请多关注” & “张三·第 9 篇·时态题连续 3 次未掌握” \\
“这个学生状态很差” & “本班今日心智化窗口偏窄，建议降低新知识负荷” \\
\bottomrule
\end{tabular}
\end{center}

也就是说：\textbf{情绪数据参与“建议的生成”，但不呈现为任何学生的标签。}

\subsubsection{R7 与“离校凭证”的例外}

R7 说学生数据不出校。但有一个例外：\textbf{学生成年后，可以带走自己的成果凭证}。

\begin{explainbox}{触发条件必须精确}
\warn{错误理解}：“离开机构就可以上链”。\\
\textbf{为什么错}：一个初二学生离开机构，他仍然是未成年人，保护要求一点没变。
“离开机构”不会让一个孩子变成成年人。\\[4pt]
\hl{正确条件}：
\begin{equation}
\text{可签发} = \big(\text{年满 18 周岁}\big) \;\lor\; \big(\text{未满 18 但监护人与本人均同意}\big)
\end{equation}
\end{explainbox}

上链的内容范围：

\begin{center}
\begin{tabular}{p{6.6cm}p{5.8cm}}
\toprule
\textbf{可以上链} & \textbf{绝不上链} \\
\midrule
学习内容范围（完成 21 篇讲义） & $\theta$、$\phi$、$\gamma$、情绪色相 $c$ \\
时间跨度 & 反思文本 \\
掌握度变化的区间（如“提升 $\geq 15\%$”） & 任何可反推心智状态的量 \\
机构签名 + 学生标识 & 原始逐题对错（加盐承诺可以） \\
\bottomrule
\end{tabular}
\end{center}

\begin{explainbox}{一句话总结这条边界}
上链的是“\textbf{学习成果的凭证}”，不是“\textbf{学习过程的监控录像}”。
\end{explainbox}

\subsubsection{加盐承诺：为什么必须加随机数}

回到第 2.6 节讲的陷阱。假设某班正确率是 62\%：

\begin{center}
\begin{tabular}{p{4.6cm}p{7.8cm}}
\toprule
\textbf{做法} & \textbf{后果} \\
\midrule
直接存 62\% & 数据完全暴露 \\
存 $H(0.62)$ & \warn{枚举 101 次就能还原}（0\% 到 100\%） \\
\textbf{存 $H(0.62 \;\|\; \text{随机数})$} & \hl{安全}。随机数有无数种可能，无法枚举 \\
\bottomrule
\end{tabular}
\end{center}

\subsection{必须排除的具体操作}

在做平台研究时，有一类“看起来有效”的做法必须排除：

\begin{longtable}{@{}p{7.0cm}p{6.6cm}@{}}
\toprule
\textbf{操作} & \textbf{排除原因} \\
\midrule
买特定号段的手机卡注册账号 & 规避平台风控 \\
连续 7 天每天刷 2 小时同品类视频“养号” & 人为制造虚假兴趣信号 \\
组织 50 人点赞、完播 & 刷量 \\
直播间安排“水军”配合逼单 & 虚假互动、误导消费者 \\
“没效果退全款”这类效果承诺 & 广告合规风险 \\
用愧疚感促使家长下单 & 操纵性营销 \\
用数字人批量生产同质内容 & 低质内容治理风险 \\
\bottomrule
\end{longtable}

\begin{explainbox}{为什么这些要单独列出来}
因为这些操作\textbf{在短期内确实能提升数据}，所以很有诱惑力。
但它们的代价是：账号风险、法律风险，以及——\hl{研究结论会被污染}。
如果数据是靠刷出来的，那么“什么内容有效”这个结论就不可信了。
\end{explainbox}

\newpage

% ============================================================

\section{全部翻车清单}

\begin{enumerate}[leftmargin=2em]
\item \textbf{第一期}：默认参数 $\eta$ 太大；场景切换惩罚没有用。
\item \textbf{第二期}：自激过程看起来很美，但安慰剂和时间重标显示是过拟合；两次改进失败，最后靠“换单位”成功。
\item \textbf{第三期}：AIC 漏计参数导致结论被推翻；振幅不可识别；19 人功效不足；“单词测试零区分度”不能推广。
\item \textbf{第四期}：E5 的 Fisher-$z$ 方差检验设定错误，改用 ANOVA；E4 的“抱团有害”假设只对了一半；MBTI 不能当天生硬约束；黄金比没有证据。
\end{enumerate}

\begin{explainbox}{为什么保留翻车清单}
因为\textbf{知道自己哪里错过，比多写一个漂亮结论更重要}。\\
这份手册的价值，不是它从来不犯错，而是它把错误也写下来。
\end{explainbox}



\part{总览、文件索引与总结}

\section{四期总览}

\begin{center}
\begin{longtable}{@{}p{1.4cm}p{2.6cm}p{3.6cm}p{4.6cm}p{2.6cm}@{}}
\toprule
\textbf{期} & \textbf{解决什么} & \textbf{核心方法} & \textbf{最重要的结果} & \textbf{产物} \\
\midrule
第一期 & 怎么建模、怎么算指标 & 会话、节律、知识追踪、$\theta$ & 会话是比“一次点击”更好的基本单位 & 全流程手册 \\
第二期 & 怎么检查模型是不是骗自己 & 安慰剂、时间重标、换单位 & 自激归零是过拟合；会话模型闭合三层结构 & 第二期手册 \\
第三期 & 怎么防止事后找理由 & 功效、预注册、元认知、metameta & 19 人只能测 $|r|\ge0.605$；元认知敏感性中位 0.585 & 第三期手册 \\
第四期 & 怎么把信号接回系统 & 犹豫、外部脚手架、耦合基线、群体拓扑 & 顺应节律更好；改变节奏有成本；小圈子与桥接各有优势 & 第四期手册 \\
\bottomrule
\end{longtable}
\end{center}


\section{全部文件与数据索引}

\subsection{四期手册}

\begin{longtable}{@{}p{8.6cm}p{5.0cm}@{}}
\toprule
\textbf{文件} & \textbf{内容} \\
\midrule
Tender零基础全流程教学手册.pdf & 第一期：建模与指标 \\
Tender零基础教学手册-第二期.pdf & 第二期：检验与防错 \\
Tender零基础教学手册-第三期.pdf & 第三期：预注册与元认知 \\
Tender零基础教学手册-第四期.pdf & 第四期：从元认知信号到耦合基线 \\
\bottomrule
\end{longtable}

\subsection{主要脚本}

\begin{longtable}{@{}p{6.4cm}p{7.0cm}@{}}
\toprule
\textbf{脚本} & \textbf{用途} \\
\midrule
run\_teacher\_student\_observer.py & E1 教师--学生--观察者 \\
run\_arrangement\_fit.py & E2 节律--安排匹配（EdNet） \\
run\_habit\_change\_cost.py & E3 习惯改变成本（OULAD） \\
run\_group\_topology.py & E4 群体拓扑（Stanford MOOCPosts） \\
run\_multi\_teacher\_trait.py & E5 多教师 metameta + E6 天性先验 \\
\bottomrule
\end{longtable}

\subsection{主要产物}

\begin{longtable}{@{}p{8.4cm}p{5.2cm}@{}}
\toprule
\textbf{目录} & \textbf{内容} \\
\midrule
outputs/related\_experiments/E1\_observer/ & E1 采集表、功效、合成自检 \\
outputs/related\_experiments/E2\_arrangement\_fit/ & EdNet 会话、匹配度、图 \\
outputs/related\_experiments/E3\_habit\_change/ & OULAD 学生--周、习惯改变、图 \\
outputs/related\_experiments/E4\_group\_topology/ & 共帖边、作者网络指标、图 \\
outputs/related\_experiments/E5\_multi\_teacher/ & 多教师模板、ANOVA 功效 \\
outputs/related\_experiments/E6\_trait\_prior/ & 天性模板、功效、恢复 \\
\bottomrule
\end{longtable}


\section{最重要的 20 个数字}

\begin{enumerate}[leftmargin=2em]
\item \textbf{1.342}：会话大小幂律指数。
\item \textbf{0.627}：管线自检恢复出的真实效应。
\item \textbf{0.605}：19 人相关分析能检测的最小 $|r|$。
\item \textbf{0.053}：OULAD 中位效应。
\item \textbf{0.289}：OULAD 最大效应。
\item \textbf{92}：检测 0.289 所需人数。
\item \textbf{0.585}：ASSISTments 响应时间的元认知敏感性中位 AUC。
\item \textbf{+0.513}：ASSISTments 元认知敏感性与准确性的相关。
\item \textbf{0.5625}：教师题目级判断 AUC。
\item \textbf{0.409}：教师观察折后可信度。
\item \textbf{0.164}：EdNet 节律匹配对参与度的效应。
\item \textbf{0.572}：对齐 vs 错位会话的配对 $d_z$。
\item \textbf{-1.134}：OULAD 习惯改变对下一周参与度的效应。
\item \textbf{-0.431}：控制当前点击后的习惯改变效应。
\item \textbf{+0.233}：OULAD 习惯改变对退课率的效应。
\item \textbf{+0.188 / -0.096}：Stanford 论坛聚类/桥接对触达的效应。
\item \textbf{+0.043}：桥接对点赞的效应。
\item \textbf{0.903}：5 位教师、$sd_\tau=0.3$ 时 ANOVA 功效。
\item \textbf{0.35}：19 人 $\times$ 10 次的 arrangement-fit MDES。
\item \textbf{1.618}：黄金比；目前只是隐喻，不是本项目结论。
\end{enumerate}


\section{一句话总结}

\begin{quote}
\large
\textbf{第一期学会“做”，第二期学会“验”，第三期学会“防”，第四期学会“接”。}\\[6pt]
\textbf{犹豫是元认知的行为证据；外部只能引出和训练，不能硬植入；}\\
\textbf{基线是用户与环境的耦合稳定点，不是固定性格；}\\
\textbf{黄金比目前只是隐喻，真正可检验的是闭环、自相似和尺度不变。}
\end{quote}

\vspace{12pt}
\hrule
\vspace{6pt}
\begin{center}
\small 本文档由四期手册和项目实际执行结果合并生成 ｜
所有数值均可在项目 outputs 目录中复核
\end{center}





\part{增补：开花、结果、森林与认知构念实验}

\section{为什么需要这一部分}

“开花、结果、森林”最早只是产品隐喻。隐喻可以帮助理解，但不能直接用于测量、预测或宣称效果。为了让它成为可检验模型，本增补做三件事：

\begin{enumerate}[leftmargin=2em]
\item 把三个隐喻拆成可观测的构念；
\item 给出公式、符号、阈值和证据等级；
\item 用公开数据和合成案例做初步检验，并明确哪些结论还不能下。
\end{enumerate}

\begin{explainbox}{先给结论}
开花不是“突然懂了”，而是知识经过重复提取、间隔、延迟保持和迁移，逐渐进入稳定认知结构。\\
结果不是“人多力量大”，而是不同认知结构在准确分工和协调下发生重组，产生超过个体基线的新颖且有用产出。\\
森林不是人数，而是不同认知能够长期共存、持续调用和重组的环境。\\
目前公开数据支持“继续验证”，不支持“已经证明”。
\end{explainbox}

\section{三个构念的正式定义}

\subsection{开花：从工作记忆表征到稳定认知结构}

设学习者在延迟时间 $t$ 的提取正确率为 $R_t$，新情境迁移成绩为 $T$，结构化解释或图式一致性评分为 $S$，认知负荷为 $C$。定义：

\begin{equation}
B=\sigma\left(\beta_0+\beta_R \bar R+\beta_T T+\beta_S S-\beta_C C\right)
\end{equation}

其中，$\sigma(z)=1/(1+e^{-z})$ 是 Logistic 函数；$\bar R$ 是 24 小时、7 天和 30 天延迟提取的加权平均；$T,S,C\in[0,1]$。默认阈值取 $B\ge0.65$，并要求 $R_{7d}\ge0.60$、$T\ge0.50$。

符号说明：
\begin{itemize}[leftmargin=2em]
\item $B$：开花得分，表示证据综合后知识进入稳定认知结构的概率性程度；
\item $\beta_0$：截距；$\beta_R,\beta_T,\beta_S,\beta_C$：各证据的权重；
\item $R_t$：延迟 $t$ 后无提示或低提示提取正确率；
\item $T$：换材料表面、换任务形式后是否仍能迁移；
\item $S$：能否用结构化解释说明知识之间的关系；
\item $C$：完成任务时的心理努力或认知负荷，越高越不利。
\end{itemize}

\warn{即时答对、熟悉感或单次自评不能代替上述证据。}

\subsection{结果：不同认知结构产生新组合}

设群体成员 $i$ 的独立新颖性为 $N(I_i)$，群体产出为 $G$，综合比例为 $X$，有用性为 $U$，交互记忆系统为 $M$，关系冲突为 $K$。定义：

\begin{equation}
Q=\sigma\left(\rho_0+\rho_N\max(0,N(G)-\max_i N(I_i))+\rho_X X+\rho_U U+\rho_M M-\rho_K K\right)
\end{equation}

默认阈值取 $Q\ge0.35$，并要求 $N(G)>\max_i N(I_i)$、$U\ge0.60$、$X\ge0.50$。

符号说明：
\begin{itemize}[leftmargin=2em]
\item $Q$：结果得分；
\item $N(G)-\max_i N(I_i)$：群体相对最佳个体的新颖性增益；
\item $X$：最终产出中能够追溯到至少两位成员独立贡献的比例；
\item $U$：盲评质量或真实任务效用；
\item $M$：群体是否知道彼此擅长什么，并能准确调用；
\item $K$：关系冲突和协调损耗。
\end{itemize}

\warn{讨论时长、想法数量、群体满意度不能单独作为“结果”成立证据。}

\subsection{森林：可持续的认知重组环境}

设认知多样性为 $D$，交互记忆为 $M$，贡献均衡为 $E$，历史重组成功为 $H$，趋同为 $X$。定义：

\begin{equation}
F=\sigma\left(\phi_0+\phi_D D+\phi_M M+\phi_E E+\phi_H H-\phi_X X\right)
\end{equation}

默认阈值取 $F\ge0.60$，并要求 $D\ge0.50$、$M\ge0.50$、$X<0.50$。

这里的森林不是人数，也不是所有成员达成一致，而是差异能够保留、知识能够被正确调用、组合能够重复发生。

\section{本项目初步实验}

\subsection{EdNet 真实时间间隔}

在 4,172 名学习者、11,495 道题和 1,247,784 条重复作答记录上，按学习者分组进行五折交叉验证并重复五次：

\begin{center}
\begin{tabular}{lrr}
\toprule
模型 & AUC & LogLoss \\
\midrule
基线特征 & 0.6648 & 0.6472 \\
加入真实时间间隔 & 0.6839 & 0.6345 \\
\midrule
增量 & +0.0191 & -0.0127 \\
\bottomrule
\end{tabular}
\end{center}

间隔四分位的观察正确率依次为 36.0\%、48.2\%、52.1\% 和 68.8\%，并非严格单调。

\subsection{ASSISTments 事件间隔代理}

在 3,913 名学习者、81 个技能和 238,869 条首次后的机会记录上，用“两次同技能机会之间的其他事件数”作为间隔代理。加入该代理只带来 $+0.00033$ AUC 和 $-0.00021$ LogLoss 的增量。

因此，事件位置不能替代真实时间。产品必须记录真实时间、延迟提取和迁移任务。

\subsection{群体构念评分模拟}

在 10,000 个已知类别的合成案例中：

\begin{center}
\begin{tabular}{lrrrr}
\toprule
构念 & AUROC & 准确率 & 敏感度 & 特异度 \\
\midrule
结果 & 0.993 & 83.3\% & 67.1\% & 99.4\% \\
森林 & 0.999 & 84.2\% & 68.4\% & 100\% \\
\bottomrule
\end{tabular}
\end{center}

这说明评分规则能够拒绝“最佳个人答案冒充群体重组”和“高趋同环境冒充森林”，但\textbf{不是}真人群体实验。

\section{项目实现与数据字段}

\begin{center}
\begin{longtable}{@{}p{3.6cm}p{5.0cm}p{6.2cm}@{}}
\toprule
对象 & 机制 & 最小字段 \\
\midrule
工作记忆负荷 & 任务难度、提示、心理努力 & load\_estimate, effort\_rating, cue\_level \\
提取与间隔 & 检索、间隔、反馈 & retrieval\_attempt, delay, accuracy, feedback\_type \\
认知结构 & 图式连接、解释、迁移 & schema\_links, explanation\_score, transfer\_context \\
开花判定 & 延迟保持 + 迁移 + 低负荷 & bloom\_probability, retention\_7d, transfer\_score \\
群体重组 & 个体基线、综合、交互记忆 & member\_baseline, synthesis\_links, tms\_score \\
结果判定 & 新颖性 + 有用性 + 组合性 & novelty\_gain, utility\_score, result\_probability \\
森林健康 & 多样性、协调、趋同 & cognitive\_diversity, contribution\_entropy, homogenization \\
\bottomrule
\end{longtable}
\end{center}

项目已经提供：

\begin{verbatim}
GET  /cognition/constructs
POST /cognition/assessments/{student_id}
GET  /cognition/assessments/{student_id}
\end{verbatim}

评分结果写入 DAOQL 的 cognitive\_construct\_assessment 事件，并进入审计链。默认权重和阈值是待校准产品先验，不是科学常数。

\section{证据边界与后续实验}

\begin{explainbox}{不能宣称的事情}
不能宣称“开花已经被实验证明”；\\
不能宣称“间隔越久越有效”；\\
不能宣称“不同人在一起就会创造”；\\
不能宣称“合成模拟证明了真实群体协作有效”。
\end{explainbox}

下一步需要真人随机对照实验：

\begin{enumerate}[leftmargin=2em]
\item 开花：集中重复、间隔提取、间隔提取 + 结构化解释 + 迁移任务三组对照；
\item 结果：最佳成员汇总、纯头脑风暴、认知画像 + 综合支架、反趋同提示四组对照；
\item 森林：多轮协作网络追踪，测量多样性、贡献均衡、交互记忆、重组和趋同。
\end{enumerate}

完整证据、公式、测量协议与参考文献见：

\begin{itemize}[leftmargin=2em]
\item Tender-认知机制学术补编-V1.2.pdf；
\item Tender-认知构念实验报告-V1.0.pdf；
\item experiments/cognitive\_constructs\_bench/REPORT.md。
\end{itemize}

\begin{quote}
\large
\textbf{开花，是知识从工作记忆进入稳定认知；}\\
\textbf{结果，是不同认知重组出单个人不会产生的新东西；}\\
\textbf{森林，是让这种重组持续发生的多样认知环境。}
\end{quote}



\part{增补二：环境型多 Agent 与真桌面弹幕}

\section{为什么不再需要一个日常操作页面}

传统 App 要求用户主动打开、查找功能、填写表单，再决定下一步。Tender 的新方向不是让用户更频繁地打开页面，而是让 Agent 存在于用户正在做的事情旁边：

\begin{quote}
\textbf{后台理解状态，在恰当时机出现一句真的有用的话。}
\end{quote}

这不表示“没有控制权”。日常操作页被取消，但隐私、授权、历史和退出入口仍保留在控制面中。

\section{日常产品面：桌面圆桌而不是操作页}

V6.6 的日常产品面不是网页，也不是要求用户每天打开的任务面板。它由三层组成：

\begin{enumerate}[leftmargin=2em]
\item 全屏桌面弹幕层：独立于浏览器运行，始终置顶、无边框、鼠标穿透、不抢焦点；
\item 功能型 Agent 圆桌：每个 Agent 从不同专业视角给出建议；
\item 最小控制面：回复、静默、清空和退出都可以通过快捷键完成。
\end{enumerate}

桌面弹幕采用浅灰半透明底和黑色文字。顶部显示本轮在场的全部 Agent，但屏幕上最多同时保留两条弹幕，其余自动排队接力。这样既保留多人协作感，又避免用户不知道该看哪一条。

\begin{quote}
一轮判断可以产生多位 Agent 的观点；用户当前只需要阅读其中最重要的两条。
\end{quote}

\section{什么时候才允许出现}

允许触发的情形包括：

\begin{itemize}[leftmargin=2em]
\item 连续卡住或同类错误反复出现；
\item 长时间空闲或频繁切换窗口；
\item 深度专注过久，适合做一次短提取；
\item 用户主动求助；
\item 高风险安全情形需要升级，但普通弹幕不承担高风险沟通。
\end{itemize}

\section{打扰预算}

\begin{longtable}{@{}p{4.2cm}p{7.5cm}@{}}
\toprule
规则 & 含义 \\
\midrule
冷却时间 & 默认每 3 分钟最多一条 \\
每日限额 & 默认每天最多 8 条 \\
静默 & 用户可静默 30 分钟 \\
负反馈 & 选择“不是这样”后立即停止当前提示并修正判断 \\
深度工作 & 默认不打扰 \\
升级 & 高风险内容进入安全流程，不使用普通弹幕 \\
\bottomrule
\end{longtable}

\section{提示层级}

\begin{enumerate}[leftmargin=2em]
\item 状态点：只表示“我在”，不打断；
\item 单条弹幕：一句话加一个最小动作；
\item 展开卡：用户主动查看证据和原因；
\item 对话：用户主动发起，或高风险情况下升级。
\end{enumerate}

\section{默认监测边界与本地归一化}

桌面端在用户授权后默认提取以下行为特征：

\begin{itemize}[leftmargin=2em]
\item 前台应用类别：学习、文档、编程、视频、聊天、游戏、浏览器；
\item 短任务提示：如“高考报名”“语法填空”“PDF”“数学”；
\item 空闲时长；
\item 前台窗口切换次数；
\item 键盘和鼠标活动强度；
\item 连续专注时长。
\end{itemize}

窗口标题只在本地用于归一化分类，例如把“某网站——高考报名”映射为 \texttt{study} 和“高考报名”。原始窗口标题不会上传服务器，也不写入 DAOQL。

桌面端不读取：

\begin{itemize}[leftmargin=2em]
\item 屏幕截图；
\item 其他窗口正文；
\item 剪贴板；
\item 摄像头与麦克风；
\item 文件内容。
\end{itemize}

如果未来增加 OCR、可访问性树或可穿戴设备，必须重新取得授权、完成隐私影响评估并保留退出机制。

\section{功能型多 Agent 圆桌}

V6.6 的圆桌成员不是同一内容的不同性格版本，而是功能型 Agent：

\begin{longtable}{@{}p{3.6cm}p{5.0cm}p{4.0cm}@{}}
\toprule
Agent & 专业功能 & 典型输出 \\
\midrule
学习策略师 & 任务拆解、优先级、复习节奏 & 两分钟动作、完成条件 \\
学科专家 & 知识点、先修关系、错因 & 卡在概念、公式、条件或计算 \\
神经认知师 & 注意力、工作记忆、认知负荷 & 切换成本、休息窗口、单窗口策略 \\
心理调节师 & 焦虑、回避、动机、自我评价 & 区分逃避与疲劳、降低羞耻 \\
营养健康师 & 睡眠、饮食、补水、身体能量 & 困、饿、头痛、眼疲劳的检查 \\
生涯规划师 & 专业兴趣、目标路径、长期证据 & 当前动作与目标专业的关系 \\
安全保护官 & 隐私、身心安全、求助边界 & 证件信息、健康风险和升级 \\
元认知教练 & 策略觉察、监控和复盘 & 当前策略、证据、下一步调整 \\
\bottomrule
\end{longtable}

\subsection{Triage Agent：先分诊，再调用专家}

V6.6 的 Triage Agent 先读取用户消息和行为信号，输出问题类型、置信度、风险等级、证据、替代解释和相关专家列表。专家圆桌负责配置可用专家，Triage Agent 决定本轮真正被调用的专家。

例如：

\begin{itemize}[leftmargin=2em]
\item 用户说“我很焦虑”，调用心理调节师和神经认知师；
\item 用户说“数学题不会”，调用学科专家、学习策略师和元认知教练；
\item 用户说“我困得看不进去”，调用营养健康师和神经认知师；
\item 用户说“报名还没确认”，调用安全保护官、学习策略师和生涯规划师。
\end{itemize}

无关专家保持静默。低置信度时，Triage Agent 先提出一个最小澄清问题，不贸然调用专业 Agent。

功能型 Agent 必须输出非冗余内容。若两条建议语义相似，应合并或只保留更具体的一条。

设第 \(t\) 轮显示 \(N_t\) 条消息，语义向量为 \(\mathbf v_i\)，则语义冗余定义为：

\[
R_t=\frac{2}{N_t(N_t-1)}\sum_{i<j}\cos(\mathbf v_i,\mathbf v_j)
\]

其中 \(R_t\) 越高，表示重复越多。产品先用最多两条并发和低冗余约束来控制认知负荷。

\section{工程接口}

\begin{verbatim}
POST /ambient/plan/{student_id}
POST /ambient/feedback/{student_id}
POST /ambient/reply/{student_id}
GET  /ambient/personas
GET  /ambient/persona/{student_id}
POST /ambient/persona/{student_id}
GET  /ambient/ensemble/{student_id}
POST /ambient/ensemble/{student_id}
GET  /ambient/status/{student_id}
\end{verbatim}

规划器必须执行冷却、每日限额和静默窗口；每个提示都要说明为什么出现，并将提示与反馈写入 DAOQL 和审计链。

\section{两个首批场景}

\subsection{高考报名}

高考报名是高利害、强截止、清单化的流程任务。Tender 只监测授权后的流程状态：距离截止还有多久、是否还有未确认事项、确认、缴费或照片审核是否完成。

弹幕示例：

\begin{quote}
报名截止前还有 2 项未确认。先处理最紧急的一项，不在这里展示证件号码。
\end{quote}

安全边界：不在弹幕中展示身份证号、家庭住址、照片或支付信息；不替代官方报名系统。

\subsection{提分}

提分是长期学习、错点复发、最小练习和结果验证的循环。Tender 可以提醒一个最常错、最影响分数的小点，并要求用户完成一个最小练习，再用真实结果验证。

弹幕示例：

\begin{quote}
时间不多，先不要铺开复习。今天只处理一个最常错、最影响分数的小点。
\end{quote}

两个场景使用不同指标：报名场景看完成率、漏项率和误提醒率；提分场景看行动恢复率、同类错点复发率和迁移表现。

\section{关键指标}

\begin{longtable}{@{}p{5.0cm}p{7.0cm}@{}}
\toprule
指标 & 说明 \\
\midrule
提示接受率 & 选择“有用”的比例 \\
误打扰率 & 选择“不是这样”的比例 \\
行动恢复率 & 提示后重新开始行动的比例 \\
同类卡点复发率 & 同一问题在多轮中再次出现的比例 \\
静默与退出率 & 用户主动关闭辅助的频率 \\
主动求助率 & 用户在没有提示时主动请求帮助的比例 \\
学习护照留存 & 长期记录是否被持续使用 \\
\bottomrule
\end{longtable}

\section{V6.6 当前实现状态与验收}

当前实现已经从页面级原型升级为 Windows 真桌面悬浮层：

\begin{itemize}[leftmargin=2em]
\item 独立进程运行，不依赖浏览器；
\item 全屏透明、始终置顶、无边框；
\item 鼠标穿透，不抢键盘焦点；
\item 浅灰半透明底、黑色文字；
\item 屏幕上最多同时两条 Agent 弹幕，其他自动排队；
\item 顶部显示全部在场 Agent；
\item 用户可按 \texttt{Ctrl+Alt+R} 回复，形成一小段功能型多 Agent 对话；
\item 用户回复、反馈和静默写入 DAOQL 与审计链。
\end{itemize}

\subsection{快捷键}

\begin{longtable}{@{}p{4.5cm}p{8.0cm}@{}}
\toprule
快捷键 & 作用 \\
\midrule
\texttt{Ctrl+Alt+Space} & 触发一轮 Agent 弹幕 \\
\texttt{Ctrl+Alt+R} & 打开回复框 \\
\texttt{Ctrl+Alt+P} & 修改 Agent 圆桌成员 \\
\texttt{Ctrl+Alt+B} & 使用当前实时行为测试一轮 \\
\texttt{Ctrl+Alt+1} & 这条有用 \\
\texttt{Ctrl+Alt+2} & 这条不合适 \\
\texttt{Ctrl+Alt+3} & 静默 30 分钟 \\
\texttt{Ctrl+Alt+0} & 清空当前弹幕 \\
\texttt{Ctrl+Alt+Q} & 退出桌面弹幕层 \\
\bottomrule
\end{longtable}

\subsection{工程文件}

\begin{verbatim}
D:\Tender\agent_layer\app\ambient.py
D:\Tender\agent_layer\app\functional_agents.py
D:\Tender\agent_layer\desktop_overlay\tender_overlay.py
D:\Tender\agent_layer\desktop_overlay\config.json
D:\Tender\agent_layer\desktop_overlay\start_overlay.ps1
\end{verbatim}

\subsection{本地监测与隐私}

桌面端只使用本地归一化后的应用类别、短任务提示、空闲时长、窗口切换和输入活跃度。原始窗口标题不上传、不写入 DAOQL；不截图、不读取其他窗口正文、不读取剪贴板、不启用摄像头和麦克风。

\subsection{验收结果}

\begin{itemize}[leftmargin=2em]
\item 桌面窗口具备置顶、鼠标穿透、不抢焦点属性；
\item 功能型 Agent 数量为 8；
\item 每个 Agent 按功能生成不同内容；
\item 两条并发、其余排队的弹幕调度正常；
\item 灰底黑字可读；
\item 回复、反馈、静默和审计事件可追溯；
\item 全量 95 项自动测试通过。
\end{itemize}

\subsection{边界}

多 Agent 不等于多个独立科学专家。营养、心理、神经认知内容属于产品级启发式，不构成诊断或医疗建议。多 Agent 的价值不在于同时说更多，而在于在低打扰条件下调用正确的专业视角。

\section{教学网站、白皮书与最新架构}

Tender Academy 教学网站已经接入本地应用，入口为 \texttt{/guide/}。教学网站按白皮书架构组织，不采用普通功能目录，而是从客观世界与黑盒 API 开始，经过 DAOQL、有限接触面与 Agent Bridge，进入生命树树冠、开花、结果与知识森林。

\begin{center}
\begin{longtable}{@{}p{3.6cm}p{8.0cm}@{}}
\toprule
\textbf{课程主题} & \textbf{要回答的问题} \\
\midrule
生命树与有限世界 & 学习、行动和结果分别属于哪一层？ \\
两条回路 & 人与世界交往和人与人交往有什么不同？ \\
DAOQL、有限接触面与 Agent Bridge & 客观事件如何变成可更正的主观体验？ \\
生命树树冠 & 注意、记忆、语言、内省、情绪与身体如何被组织？ \\
生长、开花与结果 & 内部理解如何进入外部世界并产生分化后果？ \\
语言证据与 Tender 状态 & 话语如何成为弱证据，而不是能力判决？ \\
知识森林 & 全部知识体系如何成为可探索环境？ \\
Agent、治理与路线图 & 支持者如何协作，权力和风险如何受到约束？ \\
\bottomrule
\end{longtable}
\end{center}

教学网站支持中文和 English 切换、学习进度保存、知识检查和打印。教学网站本身不包含数学公式；需要严格定义的模型公式、符号表和实验边界，请阅读《Tender 项目白皮书 V1.7》。

\subsection{配套文档}

\begin{itemize}
\item 项目白皮书 V1.7：预测加工、心智化、系统架构、语言证据、知识森林、证据与治理；
\item 产品手册 V6.21：产品架构、角色、Agent、数据契约、实验与验收；
\item 用户手册 V6.9：第一次启用、桌面弹幕、纠正、隐私与安全；
\item 本教学手册 V5.4：从零开始的环境、数据、模型、Agent 和验收流程。
\end{itemize}

\subsection{本地入口}

\begin{verbatim}
教学网站：http://127.0.0.1:8080/guide/
主应用：http://127.0.0.1:8080/ui/
白皮书：http://127.0.0.1:8080/guide/assets/docs/Tender-Whitepaper-v1.7.pdf
产品手册：http://127.0.0.1:8080/guide/assets/docs/Tender-Product-Manual-v6.21.pdf
用户手册：http://127.0.0.1:8080/guide/assets/docs/Tender-User-Guide-v6.9.pdf
\end{verbatim}

\subsection{边界}

教学网站是理解系统的入口，不是临床诊断工具、教育效果承诺或固定人格测试。学习者、家庭、教师、研究人员和工程人员应通过同一套分层数据边界理解系统能力。


\section{B 端与 C 端交付模型}

Tender 采用“一套核心、三类交付界面”。DAOQL、语言证据、Tender 状态、生命树、知识森林、Agent 运行时、授权和审计是共享核心。不同市场只改变界面、客户、权限、数据范围和商业模式。

\begin{center}
\begin{longtable}{@{}p{3.2cm}p{4.8cm}p{4.6cm}@{}}
\toprule
\textbf{交付方式} & \textbf{界面提供方} & \textbf{主要能力} \\
\midrule
B 端完整系统 & Tender 提供完整教育系统 & 招生、排课、考勤、教学、班主任、执行校长、财务、学情和组织报表 \\
B 端嵌入模块 & 合作方保留原有教育系统界面 & API、SDK、事件契约、单点登录、对象权限与 Agent 能力嵌入 \\
C 端自有界面 & Tender 提供移动端、网页和桌面环境层 & 弹幕、对话条、学习护照、长期时间线与个人 Agent \\
\bottomrule
\end{longtable}
\end{center}

\subsection{B 端两种交付方式}

B 端完整系统适用于希望由 Tender 统一承接核心教育业务的机构。Tender 直接提供教育系统，机构在同一个系统中完成教学、组织、角色、财务和学情管理。

B 端嵌入模块适用于已经拥有 LMS、SIS、CRM、教务、家校或企业培训系统的机构。合作方保留原有界面和业务记录源，Tender 只嵌入完成任务所需的心智化状态、语言证据、知识森林、Agent 与治理能力。嵌入实现必须满足最小事件、对象授权、机构与个人账号分离、用户授权后再进入个人护照，以及不复制合作方全部数据库。

\subsection{C 端自有界面}

C 端由 Tender 提供自己的产品界面，不依附于某个机构后台。用户可以在移动端、网页、桌面环境层、对话条和低频弹幕中获得连续支持。个人学习护照、长期时间线和记录导出归属于用户，不因离开机构而丢失。

\subsection{数据边界}

机构购买系统不等于拥有学习者全部心智数据；合作方保留业务系统不等于 Tender 可以读取其全部数据库；个人使用 C 端也不意味着机构可以读取个人记录。购买关系、监护关系、授权关系和数据处理关系必须分别计算。


\section{当前正式架构：八功能态度、世界树与动态大五}

\subsection{版本边界}

本手册前面的 MBTI、荣格八维与节律章节记录了产品与研究演进过程，其中的实验、功效分析和批判性结论仍然保留。自 v1.7 起，活动产品不再把 MBTI 四字母作为主算法，也不向用户展示类型标签。现行实现采用统一八功能态度测评、大五基线与大五脉搏。

\subsection{两种态度乘以四种功能}

\begin{center}
\begin{longtable}{@{}p{3.0cm}p{2.5cm}p{2.5cm}p{2.5cm}p{2.5cm}@{}}
\toprule
\textbf{态度} & \textbf{感觉} & \textbf{直觉} & \textbf{思维} & \textbf{情感} \\
\midrule
外倾 & Se 现实观察 & Ne 可能性探索 & Te 结果执行 & Fe 关系协调 \\
内倾 & Si 经验对照 & Ni 长期整合 & Ti 结构建模 & Fi 价值边界 \\
\bottomrule
\end{longtable}
\end{center}

每种功能使用三道具体场景题，共 24 题。系统保存连续分数和三题一致性；一致性不足时不强行解释。外倾功能与内倾功能的平均差异形成三种内部世界入口：客观世界优先、主观世界优先、从真实卡点桥接。该入口只影响提问、任务组织和 Agent 支持，不进入能力和成绩模型。

用户界面只显示自然语言，例如“先看现实材料”或“先建立内部框架”，不显示 MBTI 字母、功能代码或固定人格名称。

\subsection{大五基线与大五脉搏}

完整大五用于建立心理资源基线；十题短脉搏用于观察近期变化、波动和恢复。脉搏可以附带关系、角色、任务和时段等最小上下文。系统据此调整沟通和负荷支持：

\begin{itemize}
\item 外向性下降：优先异步、短互动和恢复时间；
\item 宜人性下降：先确认边界与冲突，不解释为敌意；
\item 尽责性下降：缩小承诺与并行任务；
\item 情绪稳定性下降：降低评价威胁，给可控下一步；
\item 开放性下降：先连接熟悉经验，再引入变化。
\end{itemize}

脉搏不等于人格永久改变。数据只用于沟通、协作、负荷和恢复。

\subsection{世界树数据投影}

\texttt{GET /world-tree/\{student\_id\}} 将现有 DAOQL 与 metameta 数据组织为世界树：

\begin{center}
\begin{longtable}{@{}p{3.2cm}p{5.0cm}p{5.0cm}@{}}
\toprule
\textbf{层} & \textbf{当前实现} & \textbf{边界} \\
\midrule
根系 & DAOQL 事件、来源、审计与事件类型统计 & 客观事实，不是主观体验 \\
树干 & Agent Bridge、世界入口、大五脉搏与关系支持 & 支持假设，不是诊断 \\
树冠 & 五类树枝、11 个体验节点、局部状态轴 & 来自派生体验节点，不是脑扫描 \\
开花 & 回应、反思、对话与教学任务事件 & 表示向世界表达或行动 \\
结果 & 作答、基线、元认知与反馈事件 & 结果回流到新事件 \\
共享枝 & 默认未生成；需双方授权 & 只允许共同事实、目标、任务、承诺与结果 \\
\bottomrule
\end{longtable}
\end{center}

“局部状态轴”只定义产品可测量的二能级变量，例如“保持 / 更新”“安全 / 威胁”“表达 / 保留”。它是可视化和测量语言，不代表脑内真实量子态。

\subsection{黄金比与共享枝}

黄金比在 v1.7 中只作为视觉、节律和交互比例的设计假设，用于思考外部输入与内部整合、探索与巩固、共享与私密、支持与自主之间的平衡。它不作为人格评分、关系匹配、能力预测或教育效果承诺。

两个人合作不合并完整私人生命树，只建立双方授权的共享枝。系统比较共同事实、理解差异、互补入口和需要翻译的边界，不输出“匹配分数”。

\subsection{接口与验收}

\begin{verbatim}
GET  /preassessment/items
POST /preassessment/submit/{learner_id}
GET  /preassessment/profile/{learner_id}
GET  /preassessment/pulse/items
POST /preassessment/pulse/{learner_id}
GET  /preassessment/pulse/{learner_id}
GET  /world-tree/{student_id}
\end{verbatim}

验收要求：

\begin{itemize}
\item 用户统一完成 24 题八功能态度，不自报 MBTI；
\item UI 不出现 MBTI 字母或功能代码；
\item 功能一致性进入支持逻辑；
\item 大五基线与脉搏分开；
\item 当前问题、真实行为与结果优先于先验；
\item 世界树不做脑区或量子脑诊断；
\item 未授权时不得生成共享枝；
\item 全量自动测试 95 项通过。
\end{itemize}


\end{document}

